% Gravitational Fields — Physics Upper Sixth — Cours signé OnBuch+
% Official MINESEC syllabus. Compiler avec : tectonic PHY01-gravitational-fields.tex
\newcommand{\DOCMATIERE}{Physics}
\newcommand{\DOCNIVEAU}{Upper Sixth}
\newcommand{\DOCMODULE}{Upper Sixth — Physics}
\newcommand{\DOCLECON}{1}
\newcommand{\DOCTITRE}{Gravitational Fields}
% OnBuch+ — préambule « cours signés » ENRICHI (compilé avec tectonic / XeLaTeX)
% Système visuel « Pop » d'OnBuch+ : fond crème (jamais blanc), bordures encre
% épaisses, ombres « dures » décalées, titres Archivo Black en capitales.
% Version MATHÉMATIQUES : graphes (pgfplots/tikz), boîtes pédagogiques
% (définition, propriété, méthode, attention, le savais-tu, exercice résolu,
% exercice, TP, bilan), page de garde et signature OnBuch+ ancrée en bas.
%
% Macros attendues dans main.tex AVANT \input{preamble} :
%   \DOCMATIERE  \DOCNIVEAU  \DOCMODULE  \DOCLECON (numéro)  \DOCTITRE
\documentclass[11pt]{article}

\usepackage{fontspec}
\usepackage[english]{babel}
\usepackage[a4paper,margin=2cm,top=3.4cm,bottom=2.8cm,headheight=1.4cm,footskip=1.3cm]{geometry}
\usepackage{amsmath,amssymb,amsfonts}
\usepackage{mathtools} % \xmapsto, \coloneqq... souvent utilisés par les modèles
\usepackage{cancel} % simplifications barrées (bilans de forces)
% \abs{z} (module, valeur absolue) et \norm{u} (norme), utilisés par les modèles
\providecommand{\abs}{}\let\abs\relax\DeclarePairedDelimiter{\abs}{\lvert}{\rvert}
\providecommand{\norm}{}\let\norm\relax\DeclarePairedDelimiter{\norm}{\lVert}{\rVert}
\usepackage[table]{xcolor} % [table] : fournit aussi \rowcolors (alternance de lignes)
\usepackage{pagecolor}
\usepackage{tikz}
\usetikzlibrary{arrows.meta,positioning,calc,shapes.geometric,shapes.misc,shapes.arrows,decorations.pathmorphing,decorations.pathreplacing,decorations.markings,patterns,shadows,mindmap,trees,fit,backgrounds,matrix,intersections,angles,quotes}
\usepackage{pgfplots}
\usepackage[version=4]{mhchem} % équations nucléaires \ce{^{235}_{92}U}
\usepackage[european,siunitx]{circuitikz} % schémas électriques
\pgfplotsset{compat=1.18}
\usepgfplotslibrary{fillbetween,groupplots} % aires entre courbes (calcul intégral)
\usepackage{enumitem}
\usepackage{fancyhdr}
% [most] charge toutes les bibliothèques tcolorbox (listings, minted, raster,
% documentation...) et gonfle inutilement la mémoire TeX sur un document long
% (~300 boîtes) : on ne charge que ce qui est réellement utilisé.
\usepackage[skins,breakable]{tcolorbox}
\usepackage{graphicx}
\usepackage{colortbl}
\usepackage{booktabs}
\usepackage{tabularx}
\usepackage{diagbox} % case d'en-tête diagonale des tableaux à double entrée
% Colonnes extensibles centrée / gauche / droite (C, L, R), souvent utilisées par les modèles
\newcolumntype{C}{>{\centering\arraybackslash}X}
\newcolumntype{L}{>{\raggedright\arraybackslash}X}
\newcolumntype{R}{>{\raggedleft\arraybackslash}X}
\usepackage{array}
\usepackage{multicol}
\usepackage{siunitx}
% parse-numbers=false : accepte aussi \SI{2,5 \times 10^{-3}}{...}
\sisetup{locale=UK,output-decimal-marker={.},group-minimum-digits=4,parse-numbers=false}
\DeclareSIUnit\ohm{\ensuremath{\symup{\Omega}}} % U+2126 absent de la police
\DeclareSIUnit\division{div} % réglages d'oscilloscope (ms/div, V/div)
\DeclareSIUnit\tr{tr} % tours (vitesse de rotation, ex. \SI{1500}{\tr\per\minute})
\usepackage{qrcode}
% \degree : commande fréquemment utilisée par les modèles pour un angle ou une
% température en dehors de \SI{}{} (non fournie par les paquets chargés).
\providecommand{\degree}{^{\circ}}
% \qed (fin de démonstration), utilisé par les modèles sans amsthm
\providecommand{\qed}{\hfill$\square$}
% \barre : double barre des valeurs interdites dans un tableau de variations (tkz-tab)
\providecommand{\barre}{\ensuremath{\|}}
% environnement proof (amsthm non chargé) utilisé par les modèles
\ifdefined\proof\else\newenvironment{proof}[1][Proof]{\par\noindent\textit{#1.}\ }{\qed\par}\fi
\usepackage{lastpage}
\usepackage{hyperref}
\hypersetup{hidelinks}

% ============================================================
% Palette « Pop » OnBuch+ — mêmes hex que la classe P (plus_ui.dart)
% ============================================================
\definecolor{poporange}{HTML}{F59321}
\definecolor{popdark}{HTML}{E07A0C}
\definecolor{popcream}{HTML}{FFF6E8}
\definecolor{popink}{HTML}{1C1714}
\definecolor{poppurple}{HTML}{7A5AE0}
\definecolor{popgreen}{HTML}{2E9E6A}
\definecolor{poppink}{HTML}{E0457B}
\definecolor{popblue}{HTML}{2F7DE1}
\definecolor{popgold}{HTML}{C98A12}
\definecolor{popmuted}{HTML}{8A7F76}
\definecolor{popline}{HTML}{EADFD2}
\definecolor{popgray}{HTML}{8A7F76}
\colorlet{popgrey}{popgray}
% Alias tolérants (le modèle nomme parfois les couleurs différemment)
\colorlet{popwhite}{white}
\colorlet{popblack}{popink}
\colorlet{popred}{poppink}
\colorlet{popyellow}{popgold}
% Teintes claires pour fonds de tableaux / zones
\colorlet{popblueL}{popblue!10!white}
\colorlet{poporangeL}{poporange!14!white}
\colorlet{popgreenL}{popgreen!12!white}
\colorlet{poppurpleL}{poppurple!12!white}
\colorlet{poppinkL}{poppink!12!white}
\colorlet{popgoldL}{popgold!14!white}

\pagecolor{popcream}

% ============================================================
% Typographie
% ============================================================
\defaultfontfeatures{Ligatures=TeX}
\setmainfont{PlusJakartaSans-Medium.ttf}[Path=../../../fonts/,BoldFont=PlusJakartaSans-Bold.ttf]
% Maths en sans-serif (Fira Math) avec chiffres et lettres droites de Plus
% Jakarta, pour que les formules se fondent dans le texte.
\usepackage[math-style=ISO,bold-style=ISO]{unicode-math}
\setmathfont{FiraMath-Regular.otf}[Path=../../../fonts/]
\setmathfont{PlusJakartaSans-Medium.ttf}[Path=../../../fonts/,range={up/num,up/{Latin,latin},\mathup}]
% (icomma retiré : décimales avec point en anglais)
% Caractères mathématiques Unicode absents des polices de texte (Plus Jakarta,
% Archivo Black) : rendus par leur équivalent mathématique, en texte comme en
% formule — sinon ils s'affichent en carré vide (ex. « Arithmétique dans ℤ »).
\usepackage{newunicodechar}
\newunicodechar{ℝ}{\ensuremath{\mathbb{R}}}
\newunicodechar{ℤ}{\ensuremath{\mathbb{Z}}}
\newunicodechar{ℂ}{\ensuremath{\mathbb{C}}}
\newunicodechar{ℕ}{\ensuremath{\mathbb{N}}}
\newunicodechar{ℚ}{\ensuremath{\mathbb{Q}}}
\newunicodechar{𝔻}{\ensuremath{\mathbb{D}}}
\newunicodechar{α}{\ensuremath{\alpha}}
\newunicodechar{β}{\ensuremath{\beta}}
\newunicodechar{γ}{\ensuremath{\gamma}}
\newunicodechar{δ}{\ensuremath{\delta}}
\newunicodechar{ε}{\ensuremath{\varepsilon}}
\newunicodechar{θ}{\ensuremath{\theta}}
\newunicodechar{λ}{\ensuremath{\lambda}}
\newunicodechar{μ}{\ensuremath{\mu}}
\newunicodechar{σ}{\ensuremath{\sigma}}
\newunicodechar{τ}{\ensuremath{\tau}}
\newunicodechar{φ}{\ensuremath{\varphi}}
\newunicodechar{ψ}{\ensuremath{\psi}}
\newunicodechar{ω}{\ensuremath{\omega}}
\newunicodechar{Σ}{\ensuremath{\Sigma}}
% majuscules produites par \MakeUppercase dans les titres (α-aminés -> Α)
\newunicodechar{Α}{\ensuremath{\alpha}}
\newunicodechar{Β}{\ensuremath{\beta}}
\newunicodechar{Γ}{\ensuremath{\Gamma}}
\newunicodechar{Δ}{\ensuremath{\Delta}}
\newunicodechar{Ε}{\ensuremath{\varepsilon}}
\newunicodechar{Θ}{\ensuremath{\Theta}}
\newunicodechar{Λ}{\ensuremath{\Lambda}}
\newunicodechar{Μ}{\ensuremath{\mu}}
\newunicodechar{Τ}{\ensuremath{\tau}}
\newunicodechar{Φ}{\ensuremath{\Phi}}
\newunicodechar{Ψ}{\ensuremath{\Psi}}
\newunicodechar{Ω}{\ensuremath{\Omega}}
\newunicodechar{↦}{\ensuremath{\mapsto}}
\newunicodechar{∈}{\ensuremath{\in}}
\newunicodechar{∉}{\ensuremath{\notin}}
\newunicodechar{∧}{\ensuremath{\wedge}}
\newunicodechar{⇔}{\ensuremath{\Leftrightarrow}}
\newunicodechar{⟺}{\ensuremath{\Longleftrightarrow}}
\newunicodechar{⇒}{\ensuremath{\Rightarrow}}
\newunicodechar{⟹}{\ensuremath{\Longrightarrow}}
\newunicodechar{∘}{\ensuremath{\circ}}
\newunicodechar{∪}{\ensuremath{\cup}}
\newunicodechar{∩}{\ensuremath{\cap}}
\newunicodechar{≡}{\ensuremath{\equiv}}
\newunicodechar{⊂}{\ensuremath{\subset}}
\newunicodechar{⊥}{\ensuremath{\perp}}
\newunicodechar{∃}{\ensuremath{\exists}}
\newunicodechar{∀}{\ensuremath{\forall}}
\newunicodechar{∛}{\ensuremath{\sqrt[3]{\,}}}
\newunicodechar{∜}{\ensuremath{\sqrt[4]{\,}}}
\newunicodechar{⃗}{\ensuremath{{}^{\to}}}
\newunicodechar{ⁿ}{\ifmmode^{n}\else\textsuperscript{\ensuremath{n}}\fi}
\newunicodechar{ˣ}{\ifmmode^{x}\else\textsuperscript{\ensuremath{x}}\fi}
\newunicodechar{ᵇ}{\ifmmode^{b}\else\textsuperscript{\ensuremath{b}}\fi}
\newunicodechar{ʳ}{\ifmmode^{r}\else\textsuperscript{\ensuremath{r}}\fi}
\newunicodechar{ᵘ}{\ifmmode^{u}\else\textsuperscript{\ensuremath{u}}\fi}
\newunicodechar{ʸ}{\ifmmode^{y}\else\textsuperscript{\ensuremath{y}}\fi}
\newunicodechar{ᵃ}{\ifmmode^{a}\else\textsuperscript{\ensuremath{a}}\fi}
\newunicodechar{ᵉ}{\ifmmode^{e}\else\textsuperscript{\ensuremath{e}}\fi}
\newunicodechar{ᵏ}{\ifmmode^{k}\else\textsuperscript{\ensuremath{k}}\fi}
\newunicodechar{ᵖ}{\ifmmode^{p}\else\textsuperscript{\ensuremath{p}}\fi}
\newunicodechar{ᴺ}{\ifmmode^{N}\else\textsuperscript{\ensuremath{N}}\fi}
\newunicodechar{ᶜ}{\ifmmode^{c}\else\textsuperscript{\ensuremath{c}}\fi}
\newunicodechar{ʰ}{\ifmmode^{h}\else\textsuperscript{\ensuremath{h}}\fi}
\newunicodechar{ᵠ}{\ifmmode^{q}\else\textsuperscript{\ensuremath{q}}\fi}
\newunicodechar{ₙ}{\ifmmode_{n}\else\textsubscript{\ensuremath{n}}\fi}
\newunicodechar{ₐ}{\ifmmode_{a}\else\textsubscript{\ensuremath{a}}\fi}
\newunicodechar{ᵢ}{\ifmmode_{i}\else\textsubscript{\ensuremath{i}}\fi}
\newunicodechar{ᵣ}{\ifmmode_{r}\else\textsubscript{\ensuremath{r}}\fi}
\newunicodechar{ₖ}{\ifmmode_{k}\else\textsubscript{\ensuremath{k}}\fi}
\newunicodechar{ᵦ}{\ifmmode_{\beta}\else\textsubscript{\ensuremath{\beta}}\fi}
\newunicodechar{ₓ}{\ifmmode_{x}\else\textsubscript{\ensuremath{x}}\fi}
\newfontfamily\popdisplay{ArchivoBlack-Regular.ttf}[Path=../../../fonts/]
\newfontfamily\poplogo{SpaceGrotesk-Bold.ttf}[Path=../../../fonts/]
\newfontfamily\popsemi{PlusJakartaSans-SemiBold.ttf}[Path=../../../fonts/]
\newfontfamily\popxbold{PlusJakartaSans-ExtraBold.ttf}[Path=../../../fonts/]
\newfontfamily\popxboldls{PlusJakartaSans-ExtraBold.ttf}[Path=../../../fonts/,LetterSpace=3.0]

\setlist[enumerate]{leftmargin=1.7em,itemsep=5pt,topsep=4pt}
\setlist[itemize]{leftmargin=1.7em,itemsep=4pt,topsep=4pt,label={\color{poporange}\textbullet}}
\setlength{\parindent}{0pt}
\setlength{\parskip}{5pt}
\renewcommand{\arraystretch}{1.25}

% pgfplots : style Pop par défaut
\pgfplotsset{
  every axis/.append style={axis line style={popink,line width=1pt},tick style={popink},
    grid style={popline,line width=0.5pt},label style={font=\small},tick label style={font=\footnotesize}},
  popcurve/.style={poporange,line width=1.8pt,no markers,smooth},
}
% Même style disponible dans un \draw TikZ simple (hors pgfplots)
\tikzset{popcurve/.style={poporange,line width=1.8pt}}
\tikzset{popgrid/.style={popline,very thin}}
\colorlet{popcurve}{poporange} % le nom du style est parfois utilisé comme couleur

% ============================================================
% Pill / squiggle
% ============================================================
\newcommand{\poppill}[3]{%
  \tikz[baseline=-0.55ex]\node[draw=popink,line width=1.2pt,rounded corners=2.6pt,%
    fill=#2,inner xsep=6.5pt,inner ysep=3pt]{%
    \popxboldls\fontsize{7}{7}\selectfont\color{#3}\MakeUppercase{#1}};%
}
\newcommand{\popsquiggle}[1]{%
  \tikz[baseline=-0.4ex]{
    \draw[#1,line width=2.6pt,line cap=round]
      plot [smooth] coordinates {(0,0.09) (0.34,0.01) (0.68,0.21) (1.02,0.01) (1.36,0.10)};
  }%
}
% Mot-clé surligné (marqueur orange sous le texte)
\newcommand{\cle}[1]{{\popxbold\color{popdark}#1}}

% ============================================================
% En-tête / pied
% ============================================================
\pagestyle{fancy}
\fancyhf{}
\renewcommand{\headrulewidth}{0pt}
\renewcommand{\footrulewidth}{0pt}
\lhead{\raisebox{-0.15ex}{\poplogo\fontsize{13}{13}\selectfont\color{popink}OnBuch\color{poporange}+}}
\rhead{\poppill{\DOCMATIERE\ \textbullet\ \DOCNIVEAU\ \textbullet\ Chapter \DOCLECON}{white}{popink}}
\lfoot{\popxbold\fontsize{7.6}{7.6}\selectfont\color{popmuted}\MakeUppercase{Signed course OnBuch+}\ \textbullet\ {\popsemi\DOCTITRE}}
\rfoot{\tikz[baseline=-0.6ex]\node[rounded corners=8.5pt,fill=popink,minimum height=17pt,minimum width=17pt,inner xsep=5pt,inner sep=0]{\popxbold\fontsize{8}{8}\selectfont\color{popcream}\thepage\,/\,\pageref*{LastPage}};}

% ============================================================
% Titres
% ============================================================
\newcommand{\chaptitre}[2]{%
  \begin{tcolorbox}[enhanced,colback=poporange,colframe=popink,boxrule=2.6pt,arc=11pt,
      drop shadow=popink,left=15pt,right=15pt,top=12pt,bottom=11pt,before skip=3pt,after skip=18pt]
    \poppill{#2}{popink}{popcream}\par\vspace{9pt}
    {\raggedright\hyphenpenalty=10000\exhyphenpenalty=10000\popdisplay\fontsize{22}{24}\selectfont\color{popcream}\MakeUppercase{#1}\par}
  \end{tcolorbox}
}
% Section numérotée : pastille orange avec le numéro + titre Archivo
\newcounter{popsec}
\newcommand{\coursec}[1]{%
  \stepcounter{popsec}\setcounter{popsub}{0}%
  \par\vspace{16pt}\noindent
  \begin{minipage}{\linewidth}
  \tikz[baseline=-0.7ex]\node[draw=popink,line width=1.6pt,fill=poporange,rounded corners=4pt,minimum size=22pt,inner sep=0]{\popdisplay\fontsize{12}{12}\selectfont\color{popcream}\thepopsec};%
  \hspace{8pt}{\popdisplay\fontsize{15}{17}\selectfont\color{popink}\MakeUppercase{#1}}\par
  \vspace{4pt}\noindent{\color{poporange}\rule{36pt}{3.6pt}}
  \end{minipage}\par\nopagebreak\vspace{8pt}
}
\newcounter{popsub}
\newcommand{\courssub}[1]{%
  \stepcounter{popsub}%
  \par\vspace{9pt}\noindent{\popxbold\fontsize{11.5}{13}\selectfont\color{popink}{\color{poporange}\thepopsec.\thepopsub}\hspace{6pt}#1}\par\nopagebreak\vspace{4pt}
}

% ============================================================
% Boîtes pédagogiques — même recette : fond blanc, bordure épaisse colorée,
% ombre dure encre, badge plein. Titre optionnel [#1].
% ============================================================
\tcbset{popbase/.style={breakable,enhanced,colback=white,boxrule=2.3pt,arc=9pt,
  shadow={3pt}{-3pt}{0pt}{popink},left=14pt,right=14pt,top=11pt,bottom=11pt,before skip=11pt,after skip=15pt,
  fontupper=\popsemi\fontsize{10.6}{14.5}\selectfont\color{popink},
  fontlower=\popsemi\fontsize{10.6}{14.5}\selectfont\color{popink}}}
% Coche dessinée (le glyphe ✓ n'existe pas dans les polices du document)
\newcommand{\popcheck}[1]{\tikz[baseline=-0.5ex]\draw[#1,line width=1.6pt,line cap=round,line join=round](-0.13,0.0)--(-0.04,-0.09)--(0.13,0.1);}
\providecommand{\coche}{\popcheck{popgreen}}
\providecommand{\resultat}[1]{\par\smallskip\noindent\fcolorbox{popgreen}{popgreenL}{\parbox{\dimexpr\linewidth-2\fboxsep-2\fboxrule\relax}{\textbf{Result.} #1}}\par\smallskip}
\newcommand{\popboxhead}[3]{\poppill{#1}{#2}{white}\ifx\relax#3\relax\else\hspace{6pt}{\popxbold\fontsize{10.6}{12}\selectfont\color{popink}#3}\fi\par\vspace{7pt}}

\newtcolorbox{aretenir}[1][]{popbase,colframe=popgreen,before upper={\popboxhead{Key points}{popgreen}{#1}}}
\newtcolorbox{exemplebox}[1][]{popbase,colframe=poporange,before upper={\popboxhead{Exemple}{poporange}{#1}}}
\newtcolorbox{definition}[1][]{popbase,colframe=popblue,before upper={\popboxhead{Definition}{popblue}{#1}}}
\newtcolorbox{propriete}[1][]{popbase,colframe=poppurple,before upper={\popboxhead{Property}{poppurple}{#1}}}
\newtcolorbox{methode}[1][]{popbase,colframe=popgold,before upper={\popboxhead{Method}{popgold}{#1}}}
\newtcolorbox{attention}[1][]{popbase,colframe=poppink,before upper={\popboxhead{Warning}{poppink}{#1}}}
\newtcolorbox{experience}[1][]{popbase,colframe=popblue,colback=popblueL,before upper={\popboxhead{Experiment}{popblue}{#1}}}
% « Le savais-tu ? » : carte violette pleine (contexte, culture, Cameroun)
\newtcolorbox{savaistu}[1][]{popbase,colback=poppurple,colframe=popink,
  fontupper=\popsemi\fontsize{10.4}{14.2}\selectfont\color{white},
  before upper={\poppill{Did you know?}{white}{poppurple}\ifx\relax#1\relax\else\hspace{6pt}{\popxbold\color{white}#1}\fi\par\vspace{7pt}}}
% Exercice résolu : énoncé en haut, \tcblower puis la solution sur fond crème
\newcounter{popexo}\newcounter{popexr}
\newtcolorbox{exoresolu}[1][]{popbase,colframe=poporange,colbacklower=popcream,
  segmentation style={popink,line width=1.2pt,dashed},
  before upper={\stepcounter{popexr}\popboxhead{Worked example \thepopexr}{poporange}{#1}},
  before lower={\poppill{Solution}{popink}{popcream}\par\vspace{6pt}}}
% Exercice (non résolu en place) — la correction est dans la partie Corrigés
\newtcolorbox{exercice}[1][]{popbase,colframe=popink,
  before upper={\stepcounter{popexo}\popboxhead{Exercise \thepopexo}{popink}{#1}}}
% Corrigé
\newtcolorbox{corrige}[1][]{popbase,colframe=popgreen,colback=popgreenL,
  before upper={\popboxhead{Solution}{popgreen}{#1}}}
% Objectifs / prérequis / bilan
\newtcolorbox{objectifs}{popbase,colframe=popink,colback=poporangeL,
  before upper={\popboxhead{Chapter objectives}{popink}{}}}
\newtcolorbox{prerequis}{popbase,colframe=popink,colback=popblueL,
  before upper={\popboxhead{Prerequisites}{popblue}{}}}
\newtcolorbox{bilan}{popbase,colframe=popink,colback=white,boxrule=2.8pt,
  before upper={\popboxhead{Fiche bilan}{poporange}{L'essentiel à connaître pour l'examen}}}
% Figure encadrée avec légende
\newtcolorbox{popfigure}[1][]{enhanced,breakable=false,colback=white,colframe=popline,boxrule=1.4pt,arc=8pt,
  left=10pt,right=10pt,top=10pt,bottom=8pt,before skip=10pt,after skip=12pt,halign=center,
  lower separated=false,halign lower=center,
  fontlower=\popsemi\fontsize{9}{11.5}\selectfont\color{popmuted},
  #1}
\newcommand{\legende}[1]{\par\vspace{6pt}{\popsemi\fontsize{9}{11.5}\selectfont\color{popmuted}#1}}

% ============================================================
% Signature OnBuch+
% ============================================================
\newcommand{\poplogoinline}[1]{{\poplogo\fontsize{#1}{#1}\selectfont\color{popink}OnBuch\color{poporange}+}}
\newcommand{\signatureonbuch}{%
  \begin{tcolorbox}[enhanced,colback=popink,colframe=popink,boxrule=2.2pt,arc=10pt,
      drop shadow=poporange,left=14pt,right=14pt,top=10pt,bottom=10pt,before skip=0pt,after skip=0pt]
    \tikz[baseline=-0.7ex]\node[circle,fill=popgreen,draw=popcream,line width=1.3pt,minimum size=22pt,inner sep=0]{\popcheck{white}};%
    \hspace{9pt}\begin{minipage}[c]{0.62\linewidth}
      {\popxbold\fontsize{12}{13}\selectfont\color{popcream}Signed\ \textbullet\ {\poplogo OnBuch\color{poporange}+}}\par\vspace{2pt}
      {\raggedright\popsemi\fontsize{8}{10.5}\selectfont\color{popline}\MakeUppercase{OnBuch+ teaching team}\\\MakeUppercase{Follows the official MINESEC syllabus}\par}
    \end{minipage}\hfill
    {\poplogo\fontsize{22}{22}\selectfont\color{popcream}OnBuch\color{poporange}+}
  \end{tcolorbox}%
}

% ============================================================
% Page de garde
% #1 titre · #2 matière · #3 niveau · #4 module · #5 n° leçon · #6 durée ·
% #7 description / objectifs (texte ou itemize)
% ============================================================
\newcommand{\pagegarde}[7]{%
  \thispagestyle{empty}
  \newgeometry{a4paper,margin=1.7cm,top=1.6cm,bottom=1.4cm}%
  \begin{tikzpicture}[remember picture,overlay]
    \fill[poporange!18] ([xshift=-3.2cm,yshift=-3.6cm]current page.north east) circle (4.6cm);
    \fill[poppurple!14] ([xshift=2.4cm,yshift=7.6cm]current page.south west) circle (3.8cm);
  \end{tikzpicture}%
  {\poplogo\fontsize{30}{30}\selectfont\color{popink}OnBuch\color{poporange}+}\hfill
  \raisebox{0.6ex}{\poppill{Signed course \textbullet\ Premium}{popink}{popcream}}\par
  \vspace{1pt}\popsquiggle{poppurple}\par
  \vspace{8pt}
  \begin{tcolorbox}[enhanced,colback=poporange,colframe=popink,boxrule=2.8pt,arc=13pt,
      drop shadow=popink,left=18pt,right=18pt,top=12pt,bottom=13pt,after skip=10pt]
    \poppill{#2 \textbullet\ #3}{popink}{popcream}\hspace{5pt}\poppill{#4}{white}{popink}\par\vspace{8pt}
    {\popdisplay\fontsize{14}{14}\selectfont\color{popink}CHAPTER #5}\par\vspace{4pt}
    {\raggedright\hyphenpenalty=10000\exhyphenpenalty=10000\popdisplay\fontsize{25}{27}\selectfont\color{popcream}\MakeUppercase{#1}\par}
    \vspace{8pt}
    {\popxbold\fontsize{9.5}{9.5}\selectfont\color{popink}\MakeUppercase{Suggested time: #6}}
  \end{tcolorbox}
  \begin{tcolorbox}[enhanced,colback=white,colframe=popink,boxrule=2.2pt,arc=10pt,
      drop shadow=popink,left=16pt,right=16pt,top=10pt,bottom=10pt,after skip=8pt]
    \poppill{In this chapter}{poppurple}{white}\par\vspace{5pt}
    \popsemi\fontsize{9.3}{12.6}\selectfont\color{popink}#7
  \end{tcolorbox}
  \noindent\begin{minipage}[t]{0.487\linewidth}
    \begin{tcolorbox}[enhanced,colback=popgreen,colframe=popink,boxrule=2.2pt,arc=10pt,
        drop shadow=popink,left=11pt,right=11pt,top=10pt,bottom=10pt,height=3.1cm,valign=center]
      \tcbox[colback=white,colframe=popink,boxrule=1.3pt,arc=5pt,boxsep=0pt,nobeforeafter,
        left=3pt,right=3pt,top=3pt,bottom=3pt,valign=center]{\qrcode[height=1.9cm]{https://play.google.com/store/apps/details?id=com.luvvix.onbuch&hl=en}}%
      \hspace{8pt}\begin{minipage}[c]{0.5\linewidth}
        {\popdisplay\fontsize{11}{12.5}\selectfont\color{white}\MakeUppercase{Download OnBuch}}\par\vspace{4pt}
        {\raggedright\popsemi\fontsize{8.4}{11}\selectfont\color{white}Free \textbullet\ Google Play\\AI tutor, past papers, results\par}
      \end{minipage}
    \end{tcolorbox}
  \end{minipage}\hfill
  \begin{minipage}[t]{0.487\linewidth}
    \begin{tcolorbox}[enhanced,colback=poppurple,colframe=popink,boxrule=2.2pt,arc=10pt,
        drop shadow=popink,left=11pt,right=11pt,top=10pt,bottom=10pt,height=3.1cm,valign=center]
      \tikz[baseline=-0.7ex]\node[circle,fill=white,draw=popink,line width=1.3pt,minimum size=1.6cm,inner sep=0]{\popdisplay\fontsize{24}{24}\selectfont\color{poppurple}+};%
      \hspace{8pt}\begin{minipage}[c]{0.6\linewidth}
        {\popdisplay\fontsize{11}{12.5}\selectfont\color{white}\MakeUppercase{Go OnBuch+}}\par\vspace{4pt}
        {\raggedright\popsemi\fontsize{8.4}{11}\selectfont\color{white}All signed courses, unlimited AI tutor, PDF sheets — OnBuch+ tab in the app\par}
      \end{minipage}
    \end{tcolorbox}
  \end{minipage}
  \vspace{8pt}
  \popcontenu
  \vfill
  \signatureonbuch
  \restoregeometry
  \newpage
}

% Rangée « Ce cours contient » (page de garde)
\newcommand{\poptile}[3]{% #1 couleur #2 gros texte #3 légende
  \begin{tcolorbox}[enhanced,colback=#1,colframe=popink,boxrule=2pt,arc=9pt,shadow={2.5pt}{-2.5pt}{0pt}{popink},
      left=6pt,right=6pt,top=9pt,bottom=9pt,halign=center,valign=center,height=1.75cm,width=\linewidth,nobeforeafter]
    {\popdisplay\fontsize{12.5}{13}\selectfont\color{white}\MakeUppercase{#2}}\par\vspace{3pt}
    {\centering\popsemi\fontsize{7.8}{9.5}\selectfont\color{white}#3\par}
  \end{tcolorbox}}
\newcommand{\popcontenu}{%
  \par\noindent{\popxboldls\fontsize{7.5}{7.5}\selectfont\color{popmuted}THIS SIGNED COURSE CONTAINS}\par\vspace{7pt}
  \noindent\begin{minipage}[t]{0.235\linewidth}\poptile{popblue}{Course}{complete, illustrated, step by step}\end{minipage}\hfill
  \begin{minipage}[t]{0.235\linewidth}\poptile{poporange}{Methods}{and worked examples}\end{minipage}\hfill
  \begin{minipage}[t]{0.235\linewidth}\poptile{poppink}{Exercises}{exam style + solutions}\end{minipage}\hfill
  \begin{minipage}[t]{0.235\linewidth}\poptile{popgreen}{Summary sheet}{the essentials to remember}\end{minipage}\par}

% Sommaire
\newtcolorbox{sommaire}{enhanced,colback=white,colframe=popink,boxrule=2.2pt,arc=10pt,
  drop shadow=popink,left=18pt,right=18pt,top=13pt,bottom=13pt,before skip=6pt,after skip=20pt,
  before upper={\poppill{Contents}{poppurple}{white}\par\vspace{9pt}\popsemi\fontsize{10.6}{14.5}\selectfont\color{popink}}}

% Signature de fin de cours — poussée tout en bas de la dernière page
\newcommand{\findecours}{%
  \par\vfill\nopagebreak
  \begin{center}\popsquiggle{poporange}\end{center}\vspace{4pt}
  \signatureonbuch
}
\AtBeginDocument{\def\llbracket{\mathopen{[\mkern-3.5mu[}}\def\rrbracket{\mathclose{]\mkern-3.5mu]}}} % intervalles d'entiers (glyphe ⟦ absent de la police)
% Glyphes absents de Fira Math : redéfinis à partir de glyphes disponibles
\AtBeginDocument{%
  \def\setminus{\mathbin{\backslash}}\let\smallsetminus\setminus
  \def\vdots{\mathord{\vbox{\baselineskip3.2pt\lineskiplimit0pt\kern4pt\hbox{.}\hbox{.}\hbox{.}}}}%
  \def\ddots{\mathinner{\mkern1mu\raise6.5pt\hbox{.}\mkern2mu\raise3.6pt\hbox{.}\mkern2mu\raise0.7pt\hbox{.}\mkern1mu}}%
  \def\triangle{\mathord{\scalebox{0.9}{$\symup{\Delta}$}}}%
  \def\nabla{\mathord{\raisebox{\depth}{\scalebox{1}[-1]{$\symup{\Delta}$}}}}%
  \def\checkmark{\popcheck{popdark}}%
  \def\star{\mathbin{\ast}}\def\bigstar{\mathord{\ast}}%
  \def\diamond{\mathbin{\scalebox{0.6}{$\Diamond$}}}%
  \def\bigcup{\mathop{\mathchoice{\raisebox{-0.3em}{\scalebox{1.7}{$\cup$}}}{\scalebox{1.25}{$\cup$}}{\cup}{\cup}}\displaylimits}%
  \def\bigcap{\mathop{\mathchoice{\raisebox{-0.3em}{\scalebox{1.7}{$\cap$}}}{\scalebox{1.25}{$\cap$}}{\cap}{\cap}}\displaylimits}%
  \def\lesseqqgtr{\mathrel{\lessgtr}}\def\bigoplus{\mathop{\oplus}\displaylimits}%
}
\providecommand{\ssi}{if and only if} % abréviation fréquente
\AtBeginDocument{\providecommand{\wideparen}{\overparen}} % arc de cercle

%% FIGFIX-BEGIN v5 — correctifs globaux des figures (docs/onbuch-generation/tools/figfix)
\usepackage{adjustbox}\usepackage{etoolbox}\tcbuselibrary{raster}
% toute figure TikZ/pgfplots plus large que la ligne est réduite à la largeur disponible
\BeforeBeginEnvironment{tikzpicture}{\begin{adjustbox}{max width=\linewidth}}
\AfterEndEnvironment{tikzpicture}{\end{adjustbox}}
% légendes pgfplots lisibles par-dessus les courbes
\pgfplotsset{every axis/.append style={legend style={fill=white,fill opacity=0.92,text opacity=1,draw=black!15,inner sep=3pt}}}
% étiquettes d'arêtes (syntaxe edge["..."]) avec fond blanc
\tikzset{every edge quotes/.append style={fill=white,inner sep=1.2pt}}
%% FIGFIX-END
\begin{document}

\pagegarde{Gravitational Fields}{Physics}{Upper Sixth}{Upper Sixth — Physics}{1}{4 periods}{This chapter studies gravitational fields: Newton's inverse-square law, Kepler's laws, the variation of g, gravitational potential and potential energy, escape and orbital speeds, geostationary satellites, and projectile motion in a uniform field. It provides the essential tools for the GCE Advanced Level mechanics paper.
\vspace{4pt}
\begin{multicols}{2}\small
\begin{itemize}
\item By the end of this chapter I can state and apply Newton's law of universal gravitation between two point masses
\item By the end of this chapter I can state Kepler's three laws and use the third law to relate orbital period and radius
\item By the end of this chapter I can define gravitational field strength and derive g = GM/r² for a planet
\item By the end of this chapter I can describe and calculate the variation of g with altitude, depth and latitude
\item By the end of this chapter I can define gravitational potential and potential energy and compute the work done moving a mass in a field
\item By the end of this chapter I can derive and calculate escape velocity and orbital speed for a satellite
\end{itemize}\end{multicols}}

\begin{sommaire}
\begin{multicols}{2}
\begin{enumerate}
\item Newton's Law of Gravitation and Kepler's Laws
\item Gravitational Field Strength and Variation of g
\item Gravitational Potential, Potential Energy, Escape and Orbital Speeds
\item Motion in the Gravitational Field: Projectile Motion
\item Integration activity
\item Methods and worked examples
\item Exercises
\item Solutions to the exercises
\item Summary sheet
\end{enumerate}
\end{multicols}
\end{sommaire}

\begin{objectifs}
\begin{itemize}
\item By the end of this chapter I can state and apply Newton's law of universal gravitation between two point masses
\item By the end of this chapter I can state Kepler's three laws and use the third law to relate orbital period and radius
\item By the end of this chapter I can define gravitational field strength and derive g = GM/r² for a planet
\item By the end of this chapter I can describe and calculate the variation of g with altitude, depth and latitude
\item By the end of this chapter I can define gravitational potential and potential energy and compute the work done moving a mass in a field
\item By the end of this chapter I can derive and calculate escape velocity and orbital speed for a satellite
\item By the end of this chapter I can explain the conditions for a geostationary orbit and compute its radius
\item By the end of this chapter I can analyse projectile motion in a uniform gravitational field (range, maximum height, time of flight)
\end{itemize}
\end{objectifs}

\begin{prerequis}
\begin{itemize}
\item Newton's laws of motion and the concept of weight (Lower Sixth mechanics)
\item Uniform circular motion: centripetal acceleration a = v²/r = ω²r (Lower Sixth)
\item Work, kinetic energy and potential energy in a uniform field (Lower Sixth)
\item Projectile motion basics from Lower Sixth mechanics
\item Algebraic manipulation of powers, square roots and scientific notation
\end{itemize}
\end{prerequis}

\newpage

\coursec{Newton's Law of Gravitation and Kepler's Laws}

Every evening, the CRTV signal reaches your television from a satellite parked $35\,786\ \mathrm{km}$ above the equator --- yet nobody ever launched it with a ladder or a rope. How fast must a satellite move to ``hang'' above the same point of Cameroon without falling? Why does a mango dropped in Bamenda fall slightly differently from one dropped at the top of Mount Cameroon? And why does a stone thrown horizontally from a school field in Buea follow a curved path? This chapter answers these questions with one single idea: \cle{Newton's law of universal gravitation}.

\courssub{From falling bodies to planetary motion: a historical sketch}

For centuries, people believed that the physics of the heavens was different from the physics of the Earth. A mango falls from a tree in Bamenda; the Moon circles the Earth; the planets wander across the sky. These looked like three unrelated phenomena.

The story changed in the seventeenth century. First, \cle{Galileo} showed that all bodies near the Earth's surface fall with the same acceleration $g$ (air resistance apart). Then \cle{Tycho Brahe} accumulated decades of precise observations of planetary positions, and his assistant \cle{Johannes Kepler} extracted from this data three empirical laws describing how planets move around the Sun.

The genius of \cle{Isaac Newton} (around 1687, in his \emph{Principia}) was to realise that \emph{one single force} explains all of it. The famous --- possibly legendary --- apple suggested the question: if gravity pulls the apple down, does the same pull reach the Moon? Newton's answer: yes. The Moon is continuously ``falling'' towards the Earth, but its sideways velocity is so large that it perpetually misses. The same attractive force that drops a mango also steers the Moon and the planets. This is why the law is called \emph{universal}.

\begin{savaistu}[A force you use every day without knowing it]
Every GPS position, every satellite TV channel, every tide table on the coast of Kribi and Limbe is computed from Newton's law of gravitation, written down more than three centuries ago. The same equation that describes a falling mango describes a satellite $35\,786\ \mathrm{km}$ above the Atlantic.
\end{savaistu}

\courssub{Newton's law of universal gravitation}

\begin{definition}[Newton's law of universal gravitation]
Every particle of matter in the universe attracts every other particle with a force that is directly proportional to the product of their masses and inversely proportional to the square of the distance between them:
\[
F = \frac{G\,M\,m}{r^{2}}
\]
where $F$ is the magnitude of the gravitational force (in $\mathrm{N}$), $M$ and $m$ are the two masses (in $\mathrm{kg}$), $r$ is the distance between their centres (in $\mathrm{m}$), and $G$ is the \cle{universal gravitational constant}:
\[
G = 6.67\times 10^{-11}\ \mathrm{N\,m^{2}\,kg^{-2}}.
\]
\end{definition}

Because the force varies as $1/r^{2}$, we say that gravitation obeys an \cle{inverse-square law}: doubling the separation divides the force by four, tripling it divides the force by nine. This $1/r^{2}$ dependence is the single most important feature of the law, and it reappears everywhere in this chapter (in $g$, in the potential, in orbital mechanics).

The law is stated for \cle{point masses}. Newton proved (using his newly invented calculus) that a \emph{uniform spherical body} attracts external objects exactly as if all its mass were concentrated at its centre. This is why we may treat the Earth, the Sun and the planets as point masses located at their centres.

\begin{attention}[Distance measured between centres, not surfaces]
In $F = GMm/r^{2}$, the distance $r$ is measured \textbf{between the centres} of the two bodies, never between their surfaces. For a satellite at altitude $h$ above the Earth, $r = R_{E} + h$, where $R_{E} = 6.37\times 10^{6}\ \mathrm{m}$ is the Earth's radius. Forgetting to add $R_{E}$ is one of the most frequent errors at the GCE Advanced Level.
\end{attention}

\begin{popfigure}
\begin{center}
\begin{tikzpicture}[scale=1.0]
  % mass M
  \shade[ball color=popblue] (0,0) circle (0.55);
  \node[below=0.65cm of {(0,0)}, popink] {$M$};
  % mass m
  \shade[ball color=poporange] (6,0) circle (0.35);
  \node[below=0.5cm of {(6,0)}, popink] {$m$};
  % distance line
  \draw[dashed, popmuted] (0,0) -- (6,0);
  \draw[<->, popmuted] (0,-0.9) -- (6,-0.9) node[midway, below, popink] {$r$ (centre to centre)};
  % force vectors
  \draw[->, very thick, popdark] (0.6,0) -- (2.1,0) node[midway, above, popink] {$\vec{F}_{m\to M}$};
  \draw[->, very thick, popdark] (5.4,0) -- (3.9,0) node[midway, above, popink] {$\vec{F}_{M\to m}$};
\end{tikzpicture}
\end{center}
\legende{Two point masses $M$ and $m$ separated by $r$. The gravitational forces form an action--reaction pair: equal in magnitude, opposite in direction, always attractive.}
\end{popfigure}

\textbf{Vector nature and superposition.}\par
Gravitation is a \emph{vector} force. On mass $m$ at position $\vec{r}$ relative to mass $M$,
\[
\vec{F} = -\frac{GMm}{r^{2}}\,\hat{r},
\]
where the minus sign and the unit vector $\hat{r}$ (pointing from $M$ towards $m$) express that the force is always \cle{attractive}, directed along the line joining the centres. The two forces on $M$ and on $m$ form a Newton's-third-law pair: equal magnitudes, opposite directions. When several masses are present, the total force on one body is the \emph{vector sum} of the individual forces (principle of \cle{superposition}). For example, the force on the Moon is the vector sum of the pulls of the Earth and the Sun.

\textbf{The constant $G$ and its measurement.}\par
$G$ is a \emph{universal} constant: it has the same value everywhere in the universe, for all pairs of masses. Its unit follows directly from the law: $[G] = [F][r^{2}]/[m^{2}]$, that is $\mathrm{N\,m^{2}\,kg^{-2}}$. Its smallness ($\sim 10^{-11}$) explains why gravitational attraction between everyday objects is utterly negligible compared with the pull of the whole Earth. $G$ was first measured by \cle{Cavendish} (1798) with a torsion balance: two small lead spheres are hung from a fine fibre, two large lead spheres are brought near them, and the tiny twist of the fibre --- calibrated against a known restoring torque --- gives the attractive force, hence $G$.

\begin{exoresolu}[Everyday masses versus astronomical masses]
\textbf{Statement.} (a) Calculate the gravitational force between two students of masses $60\ \mathrm{kg}$ and $55\ \mathrm{kg}$ standing $1.0\ \mathrm{m}$ apart. (b) Calculate the force between the Earth ($M_{E} = 5.97\times 10^{24}\ \mathrm{kg}$) and the same $60\ \mathrm{kg}$ student standing on its surface ($R_{E} = 6.37\times 10^{6}\ \mathrm{m}$). Comment.
\tcblower
\textbf{Solution.}
(a) Treating the students as point masses:
\[
F = \frac{G m_{1} m_{2}}{r^{2}} = \frac{(6.67\times 10^{-11})(60)(55)}{(1.0)^{2}}
   = 2.2\times 10^{-7}\ \mathrm{N}.
\]
This is about the weight of a grain of sand --- impossible to feel.
(b) The student is at distance $r = R_{E}$ from the Earth's centre:
\[
F = \frac{(6.67\times 10^{-11})(5.97\times 10^{24})(60)}{(6.37\times 10^{6})^{2}}
  = \frac{2.39\times 10^{16}}{4.06\times 10^{13}}
  = 5.9\times 10^{2}\ \mathrm{N}.
\]
This is simply the student's weight, $mg \approx 60 \times 9.81 = 5.9\times 10^{2}\ \mathrm{N}$. The Earth's enormous mass compensates the smallness of $G$.
\end{exoresolu}

\courssub{Weight and the link between $g$, $G$, $M_{E}$ and $R_{E}$}

The \cle{weight} of a body of mass $m$ at the Earth's surface is nothing other than the gravitational force exerted on it by the whole Earth:
\[
mg = \frac{G M_{E} m}{R_{E}^{2}}
\quad\Longrightarrow\quad
\boxed{\,g = \frac{G M_{E}}{R_{E}^{2}}\,}.
\]
This relation is extremely useful: knowing $g$, $G$ and $R_{E}$, Cavendish's experiment effectively allowed scientists to ``weigh the Earth'', since $M_{E} = gR_{E}^{2}/G$. Numerically, $g = (6.67\times 10^{-11})(5.97\times 10^{24})/(6.37\times 10^{6})^{2} \approx 9.81\ \mathrm{m\,s^{-2}}$, as expected.

\begin{aretenir}[$G$ is universal, $g$ is local]
$G = 6.67\times 10^{-11}\ \mathrm{N\,m^{2}\,kg^{-2}}$ has the same value everywhere in the universe. The gravitational field strength $g$ is \emph{local}: it depends on the mass and radius of the attracting body and on the position. At the Earth's surface $g \approx 9.81\ \mathrm{m\,s^{-2}}$; on the Moon $g \approx 1.6\ \mathrm{m\,s^{-2}}$; at altitude, $g$ decreases. Never confuse the two symbols in an examination.
\end{aretenir}

\courssub{Kepler's laws of planetary motion}

Before Newton, Kepler had summarised Brahe's observations in three laws. Newton later showed that all three follow from $F = GMm/r^{2}$ --- a triumph of the universal law.

\begin{propriete}[Kepler's three laws]
\begin{enumerate}
  \item \textbf{Law of orbits:} each planet moves on an \cle{ellipse} with the Sun at one \emph{focus}.
  \item \textbf{Law of areas:} the line joining the Sun to a planet sweeps out \emph{equal areas in equal times}. Consequently a planet moves \emph{faster} near the Sun (perihelion) and \emph{slower} far from it (aphelion).
  \item \textbf{Law of periods:} the square of the orbital period is proportional to the cube of the orbit's semi-major axis:
  \[
  T^{2} = \frac{4\pi^{2}}{GM_{S}}\,r^{3},
  \]
  where $M_{S}$ is the mass of the central body (the Sun for planets, the Earth for satellites).
\end{enumerate}
\end{propriete}

\begin{popfigure}
\begin{center}
\begin{tikzpicture}[scale=1.1]
  % ellipse
  \draw[thick, popblue] (0,0) ellipse (3.2 and 2.0);
  % foci
  \fill[popgold] (1.8,0) circle (0.12);
  \node[below, popink] at (1.8,-0.15) {Sun (focus)};
  \fill[popmuted] (-1.8,0) circle (0.05);
  \node[below, popmuted] at (-1.8,-0.1) {empty focus};
  % planet positions
  \fill[poporange] (3.2,0) circle (0.09);
  \node[right, popink] at (3.25,0) {aphelion (slow)};
  \fill[poporange] (-3.2,0) circle (0.09);
  \node[left, popink] at (-3.25,0) {perihelion (fast)};
  % swept areas
  \fill[popgreenL, opacity=0.8] (1.8,0) -- (-3.2,0) arc (180:215:3.2 and 2.0) -- cycle;
  \fill[popgreenL, opacity=0.8] (1.8,0) -- (3.2,0) arc (0:-22:3.2 and 2.0) -- cycle;
  \node[popink] at (-2.2,-0.7) {$A_{1}$};
  \node[popink] at (2.3,-0.55) {$A_{2}$};
  \node[popink] at (0,-2.6) {$A_{1} = A_{2}$ for equal time intervals $\Delta t$};
\end{tikzpicture}
\end{center}
\legende{Kepler's first and second laws: the orbit is an ellipse with the Sun at one focus; the radius vector sweeps equal areas in equal times, so the planet travels faster near perihelion.}
\end{popfigure}

\begin{methode}[Deriving Kepler's third law for a circular orbit (examinable)]
For a circular orbit of radius $r$ around a central mass $M$:
\begin{enumerate}
  \item The gravitational force provides the \emph{centripetal force}:
  \[
  \frac{GMm}{r^{2}} = \frac{mv^{2}}{r}.
  \]
  \item Simplify: $v^{2} = \dfrac{GM}{r}$, so the orbital speed is $v = \sqrt{GM/r}$. Note that the satellite's mass $m$ cancels: the orbit depends only on the \emph{central} mass.
  \item Write the speed in terms of the period: $v = \dfrac{2\pi r}{T}$.
  \item Substitute and square:
  \[
  \frac{4\pi^{2}r^{2}}{T^{2}} = \frac{GM}{r}
  \quad\Longrightarrow\quad
  T^{2} = \frac{4\pi^{2}}{GM}\,r^{3}.
  \]
\end{enumerate}
Hence $T^{2}/r^{3} = 4\pi^{2}/(GM)$ is the \emph{same constant} for all bodies orbiting the same central mass.
\end{methode}

\begin{attention}[SI units in Kepler's third law]
Before substituting in $T^{2} = 4\pi^{2}r^{3}/(GM)$, convert $r$ to \textbf{metres} and expect $T$ in \textbf{seconds}. Only convert to minutes, hours or days \emph{after} the calculation. Mixing kilometres with $G$ in SI units is a classic source of wrong answers.
\end{attention}

\begin{popfigure}
\begin{center}
\begin{tikzpicture}
\begin{loglogaxis}[
    width=0.72\linewidth, height=6.2cm,
    xlabel={$r^{3}$ / $\mathrm{m^{3}}$},
    ylabel={$T^{2}$ / $\mathrm{s^{2}}$},
    grid=both, grid style={popline, very thin},
    legend style={at={(0.03,0.97)}, anchor=north west, font=\small},
]
\addplot[domain=1e32:1e40, samples=50, thick, popblue] {3.0e-19*x};
\addlegendentry{$T^{2}=\dfrac{4\pi^{2}}{GM_{S}}r^{3}$}
\addplot[only marks, mark=*, mark size=1.6pt, poporange] coordinates {
(2.0e32,5.9e13)
  (1.4e33,3.7e14)
  (7.2e33,9.7e14)
  (1.6e34,3.4e15)
  (4.0e35,1.4e17)
  (3.3e36,8.7e17)
  (1.1e37,6.2e18)
  (5.6e37,3.1e19)
};
\addlegendentry{planets (Mercury to Neptune)}
\end{loglogaxis}
\end{tikzpicture}
\end{center}
\legende{Log--log plot of $T^{2}$ against $r^{3}$ for the planets of the solar system. All points lie on a straight line of slope $4\pi^{2}/(GM_{S})$, confirming Kepler's third law.}
\end{popfigure}

\begin{exoresolu}[Period of the Moon from Kepler's third law]
\textbf{Statement.} The Moon orbits the Earth at a mean distance $r = 3.84\times 10^{8}\ \mathrm{m}$ from the Earth's centre. Taking $M_{E} = 5.97\times 10^{24}\ \mathrm{kg}$, calculate its orbital period in days.
\tcblower
\textbf{Solution.} Apply Kepler's third law with the Earth as central body:
\begin{align*}
T^{2} &= \frac{4\pi^{2} r^{3}}{G M_{E}}
      = \frac{4\pi^{2}\,(3.84\times 10^{8})^{3}}{(6.67\times 10^{-11})(5.97\times 10^{24})}\\
      &= \frac{39.48 \times 5.67\times 10^{25}}{3.98\times 10^{14}}
      = 5.62\times 10^{12}\ \mathrm{s^{2}}.\\
T &= 2.37\times 10^{6}\ \mathrm{s} = \frac{2.37\times 10^{6}}{86\,400}\ \text{days} \approx 27.4\ \text{days}.
\end{align*}
This matches the known sidereal period of the Moon --- a beautiful confirmation of the law.
\end{exoresolu}

\begin{exercice}[Force and period]
\begin{enumerate}
  \item Calculate the gravitational force between the Sun ($M_{S} = 1.99\times 10^{30}\ \mathrm{kg}$) and the Earth ($M_{E} = 5.97\times 10^{24}\ \mathrm{kg}$), separated by $r = 1.50\times 10^{11}\ \mathrm{m}$.
  \item A satellite orbits the Earth at $r = 7.0\times 10^{6}\ \mathrm{m}$ from the centre. Using Kepler's third law, find its period in minutes.
\end{enumerate}
\end{exercice}

\begin{corrige}[Exercise 1]
\begin{enumerate}
  \item $F = \dfrac{(6.67\times 10^{-11})(1.99\times 10^{30})(5.97\times 10^{24})}{(1.50\times 10^{11})^{2}} = \dfrac{7.93\times 10^{44}}{2.25\times 10^{22}} = 3.52\times 10^{22}\ \mathrm{N}$.
  \item $T^{2} = \dfrac{4\pi^{2}(7.0\times 10^{6})^{3}}{(6.67\times 10^{-11})(5.97\times 10^{24})} = \dfrac{1.354\times 10^{22}}{3.98\times 10^{14}} = 3.40\times 10^{7}\ \mathrm{s^{2}}$, so $T = 5.83\times 10^{3}\ \mathrm{s} \approx 97\ \mathrm{min}$.
\end{enumerate}
\end{corrige}

\begin{aretenir}[Key points of this section]
\begin{itemize}
  \item Newton's law: $F = GMm/r^{2}$, attractive, along the line of centres; $r$ is measured \emph{centre to centre}. The force obeys an \emph{inverse-square law}.
  \item $G$ is universal; $g = GM_{E}/R_{E}^{2}$ is local.
  \item Kepler: elliptical orbits (Sun at a focus); equal areas in equal times; $T^{2} = 4\pi^{2}r^{3}/(GM)$, derived by equating gravitation to centripetal force.
\end{itemize}
\end{aretenir}

\coursec{Gravitational Field Strength and Variation of $g$}

In the previous section we established Newton's law of universal gravitation: any two masses attract each other with a force $F = GMm/r^{2}$, and we saw how this single law explains Kepler's empirical laws of planetary motion. We now change our point of view. Instead of always speaking of a \emph{pair} of masses, we focus on \emph{one} mass --- the Earth --- and describe the ``influence'' it creates in the whole space around it. This influence is what physicists call a \cle{gravitational field}. The idea is powerful: once we know the field, we can predict the force on \emph{any} object placed at \emph{any} point, without referring back to the Earth each time.

\courssub{Qualitative Description of the Earth's Gravitational Field}

Drop a stone anywhere near the Earth's surface: it falls towards the ground. Release a stone above Buea, above Douala, or above the ocean: in every case it accelerates towards the \emph{centre of the Earth}. The Moon, far away, is also continuously ``falling'' towards the Earth, which is what keeps it in orbit. These observations show that:

\begin{itemize}
\item the Earth exerts an attractive influence on every mass around it, near or far;
\item this influence is directed towards the centre of the Earth at every point;
\item the influence becomes weaker as we move further away.
\end{itemize}

We picture this influence using \cle{field lines}. Because the Earth is (to a good approximation) spherical, the field lines are \textbf{radial}, like the spokes of a wheel, and they are always directed \textbf{inwards}, towards the centre, because gravity is always attractive. Where the lines are close together (near the surface), the field is strong; where they spread out (far away), the field is weak.

\begin{popfigure}
\begin{center}
\begin{tikzpicture}[scale=1.0]
% Earth
\shade[ball color=popblue] (0,0) circle (0.9);
\node[white] at (0,0) {\textbf{Earth}};
% radial field lines with inward arrows
\foreach \a in {0,45,90,135,180,225,270,315}{
  \draw[->, thick, poporange] ({3.0*cos(\a)},{3.0*sin(\a)}) -- ({1.0*cos(\a)},{1.0*sin(\a)});
}
% labels
\node[popdark, align=center] at (0,-3.6) {field lines point inwards};
\node[popdark, align=center, text width=2.6cm] at (4.3,1.6) {lines far apart:\\ weak field};
\node[popdark, align=center, text width=2.6cm] at (-4.3,-1.6) {lines close together:\\ strong field};
\end{tikzpicture}
\end{center}
\legende{Radial gravitational field lines around the Earth. Every arrow points towards the centre; the density of lines decreases with distance, showing that the field weakens as $1/r^{2}$.}
\end{popfigure}

\begin{attention}[Field lines are a picture, not a physical object]
Nothing material exists along a field line. Field lines are a \emph{representation} invented to visualise the direction and relative strength of the field. The field itself exists at \emph{every} point of space, whether or not a line is drawn through it.
\end{attention}

\courssub{Definition of Gravitational Field Strength}

Suppose we want to compare the strength of gravity at two different places, say at sea level and at the summit of Mount Cameroon. We cannot simply compare the \emph{forces} on two objects, because the force depends on the object's mass: a heavy lorry is pulled harder than a small stone, even at the same place. To characterise the field \emph{itself}, we must remove the mass of the test object from the comparison. This leads to the following definition.

\begin{definition}[Gravitational field strength]
The \cle{gravitational field strength} $\vec{g}$ at a point is the gravitational force per unit mass experienced by a small test mass $m$ placed at that point:
\[
\vec{g} = \frac{\vec{F}}{m} \qquad \text{or in magnitude} \qquad g = \frac{F}{m}.
\]
Its SI unit is the newton per kilogram, $\SI{1}{N.kg^{-1}}$. Since $\SI{1}{N} = \SI{1}{kg.m.s^{-2}}$, we have the equivalent unit
\[
\SI{1}{N.kg^{-1}} = \SI{1}{m.s^{-2}},
\]
which shows that $g$ is none other than the \textbf{acceleration of free fall} at that point. The vector $\vec{g}$ always points in the direction of the force, i.e.\ towards the centre of the attracting mass.
\end{definition}

\begin{aretenir}
$g$ is a property of the \emph{point in space}, not of the object placed there. A feather and a hammer at the same point experience the same $g$ (and, in vacuum, fall with the same acceleration), even though the forces on them are very different.
\end{aretenir}

\courssub{Field Strength Due to a Point or Spherical Mass}

Consider a point mass $M$. Place a test mass $m$ at a distance $r$ from it. By Newton's law of gravitation, the force on $m$ is
\[
F = \frac{GMm}{r^{2}}.
\]
Applying the definition $g = F/m$, the mass $m$ cancels:

\begin{propriete}[Field of a point or spherical mass]
The gravitational field strength at distance $r$ from a point mass $M$ --- or \emph{outside} a uniform spherical mass $M$, at distance $r$ from its centre --- is
\[
g = \frac{GM}{r^{2}}.
\]
The field is radial and directed towards $M$. (The proof for a sphere requires integration over the whole sphere and is \emph{not} examinable; the result, due to Newton, is that a uniform sphere attracts external bodies exactly as if all its mass were concentrated at its centre.)
\end{propriete}

Applying this to the Earth, of mass $M_{E}$ and radius $R_{E}$, the field strength at its surface is
\[
g_{0} = \frac{GM_{E}}{R_{E}^{2}}.
\]

\begin{exemplebox}[The Earth's surface field]
With $G = \SI{6.67e-11}{N.m^{2}.kg^{-2}}$, $M_{E} = \SI{5.98e24}{kg}$ and $R_{E} = \SI{6.37e6}{m}$:
\[
g_{0} = \frac{(6.67\times10^{-11})(5.98\times10^{24})}{(6.37\times10^{6})^{2}} \approx \SI{9.83}{N.kg^{-1}}.
\]
This is the familiar acceleration of free fall, $g_{0} \approx \SI{9.81}{m.s^{-2}}$.
\end{exemplebox}

\courssub{The Uniform Field Approximation Near the Surface}

Since $g = GM_{E}/r^{2}$ changes with $r$, the field is strictly \emph{not} uniform. However, near the surface the change is tiny: climbing from sea level to the top of Mount Cameroon ($h \approx \SI{4.1}{km}$) changes $r$ from $\SI{6371}{km}$ to $\SI{6375}{km}$, a relative change of less than $0.1\%$. Over such small heights we may treat $\vec{g}$ as \textbf{constant in magnitude and direction}: the field is approximately \cle{uniform}, with parallel, equally spaced field lines pointing ``downwards''. This is exactly the approximation used in Lower Sixth projectile motion, and we shall use it again in Section 4.

\courssub{Variation of $g$ with Altitude (Outside the Earth)}

At altitude $h$ above the surface, the distance from the centre is $r = R_{E} + h$. Hence
\[
g(h) = \frac{GM_{E}}{(R_{E}+h)^{2}}.
\]
Dividing by $g_{0} = GM_{E}/R_{E}^{2}$ eliminates $G$ and $M_{E}$:

\begin{propriete}[Variation of $g$ with altitude]
\[
\boxed{\;g(h) = g_{0}\,\frac{R_{E}^{2}}{(R_{E}+h)^{2}}\;}
\]
For small altitudes ($h \ll R_{E}$), the binomial approximation $(1+x)^{-2} \approx 1 - 2x$ gives
\[
g(h) \approx g_{0}\left(1 - \frac{2h}{R_{E}}\right).
\]
\end{propriete}

\begin{methode}[Computing $g$ at altitude using ratios]
\begin{enumerate}
\item Write the ratio $\dfrac{g}{g_{0}} = \dfrac{R_{E}^{2}}{(R_{E}+h)^{2}}$ --- no need for $G$ or $M_{E}$.
\item Convert $h$ and $R_{E}$ to the \emph{same unit} (usually metres).
\item Compute the ratio, then multiply by $g_{0} = \SI{9.81}{m.s^{-2}}$.
\item Check: if $h \ll R_{E}$, verify with $g \approx g_{0}(1 - 2h/R_{E})$.
\end{enumerate}
\end{methode}

\begin{exoresolu}[Field strength at the summit of Mount Cameroon and at satellite altitude]
Take $R_{E} = \SI{6.37e6}{m}$ and $g_{0} = \SI{9.81}{m.s^{-2}}$.\\
\textbf{(a)} Find $g$ at the summit of Mount Cameroon, $h = \SI{4.10e3}{m}$.\\
\textbf{(b)} Find $g$ at the altitude of a geostationary satellite, $h = \SI{3.58e7}{m}$.
\tcblower
\textbf{(a)} Using the ratio method:
\[
\frac{g}{g_{0}} = \left(\frac{R_{E}}{R_{E}+h}\right)^{2} = \left(\frac{6.37\times10^{6}}{6.37\times10^{6} + 4.10\times10^{3}}\right)^{2} = \left(\frac{6.370}{6.3741}\right)^{2} \approx 0.9987.
\]
\[
g \approx 9.81 \times 0.9987 \approx \SI{9.80}{m.s^{-2}}.
\]
Check with the approximation: $g \approx g_{0}(1 - 2h/R_{E}) = 9.81\left(1 - \dfrac{2 \times 4.10\times10^{3}}{6.37\times10^{6}}\right) \approx 9.81(1 - 0.00129) \approx \SI{9.80}{m.s^{-2}}$. \checkmark\ The decrease is only about $\SI{0.013}{m.s^{-2}}$: imperceptible in everyday life.

\textbf{(b)} Here $h$ is \emph{not} small compared with $R_{E}$, so the approximation is forbidden:
\[
\frac{g}{g_{0}} = \left(\frac{6.37\times10^{6}}{6.37\times10^{6} + 3.58\times10^{7}}\right)^{2} = \left(\frac{6.37}{42.17}\right)^{2} \approx (0.1511)^{2} \approx 0.0228.
\]
\[
g \approx 9.81 \times 0.0228 \approx \SI{0.22}{m.s^{-2}}.
\]
At geostationary altitude the Earth's pull is only about $2.3\%$ of its surface value --- yet it is exactly enough to keep the satellite in orbit, as we shall see in Section 3.
\end{exoresolu}

\courssub{Variation of $g$ with Depth (Inside the Earth)}

What happens if we descend into a deep mine shaft, at depth $d$ below the surface? The distance from the centre is now $r = R_{E} - d$, which is \emph{less} than $R_{E}$, so one might naively expect $g$ to \emph{increase} by the inverse-square law. This is wrong, and the reason is instructive.

When the test mass is \emph{inside} the Earth, the shell of material lying \emph{above} it (at radii greater than $r$) pulls it in all directions, and it can be shown that the contributions of a uniform spherical shell cancel exactly: \textbf{a uniform shell exerts no net force on a mass inside it}. Only the sphere of radius $r$ \emph{below} the mass contributes. Assuming the Earth has \cle{uniform density} $\rho$, this inner mass is
\[
M_{\text{in}} = \rho \cdot \frac{4}{3}\pi r^{3}.
\]
Therefore
\[
g(r) = \frac{G M_{\text{in}}}{r^{2}} = \frac{G}{r^{2}}\cdot \rho \frac{4}{3}\pi r^{3} = \frac{4}{3}\pi G \rho\, r.
\]

\begin{propriete}[Variation of $g$ inside the Earth (uniform density model)]
Inside a uniform spherical Earth, the field strength is \emph{proportional} to the distance from the centre:
\[
g(r) = \frac{4}{3}\pi G\rho\, r \qquad (r \le R_{E}),
\]
or, in terms of the depth $d = R_{E} - r$ and the surface value $g_{0}$:
\[
\boxed{\;g(d) = g_{0}\left(1 - \frac{d}{R_{E}}\right)\;}
\]
At the centre ($d = R_{E}$, $r = 0$): $g = 0$. (The cancellation of the shell's pull is not examinable, but the linear result and its derivation from $M_{\text{in}} = \frac{4}{3}\pi\rho r^{3}$ are.)
\end{propriete}

\begin{attention}[$g$ is maximum at the surface, not at the centre]
A very common mistake is to think that ``closer to the centre means stronger gravity''. In fact, as you descend, there is less and less mass below you and the pull of the shell above cancels out. $g$ \emph{decreases linearly} with depth and is exactly \textbf{zero at the centre}: a mass placed there is pulled equally in all directions. The maximum value of $g$ is at the \emph{surface}.
\end{attention}

\begin{exemplebox}[$g$ at the centre of the Earth]
At the centre, $d = R_{E}$, so $g = g_{0}(1 - 1) = \SI{0}{m.s^{-2}}$. Halfway down ($d = R_{E}/2$), $g = g_{0}/2 \approx \SI{4.9}{m.s^{-2}}$.
\end{exemplebox}

\courssub{Graphical Summary: $g$ Against Distance from the Centre}

Combining the two results:
\[
g(r) = \begin{cases} g_{0}\,\dfrac{r}{R_{E}} & \text{for } r \le R_{E} \quad \text{(linear rise)},\\[6pt] g_{0}\,\dfrac{R_{E}^{2}}{r^{2}} & \text{for } r \ge R_{E} \quad \text{(inverse-square fall)}. \end{cases}
\]

\begin{popfigure}
\begin{center}
\begin{tikzpicture}
\begin{axis}[
  width=11cm, height=7cm,
  xlabel={$r/R_{E}$}, ylabel={$g/g_{0}$},
  xmin=0, xmax=3, ymin=0, ymax=1.15,
  xtick={0,1,2,3}, ytick={0,0.5,1},
  grid=both, grid style={popline, very thin},
  axis line style={popdark},
  legend style={draw=none, at={(0.97,0.95)}, anchor=north east, font=\small},
]
\addplot[domain=0:1, samples=50, very thick, popblue] {x};
\addlegendentry{inside: $g \propto r$}
\addplot[domain=1:3, samples=100, very thick, poporange] {1/x^2};
\addlegendentry{outside: $g \propto 1/r^{2}$}
\addplot[mark=*, mark size=1.5pt, popdark, only marks] coordinates {(1,1)};
\node[popdark, anchor=south west, font=\small] at (axis cs:1.03,0.98) {surface: maximum};
\end{axis}
\end{tikzpicture}
\end{center}
\legende{Variation of $g$ with distance $r$ from the centre of a uniform-density Earth: linear increase inside, inverse-square decrease outside, maximum at the surface $r = R_{E}$.}
\end{popfigure}

\begin{popfigure}
\begin{center}
\begin{tikzpicture}[scale=0.85]
% Earth cross-section
\fill[popblueL] (0,0) circle (2);
\draw[thick, popblue] (0,0) circle (2);
\fill[popink] (0,0) circle (0.04);
\node[popink, below right=1pt] at (0,0) {centre};
% radius
\draw[->, thick, popdark] (0,0) -- (2,0) node[midway, below] {$R_{E}$};
% altitude point
\fill[poporange] (3.1,0) circle (0.06);
\draw[<->, thick, poporange] (2.05,0.25) -- (3.05,0.25) node[midway, above] {$h$};
\node[poporange, right] at (3.15,0) {satellite};
\draw[->, thick, poporange] (0,0) -- (3.1,-0.35) node[midway, below, sloped] {$r = R_{E}+h$};
% depth point
\fill[poppurple] (-1.2,0) circle (0.06);
\draw[<->, thick, poppurple] (-1.95,-0.3) -- (-1.25,-0.3) node[midway, below] {$d$};
\node[poppurple, above left] at (-1.2,0.1) {mine};
\draw[->, thick, poppurple] (0,0.3) -- (-1.2,0.3) node[midway, above] {$r = R_{E}-d$};
\node[popdark] at (0,-2.6) {altitude $h$ measured upwards from the surface; depth $d$ downwards};
\end{tikzpicture}
\end{center}
\legende{Conventions: outside the Earth $r = R_{E} + h$; inside, $r = R_{E} - d$. In every formula, $r$ is measured from the \emph{centre}.}
\end{popfigure}

\courssub{Effect of the Earth's Rotation and Shape on Apparent $g$}

The value of $g$ actually measured at the surface is not perfectly uniform, for two reasons (qualitative treatment only):

\begin{itemize}
\item \textbf{Shape.} The Earth is not a perfect sphere: it is slightly flattened at the poles and bulges at the equator, so $R_{E}$ is slightly larger at the equator. Since $g_{0} = GM_{E}/R_{E}^{2}$, the true field is slightly \emph{smaller} at the equator than at the poles.
\item \textbf{Rotation.} A body resting on the equator is carried in a circle by the Earth's rotation and needs a small centripetal acceleration $a = \omega^{2}R_{E}$. Part of the gravitational pull is ``used up'' to provide it, so the \emph{apparent} weight (what a spring balance reads) is reduced. At the poles there is no such effect.
\end{itemize}

Both effects act in the same sense: the \textbf{apparent $g$ is smallest at the equator and largest at the poles}. The difference is small (a few tenths of a per cent), which is why we usually quote a single average value $g_{0} \approx \SI{9.81}{m.s^{-2}}$.

\begin{savaistu}[Weighing less at the equator]
Because of the combined effects of rotation and the equatorial bulge, you would weigh slightly less in Douala (near the equator) than at the poles --- even without moving! Precise gravimeters easily detect this difference, and geophysicists use tiny local variations of $g$ to prospect for oil and minerals, since dense underground rocks slightly increase the local field.
\end{savaistu}

\begin{attention}[GCE exam tip]
When a question does \emph{not} give you $G$ and $M_{E}$, do not try to invent them. Express everything in terms of $g_{0}$ and $R_{E}$ using the ratio
\[
\frac{g}{g_{0}} = \frac{R_{E}^{2}}{r^{2}} \quad (\text{outside}), \qquad \frac{g}{g_{0}} = \frac{r}{R_{E}} \quad (\text{inside}).
\]
Also remember: $r$ is always measured from the \emph{centre} of the Earth, never from the surface. Writing $g = GM/h^{2}$ with $h$ the altitude is a classic error that costs marks.
\end{attention}

\begin{exercice}[Variation of $g$]
Take $R_{E} = \SI{6.37e6}{m}$ and $g_{0} = \SI{9.81}{m.s^{-2}}$.
\begin{enumerate}
\item At what altitude above the surface is $g$ equal to $g_{0}/4$?
\item A very deep mine reaches a depth $d = \SI{100}{km}$. Using the uniform-density model, find $g$ at the bottom of the mine.
\item Explain, without calculation, why $g$ at the centre of the Earth is zero even though the distance to the centre is zero.
\end{enumerate}
\end{exercice}

\begin{corrige}[Exercise 1]
\textbf{1.} Outside the Earth, $\dfrac{g}{g_{0}} = \dfrac{R_{E}^{2}}{(R_{E}+h)^{2}} = \dfrac{1}{4}$, so $\dfrac{R_{E}}{R_{E}+h} = \dfrac{1}{2}$, giving $R_{E} + h = 2R_{E}$, hence
\[
h = R_{E} = \SI{6.37e6}{m}.
\]
(Note: $g$ falls to a quarter when the \emph{distance from the centre} doubles, not the altitude.)

\textbf{2.} Inside the Earth, $g = g_{0}\left(1 - \dfrac{d}{R_{E}}\right) = 9.81\left(1 - \dfrac{1.0\times10^{5}}{6.37\times10^{6}}\right) = 9.81 \times 0.984 \approx \SI{9.66}{m.s^{-2}}$.

\textbf{3.} At the centre, the mass of the Earth surrounds the body symmetrically in all directions. The gravitational pulls from all the surrounding shells cancel each other exactly, so the net force --- and therefore $g = F/m$ --- is zero. The formula $g = GM/r^{2}$ does \emph{not} apply inside the Earth, because the mass is no longer entirely ``below'' the body.
\end{corrige}

\begin{aretenir}[Key points of this section]
\begin{itemize}
\item Gravitational field strength: $g = F/m$, unit $\mathrm{N\,kg^{-1}} = \mathrm{m\,s^{-2}}$; it is the acceleration of free fall.
\item Outside a spherical mass: $g = GM/r^{2}$ with $r$ measured from the \emph{centre}; at the surface $g_{0} = GM_{E}/R_{E}^{2}$.
\item With altitude: $g = g_{0}R_{E}^{2}/(R_{E}+h)^{2} \approx g_{0}(1 - 2h/R_{E})$ for $h \ll R_{E}$.
\item With depth (uniform density): $g = g_{0}(1 - d/R_{E})$; $g$ falls linearly to zero at the centre.
\item $g$ is \textbf{maximum at the surface}; rotation and the equatorial bulge make apparent $g$ smallest at the equator.
\end{itemize}
\end{aretenir}

\coursec{Gravitational Potential, Potential Energy, Escape and Orbital Speeds}

In the previous section we described the Earth's gravitational field through the \cle{field strength} $\vec{g}$, a \emph{force per unit mass}. Force is not the only way to describe a field: it is often far simpler to work with \emph{energy}. Just as knowing the height of a hill tells you how much energy a cyclist needs to climb it, knowing the \cle{gravitational potential} at a point tells you how much energy a mass needs to get there. In this section we build the energy description of gravitation, and we use it to answer three questions of great practical importance: how fast must a rocket go to escape the Earth, how fast must a satellite go to stay in orbit, and where must a satellite be placed so that it appears fixed above Yaoundé or Douala?

\courssub{Work Done by a Gravitational Force: a Conservative Field}

\begin{experience}[Lifting a mass away from the Earth]
Imagine a small mass $m$ resting on the ground. To raise it slowly to a height $h$ you must pull it upwards against the Earth's attraction: you do \emph{positive} work, and the gravitational force does \emph{negative} work. If you then release the mass, gravity pulls it back down and returns exactly the energy you stored. No energy is lost along the way. This is the signature of a \cle{conservative force}: the work it does depends only on the starting and finishing points, never on the path taken.
\end{experience}

Consider a mass $m$ moved from point $A$ to point $B$ in the field of a mass $M$ (the Earth). The gravitational force $\vec{F} = -\dfrac{GMm}{r^{2}}\,\vec{e}_{r}$ varies with distance, so the work must be computed by integration along a radial path:
\[
W_{A\to B} = \int_{A}^{B} \vec{F}\cdot d\vec{r} = \int_{r_A}^{r_B} -\frac{GMm}{r^{2}}\,dr = GMm\left(\frac{1}{r_B} - \frac{1}{r_A}\right).
\]
The result depends only on $r_A$ and $r_B$. Whatever path is chosen — straight, curved, zig-zag — the work is the same, and the work over any closed path is zero.

\begin{propriete}[The gravitational field is conservative]
The gravitational force is a \cle{conservative force}: its work between two points is independent of the path followed. Consequently a \emph{potential energy} $E_p$ can be defined such that
\[
W_{A\to B}(\vec{F}_{\text{grav}}) = -\Delta E_p = E_{p,A} - E_{p,B},
\]
and the total mechanical energy $E = E_k + E_p$ of a body moving under gravity alone is conserved.
\end{propriete}

\courssub{Gravitational Potential Energy}

To define potential energy we must choose a reference level where $E_p = 0$. For objects near the ground we used ``$E_p = 0$ at the Earth's surface'', but for satellites and rockets the natural choice is a point \emph{infinitely far away}, where the gravitational force vanishes. Bringing a mass $m$ from infinity to a distance $r$ from $M$, the work done \emph{by gravity} is
\[
W_{\infty \to r} = \frac{GMm}{r} - 0 = \frac{GMm}{r} = -\left(E_p(r) - E_p(\infty)\right),
\]
so, with $E_p(\infty) = 0$:

\begin{definition}[Gravitational potential energy]
The \cle{gravitational potential energy} of a mass $m$ at a distance $r$ from the centre of a spherical mass $M$ is
\[
\boxed{E_p = -\frac{GMm}{r}} \qquad (\text{zero chosen at infinity}).
\]
$E_p$ is measured in joules (J) and is a property of the \emph{system} $(M, m)$, not of $m$ alone.
\end{definition}

\textbf{Why the negative sign?} The minus sign is not a mathematical accident; it carries deep physical meaning. At infinity the system has $E_p = 0$ and the two bodies are \emph{free} of each other. Any finite separation gives $E_p < 0$: the system has \emph{less} energy than a free system. To separate the bodies to infinity you must \emph{supply} energy $\dfrac{GMm}{r}$. We say the mass $m$ is \cle{bound} to $M$, trapped in a ``potential well''. A satellite in orbit, the Moon around the Earth, the Earth around the Sun: all have negative total energy, all are bound.

\begin{attention}[Potential energy is negative — do not ``correct'' it]
Many students drop the minus sign in $E_p = -GMm/r$ because ``energy cannot be negative''. It can! Only \emph{differences} of potential energy are physical, and the zero is a convention. A satellite in orbit is bound: \emph{raising} it to a higher orbit (larger $r$) makes $E_p$ less negative, which requires \emph{positive} work from the rocket engines. Never write $E_p = +GMm/r$.
\end{attention}

\courssub{Gravitational Potential}

Potential energy depends on the test mass $m$. To characterise the field itself, at each point of space, we divide by $m$ — exactly as we divided force by mass to obtain $g$.

\begin{definition}[Gravitational potential]
The \cle{gravitational potential} at a distance $r$ from the centre of a spherical mass $M$ is the potential energy per unit mass placed at that point:
\[
\boxed{V = \frac{E_p}{m} = -\frac{GM}{r}} \qquad \text{unit: } \mathrm{J\,kg^{-1}}.
\]
The zero of potential is at infinity. Potentials from several masses add algebraically (scalars — no vector addition needed!).
\end{definition}

The \cle{potential difference} between two points gives the work directly: moving a mass $m$ from $A$ to $B$ requires an external agent to supply
\[
W_{\text{ext}} = \Delta E_p = m\,\Delta V = m\left(V_B - V_A\right).
\]

\begin{exemplebox}[Potential at the Earth's surface]
With $M_E = 5.98\times 10^{24}\ \mathrm{kg}$, $R_E = 6.37\times 10^{6}\ \mathrm{m}$:
\[
V_{\text{surface}} = -\frac{GM_E}{R_E} = -\frac{6.67\times 10^{-11}\times 5.98\times 10^{24}}{6.37\times 10^{6}} \approx -6.26\times 10^{7}\ \mathrm{J\,kg^{-1}}.
\]
Every kilogramme sitting on the ground is ``down'' by about $63\ \mathrm{MJ}$ compared with a kilogramme at infinity. That is the depth of the Earth's potential well.
\end{exemplebox}

\textbf{Equipotential surfaces.} Around a spherical mass, $V$ depends only on $r$, so the surfaces of constant potential are concentric \emph{spheres} centred on $M$. The field lines (radial, pointing inward) are everywhere perpendicular to the equipotentials, and no work is needed to move a mass \emph{along} an equipotential surface. Close to the Earth the equipotentials are closely spaced (strong field); far away they spread apart (weak field).

\begin{popfigure}
\begin{center}
\begin{tikzpicture}
\begin{axis}[
  width=12cm, height=7.5cm,
  xlabel={$r$ (arbitrary units)}, ylabel={$V(r)$, $E_p(r)$ (negative)},
  xmin=0.8, xmax=10, ymin=-1.15, ymax=0.15,
  grid=both, grid style={popline, dashed},
  legend style={at={(0.97,0.35)}, anchor=east, font=\small},
  axis lines=left, thick
]
\addplot[domain=1:10, samples=120, very thick, popblue] {-1/x};
\addlegendentry{$V(r) = -GM/r$}
\addplot[domain=1:10, samples=120, very thick, poporange, dashed] {-2/x};
\addlegendentry{$E_p(r) = -GMm/r$ (with $m=2$ units)}
\addplot[domain=0.8:10, samples=2, thin, gray] {0};
\node[font=\small, popmuted, align=center] at (axis cs:6.5,-0.12) {both tend to $0$\\ as $r \to \infty$};
\end{axis}
\end{tikzpicture}
\end{center}
\legende{Gravitational potential $V(r)$ and potential energy $E_p(r)$: both are negative everywhere and rise towards zero at infinity. The system is bound inside this ``potential well''.}
\end{popfigure}

\courssub{Variation of Potential Inside the Earth (Extension)}

The formulae $V = -GM/r$ and $g = GM/r^2$ apply \emph{outside} a spherical mass ($r \ge R$). Inside a \emph{uniform} sphere of radius $R$ and mass $M$, the field strength grows linearly with $r$ and the potential is parabolic:

\begin{propriete}[Field and potential inside a uniform sphere]
For $r \le R$ (inside a uniform sphere of mass $M$, radius $R$):
\[
g(r) = \frac{GM}{R^3}\,r, \qquad
V(r) = -\frac{GM}{2R^3}\left(3R^2 - r^2\right).
\]
At the centre ($r=0$), $g=0$ and $V_{\text{centre}} = -\dfrac{3GM}{2R} = \dfrac{3}{2}V_{\text{surface}}$. The potential is continuous and smooth at $r=R$.
\end{propriete}

\begin{attention}[Syllabus focus]
The GCE Advanced Level examination focuses almost exclusively on the \emph{external} field ($r \ge R$). The internal formulae above are provided for completeness and deeper understanding; they are rarely examined directly.
\end{attention}

\courssub{Escape Velocity}

\textbf{Intuition.} Throw a stone upwards: it falls back. Throw it faster: it rises higher before falling back. Is there a speed for which it \emph{never} falls back? Yes — the \cle{escape velocity}, the minimum launch speed for which the projectile just reaches infinity with zero speed left.

\begin{propriete}[Escape velocity — proof examinable]
The \cle{escape velocity} from the surface of a spherical body of mass $M$ and radius $R$ is
\[
\boxed{v_e = \sqrt{\frac{2GM}{R}} = \sqrt{2g_0 R}}
\]
where $g_0 = GM/R^{2}$ is the field strength at the surface.
\end{propriete}

\textbf{Derivation (energy conservation).} Launch a body of mass $m$ from the surface with speed $v_e$, ignoring air resistance. Mechanical energy is conserved:
\[
E_{\text{surface}} = \tfrac{1}{2}mv_e^{2} - \frac{GMm}{R}.
\]
The \emph{minimum} escape condition is that the body arrives at infinity with zero kinetic energy (any leftover speed would mean we launched too fast). At infinity $E_k = 0$ and $E_p = 0$, so $E_{\infty} = 0$. Conservation gives
\[
\tfrac{1}{2}mv_e^{2} - \frac{GMm}{R} = 0
\quad\Longrightarrow\quad
v_e = \sqrt{\frac{2GM}{R}} = \sqrt{2g_0 R}.
\]

\begin{exemplebox}[Escape velocity of the Earth]
\[
v_e = \sqrt{2g_0 R_E} = \sqrt{2\times 9.81 \times 6.37\times 10^{6}} \approx 1.12\times 10^{4}\ \mathrm{m\,s^{-1}} \approx 11.2\ \mathrm{km\,s^{-1}}.
\]
Note that $v_e$ is independent of the mass of the escaping object: a stone and a rocket need the same \emph{speed} (but very different energies).
\end{exemplebox}

\begin{savaistu}[Why the Moon has no atmosphere]
Escape velocity depends on $M$ and $R$. For the Moon, $v_e \approx 2.4\ \mathrm{km\,s^{-1}}$, low enough that the fastest molecules of any gas can leak away over geological time. The Earth, with $v_e \approx 11.2\ \mathrm{km\,s^{-1}}$, comfortably retains its nitrogen and oxygen. The same physics explains why hydrogen, the lightest and fastest gas, is rare in our atmosphere.
\end{savaistu}

\courssub{Satellites in Circular Orbit: Speed, Period, Energy}

\textbf{Setting up the motion.} A satellite of mass $m$ in a circular orbit of radius $r$ around the Earth (mass $M_E$) is in uniform circular motion: the only force is gravity, which plays the role of the centripetal force.

\begin{popfigure}
\begin{center}
\begin{tikzpicture}[scale=1.0]
\shade[inner color=popblue, outer color=popblueL] (0,0) circle (0.9);
\node[white, font=\small\bfseries] at (0,0) {Earth};
\draw[thick, popmuted, dashed] (0,0) circle (2.6);
\fill[poporange] (2.6,0) circle (0.10);
\node[popdark, font=\small, anchor=west] at (2.75,0.25) {satellite $m$};
\draw[-{Stealth[length=3mm]}, very thick, popgreen] (2.6,0) -- (2.6,1.5)
  node[midway, right, font=\small] {$\vec{v}$ (tangent)};
\draw[-{Stealth[length=3mm]}, very thick, poppink] (2.6,0) -- (1.1,0)
  node[midway, above, font=\small] {$\vec{F}_g$};
\draw[-{Stealth[length=2mm]}, popmuted] (0.45,0) -- (2.15,0) node[midway, below, font=\small] {$r$};
\node[font=\small, popmuted, align=center] at (-2.6,2.2) {circular orbit\\ radius $r$};
\end{tikzpicture}
\end{center}
\legende{A satellite in circular orbit: the gravitational force $\vec{F}_g$ points to the centre and provides the centripetal force; the velocity $\vec{v}$ is tangent to the orbit.}
\end{popfigure}

Newton's second law projected on the radial (centripetal) direction gives
\[
\frac{GM_E m}{r^{2}} = \frac{mv^{2}}{r}
\quad\Longrightarrow\quad
\boxed{v = \sqrt{\frac{GM_E}{r}}}.
\]
The orbital speed depends only on $r$: the higher the orbit, the \emph{slower} the satellite. The period follows from $T = 2\pi r / v$:
\[
\boxed{T = 2\pi\sqrt{\frac{r^{3}}{GM_E}}}
\qquad\Longleftrightarrow\qquad
\frac{T^{2}}{r^{3}} = \frac{4\pi^{2}}{GM_E} \ \text{(constant)},
\]
which is exactly Kepler's third law, now \emph{derived} rather than merely observed.

\begin{propriete}[Energy of a circular orbit]
For a satellite in a circular orbit of radius $r$:
\[
E_k = \frac{GM_E m}{2r}, \qquad
E_p = -\frac{GM_E m}{r}, \qquad
\boxed{E = E_k + E_p = -\frac{GM_E m}{2r} = -E_k = \frac{E_p}{2}}.
\]
The total energy is negative: the satellite is \emph{bound}. To move it to a higher orbit (larger $r$, energy less negative), the engines must do positive work $\Delta E = \dfrac{GM_E m}{2}\left(\dfrac{1}{r_1} - \dfrac{1}{r_2}\right)$.
\end{propriete}

\textbf{Proof.} From the orbit condition, $mv^{2} = GM_E m/r$, so $E_k = \tfrac{1}{2}mv^{2} = \dfrac{GM_E m}{2r}$. Adding $E_p = -\dfrac{GM_E m}{r}$ gives $E = -\dfrac{GM_E m}{2r}$. $\square$

\begin{exoresolu}[A low-orbit observation satellite]
A satellite of mass $m = 800\ \mathrm{kg}$ circles the Earth at an altitude $h = 600\ \mathrm{km}$. Take $M_E = 5.98\times 10^{24}\ \mathrm{kg}$, $R_E = 6.37\times 10^{6}\ \mathrm{m}$. Find (a) its orbital speed, (b) its period, (c) its total mechanical energy, (d) the energy needed to raise it to an altitude of $1000\ \mathrm{km}$.
\tcblower
The orbit radius is $r = R_E + h = 6.37\times 10^{6} + 0.600\times 10^{6} = 6.97\times 10^{6}\ \mathrm{m}$.

\textbf{(a) Speed:}
\[
v = \sqrt{\frac{GM_E}{r}} = \sqrt{\frac{6.67\times 10^{-11}\times 5.98\times 10^{24}}{6.97\times 10^{6}}} \approx 7.57\times 10^{3}\ \mathrm{m\,s^{-1}} \approx 7.6\ \mathrm{km\,s^{-1}}.
\]

\textbf{(b) Period:}
\[
T = 2\pi\sqrt{\frac{r^{3}}{GM_E}} = 2\pi\sqrt{\frac{(6.97\times 10^{6})^{3}}{3.99\times 10^{14}}} \approx 5.80\times 10^{3}\ \mathrm{s} \approx 96.6\ \mathrm{min}.
\]

\textbf{(c) Total energy:}
\[
E = -\frac{GM_E m}{2r} = -\frac{3.99\times 10^{14}\times 800}{2\times 6.97\times 10^{6}} \approx -2.29\times 10^{10}\ \mathrm{J}.
\]
Negative, as expected for a bound satellite.

\textbf{(d) Energy to raise the orbit} to $r_2 = 7.37\times 10^{6}\ \mathrm{m}$:
\[
\Delta E = \frac{GM_E m}{2}\left(\frac{1}{r_1} - \frac{1}{r_2}\right)
= \frac{3.99\times 10^{14}\times 800}{2}\left(\frac{1}{6.97\times 10^{6}} - \frac{1}{7.37\times 10^{6}}\right)
\approx 1.24\times 10^{9}\ \mathrm{J}.
\]
About $1.2\ \mathrm{GJ}$ must be supplied by the engines — positive work, because the new orbit is less tightly bound.
\end{exoresolu}

\courssub{Geostationary Satellites}

\textbf{The idea.} A \cle{geostationary satellite} appears to hang motionless above one point of the equator. Television dishes in Yaoundé, Buea or Bamenda can therefore be bolted in a fixed direction and still receive CRTV or Canal+ relays day and night. For this to happen, three conditions must hold simultaneously:

\begin{enumerate}
\item the orbit is \emph{circular} and lies \emph{in the equatorial plane};
\item the period equals the Earth's rotation period, $T = 24\ \mathrm{h}$ (more precisely one sidereal day, $\approx 86\,164\ \mathrm{s}$);
\item the satellite revolves in the \emph{same sense} as the Earth's rotation (west to east).
\end{enumerate}

\begin{methode}[Finding the geostationary radius]
\begin{enumerate}
\item Write Kepler's third law for the orbit: $T^{2} = \dfrac{4\pi^{2} r^{3}}{GM_E}$.
\item Impose the geostationary condition $T = 24\ \mathrm{h} = 8.64\times 10^{4}\ \mathrm{s}$ (or $86\,164\ \mathrm{s}$ for the sidereal day).
\item Solve for $r$: $r = \left(\dfrac{GM_E T^{2}}{4\pi^{2}}\right)^{1/3}$.
\item Subtract the Earth's radius to get the altitude: $h = r - R_E$.
\end{enumerate}
\end{methode}

\begin{exemplebox}[Radius and altitude of the geostationary orbit]
Using $T = 8.64\times 10^{4}\ \mathrm{s}$:
\[
r = \left(\frac{3.99\times 10^{14}\times (8.64\times 10^{4})^{2}}{4\pi^{2}}\right)^{1/3}
\approx 4.22\times 10^{7}\ \mathrm{m} \approx 42\,200\ \mathrm{km},
\]
\[
h = r - R_E \approx 42\,200 - 6\,370 \approx 35\,800\ \mathrm{km}.
\]
(Using the sidereal day $T = 86\,164\ \mathrm{s}$ gives the accepted precise value $h \approx 35\,786\ \mathrm{km}$.) The orbital speed there is
\[
v = \sqrt{\frac{GM_E}{r}} \approx 3.07\times 10^{3}\ \mathrm{m\,s^{-1}} \approx 3.1\ \mathrm{km\,s^{-1}},
\]
much slower than a low-orbit satellite, exactly as $v = \sqrt{GM/r}$ predicts.
\end{exemplebox}

\begin{popfigure}
\begin{center}
\begin{tikzpicture}[scale=0.62]
\shade[inner color=popblue, outer color=popblueL] (0,0) circle (1.0);
\node[white, font=\footnotesize\bfseries] at (0,0) {Earth};
\draw[thick, popgreen] (0,0) circle (1.7);
\fill[popgreen] (1.7,0) circle (0.09);
\node[popgreen!60!black, font=\footnotesize, anchor=south west, align=center] at (1.6,0.1) {low orbit\\ $h \approx 600$ km};
\draw[thick, poporange] (0,0) circle (6.63);
\fill[poporange] (-6.63,0) circle (0.12);
\node[poporange!70!black, font=\footnotesize, anchor=south east, align=center] at (-6.5,0.15) {geostationary\\ $r \approx 42\,200$ km};
\draw[{Stealth[length=2mm]}-{Stealth[length=2mm]}, popmuted] (0,-1.35) -- (6.63,-1.35) node[midway, below, font=\footnotesize] {$r_{\text{geo}} \approx 6.6\,R_E$};
\end{tikzpicture}
\end{center}
\legende{To scale: a low Earth orbit ($r \approx 1.1\,R_E$) compared with the geostationary orbit ($r \approx 6.6\,R_E$).}
\end{popfigure}

\begin{savaistu}[Why above the equator? — enrichment beyond the strict syllabus]
The centre of any satellite's orbit must coincide with the centre of the Earth, because gravity pulls towards that centre. An orbit tilted relative to the equator is perfectly possible, but its ground track drifts north and south during the day: seen from the ground the satellite traces a figure of eight and is \emph{not} stationary. Only an orbit \emph{in the equatorial plane} keeps the satellite permanently above one fixed point. A consequence: geostationary satellites sit on the equatorial ring, so they see the polar regions only at a very grazing angle and cannot serve them properly. Polar observation and high-latitude communications use inclined or \emph{polar} low orbits instead — as do GPS-type constellations, whose medium-altitude satellites ($T \approx 12\ \mathrm{h}$) are not geostationary at all.
\end{savaistu}

\begin{attention}[Frequent examination traps]
\begin{itemize}
\item Confusing \emph{altitude} $h$ with \emph{orbit radius} $r = R_E + h$. All formulas ($v$, $T$, $E$) use $r$, never $h$.
\item Forgetting to convert the period to seconds before using Kepler's third law.
\item Writing $E_k = -E$ as $E_k = +\dfrac{GMm}{2r}$ but then claiming total energy is positive: for any bound orbit, $E < 0$.
\item Believing a geostationary satellite is ``weightless because far away'': at $r = 42\,200\ \mathrm{km}$, $g \approx 0.22\ \mathrm{m\,s^{-2}}$, small but \emph{not zero} — it is precisely this force that holds the orbit.
\item Forgetting that escape velocity is independent of launch direction (in the absence of atmosphere); only the speed matters.
\end{itemize}
\end{attention}

\begin{exercice}[Between two orbits]
A $500\ \mathrm{kg}$ satellite is transferred from a circular orbit of radius $r_1 = 7.0\times 10^{6}\ \mathrm{m}$ to the geostationary radius $r_2 = 4.22\times 10^{7}\ \mathrm{m}$. Using $GM_E = 3.99\times 10^{14}\ \mathrm{m^{3}\,s^{-2}}$:
\begin{enumerate}
\item Compute its speed on each orbit.
\item Compute its total mechanical energy on each orbit.
\item Deduce the work the engines must supply. Comment on its sign.
\end{enumerate}
\end{exercice}

\begin{corrige}[Exercise 1]
\textbf{1.} $v_1 = \sqrt{GM_E/r_1} = \sqrt{3.99\times 10^{14}/7.0\times 10^{6}} \approx 7.55\times 10^{3}\ \mathrm{m\,s^{-1}}$; $v_2 = \sqrt{3.99\times 10^{14}/4.22\times 10^{7}} \approx 3.07\times 10^{3}\ \mathrm{m\,s^{-1}}$. The higher orbit is slower.

\textbf{2.} $E_1 = -\dfrac{GM_E m}{2r_1} = -\dfrac{3.99\times 10^{14}\times 500}{1.4\times 10^{7}} \approx -1.43\times 10^{10}\ \mathrm{J}$; $E_2 = -\dfrac{3.99\times 10^{14}\times 500}{8.44\times 10^{7}} \approx -2.36\times 10^{9}\ \mathrm{J}$.

\textbf{3.} $W = \Delta E = E_2 - E_1 \approx -2.36\times 10^{9} - (-1.43\times 10^{10}) \approx +1.19\times 10^{10}\ \mathrm{J}$. The work is \emph{positive}: raising a satellite to a less bound (higher) orbit always costs energy, even though the satellite ends up moving \emph{slower} — the gain in potential energy exceeds the loss in kinetic energy.
\end{corrige}

\begin{aretenir}[Key points of the section]
\begin{itemize}
\item Gravity is conservative: $E = E_k + E_p$ is conserved, with $E_p = -\dfrac{GMm}{r}$ and zero at infinity. Negative $E_p$ means a \emph{bound} system.
\item Gravitational potential: $V = -\dfrac{GM}{r}$ (in $\mathrm{J\,kg^{-1}}$), with $E_p = mV$ and $W_{\text{ext}} = m\,\Delta V$.
\item Escape velocity: $v_e = \sqrt{\dfrac{2GM}{R}} = \sqrt{2g_0 R} \approx 11.2\ \mathrm{km\,s^{-1}}$ for the Earth.
\item Circular orbit: $v = \sqrt{\dfrac{GM}{r}}$, $T = 2\pi\sqrt{\dfrac{r^{3}}{GM}}$, $E = -\dfrac{GMm}{2r} = -E_k = \dfrac{E_p}{2}$.
\item Geostationary orbit: equatorial, $T = 24\ \mathrm{h}$ (sidereal), same sense as the Earth's rotation; $r \approx 42\,200\ \mathrm{km}$, $h \approx 35\,786\ \mathrm{km}$.
\end{itemize}
\end{aretenir}

\coursec{Motion in the Gravitational Field: Projectile Motion}

In the previous sections we treated gravity on a grand scale: potentials, escape speeds and satellites orbiting the whole Earth. There, the field was \emph{radial} and $g$ varied significantly with distance. We now zoom in to a completely different scale: a football kicked on a field in Buea, a stone thrown across a stream, a mango falling from a tree. Over distances of a few tens of metres, the curvature of the Earth is utterly negligible and the gravitational field can be treated as \cle{uniform}. This simplification turns the motion of any thrown object into one of the most elegant and examinable problems in all of mechanics: \cle{projectile motion}.

\courssub{The Uniform Field Approximation}

\begin{definition}[Uniform gravitational field]
A gravitational field is \cle{uniform} over a region if the field vector $\vec{g}$ has the \textbf{same magnitude and the same direction} at every point of that region. Near the Earth's surface, over distances small compared with the Earth's radius $R_E$, $\vec{g}$ is vertical (pointing towards the centre of the Earth) with magnitude $g \approx \SI{9.81}{m.s^{-2}}$.
\end{definition}

Why is this legitimate? At altitude $h$ above the surface, $g(h) = \dfrac{GM_E}{(R_E+h)^2}$. For $h = \SI{100}{m}$ and $R_E \approx \SI{6.37e6}{m}$, the relative change in $g$ is about $2h/R_E \approx 3\times 10^{-5}$ — far below the precision of any school measurement. So for any projectile problem at the GCE, we take $\vec{g}$ constant and vertical.

\begin{attention}[When the uniform model fails]
The uniform-field model is only valid for \emph{small-scale} motion near the surface. For satellites, intercontinental ballistic motion, or anything rising a significant fraction of $R_E$, you must return to the inverse-square law of Sections 1--3.
\end{attention}

\courssub{Setting Up the Equations: Independence of the Motions}

Consider a projectile launched from the origin with speed $v_0$ at angle $\theta$ above the horizontal. Choose axes: $Ox$ horizontal in the plane of motion, $Oy$ vertically upward. The only force acting (we neglect air resistance) is the weight $m\vec{g}$, directed along $-Oy$.

Newton's second law gives:
\[
m\vec{a} = m\vec{g} \quad\Longrightarrow\quad \vec{a} = \vec{g}
\]
so that, component by component:
\[
a_x = 0 \qquad\text{and}\qquad a_y = -g.
\]

\begin{propriete}[Independence of horizontal and vertical motions]
In a uniform gravitational field, the motion of a projectile is the \textbf{superposition of two independent motions}:
\begin{itemize}
\item a \cle{uniform motion} horizontally: $a_x = 0$, so $v_x = v_0\cos\theta$ is constant;
\item a \cle{uniformly accelerated motion} vertically: $a_y = -g$ (constant).
\end{itemize}
The horizontal motion does not influence the vertical motion, and vice versa. This is the key that unlocks every projectile problem.
\end{propriete}

Integrating with the initial conditions $v_x(0) = v_0\cos\theta$, $v_y(0) = v_0\sin\theta$, $x(0) = y(0) = 0$:

\begin{align*}
v_x &= v_0\cos\theta, & v_y &= v_0\sin\theta - gt,\\
x &= (v_0\cos\theta)\,t, & y &= (v_0\sin\theta)\,t - \tfrac{1}{2}gt^2.
\end{align*}

\begin{experience}[Seeing the independence with your own eyes]
Drop one coin vertically from the edge of a table and, at the same instant, flick a second coin horizontally from the same height. Listen: \textbf{both coins hit the floor at the same time}. The horizontal coin travels further, but its \emph{fall} is identical to that of the dropped coin. Horizontal and vertical motions truly are independent.
\end{experience}

\courssub{The Trajectory: a Parabola}

Eliminate $t$ between the parametric equations. From $x = (v_0\cos\theta)t$, we get $t = \dfrac{x}{v_0\cos\theta}$. Substituting into $y(t)$:

\begin{propriete}[Trajectory equation]
\[
\boxed{\;y = x\tan\theta - \frac{g\,x^2}{2v_0^{\,2}\cos^2\theta}\;}
\]
This is of the form $y = ax - bx^2$ with $b>0$: the trajectory is a \cle{parabola} opening \emph{downward}.
\end{propriete}

\begin{popfigure}
\begin{center}
\begin{tikzpicture}[scale=1.0, >=stealth]
% axes
\draw[->, popink, thick] (0,0) -- (8.6,0) node[below left] {$x$};
\draw[->, popink, thick] (0,0) -- (0,3.6) node[below left] {$y$};
% parabola y = x tan45 - g x^2/(2 v0^2 cos^2 45), choose v0^2/g = 8, so y = x - x^2/8
\draw[popblue, very thick, domain=0:8, samples=60] plot (\x, {\x - \x*\x/8});
% launch angle and velocity vector
\draw[->, poporange, very thick] (0,0) -- (1.6,1.6) node[midway, above left] {$\vec{v}_0$};
\draw[->, popgreen, thick] (0,0) -- (1.6,0) node[midway, below] {$v_0\cos\theta$};
\draw[->, popgreen, thick] (1.6,0) -- (1.6,1.6) node[midway, right] {$v_0\sin\theta$};
\draw[popdark] (0.9,0) arc (0:45:0.9) node[midway, right] {$\theta$};
% max height
\draw[dashed, popmuted] (0,2) -- (4,2);
\draw[<->, poppurple, thick] (-0.25,0) -- (-0.25,2) node[midway, left] {$H$};
\fill[popblue] (4,2) circle (2pt) node[above right, popink] {top: $v_y = 0$};
% range
\draw[<->, poppurple, thick] (0,-0.35) -- (8,-0.35) node[midway, below] {$R$};
\fill[popblue] (8,0) circle (2pt);
\fill[popblue] (0,0) circle (2pt);
\end{tikzpicture}
\end{center}
\legende{Parabolic trajectory of a projectile launched at speed $v_0$ and angle $\theta$. The initial velocity splits into a constant horizontal component $v_0\cos\theta$ and a vertical component $v_0\sin\theta$ that gravity steadily reduces.}
\end{popfigure}

\courssub{Time of Flight, Maximum Height and Range}

We now derive the three classic results for a launch \textbf{from ground level, landing at the same level}.

\textbf{Time of flight.} Landing means $y = 0$ with $t>0$:
\[
(v_0\sin\theta)t - \tfrac{1}{2}gt^2 = 0 \;\Longrightarrow\; t\left(v_0\sin\theta - \tfrac{1}{2}gt\right) = 0
\;\Longrightarrow\; T = \frac{2v_0\sin\theta}{g}.
\]

\textbf{Maximum height.} At the top, $v_y = 0$: $v_0\sin\theta - gt_{\text{top}} = 0$, so $t_{\text{top}} = \dfrac{v_0\sin\theta}{g} = \dfrac{T}{2}$ (the motion is symmetric in time). Then
\[
H = (v_0\sin\theta)t_{\text{top}} - \tfrac{1}{2}g t_{\text{top}}^2 = \frac{v_0^{\,2}\sin^2\theta}{2g}.
\]

\textbf{Range.} Substituting $T$ into $x = (v_0\cos\theta)t$:
\[
R = (v_0\cos\theta)\cdot\frac{2v_0\sin\theta}{g} = \frac{v_0^{\,2}\cdot 2\sin\theta\cos\theta}{g} = \frac{v_0^{\,2}\sin 2\theta}{g}.
\]

\begin{aretenir}[Ground-level launch: the three key results]
For a projectile launched at speed $v_0$ and angle $\theta$, landing at the \textbf{same height}:
\[
T = \frac{2v_0\sin\theta}{g}, \qquad
H = \frac{v_0^{\,2}\sin^2\theta}{2g}, \qquad
R = \frac{v_0^{\,2}\sin 2\theta}{g}.
\]
\end{aretenir}

\textbf{Maximum range and complementary angles.} Since $\sin 2\theta \le 1$ with maximum at $2\theta = 90\degree$:
\[
\theta = 45\degree \quad\Longrightarrow\quad R_{\max} = \frac{v_0^{\,2}}{g}.
\]
Because $\sin 2\theta = \sin(180\degree - 2\theta)$, two complementary angles $\theta$ and $90\degree - \theta$ (e.g. $30\degree$ and $60\degree$) give the \textbf{same range} — but different times of flight and different maximum heights: the high lob ($60\degree$) stays in the air twice as long as the flat throw ($30\degree$) for the same landing point.

\begin{popfigure}
\begin{center}
\begin{tikzpicture}
\begin{axis}[
    width=0.85\linewidth, height=6.2cm,
    xlabel={$x$ (m)}, ylabel={$y$ (m)},
    xmin=0, xmax=105, ymin=0, ymax=55,
    grid=both, grid style={popline, very thin},
    legend style={at={(0.98,0.97)}, anchor=north east, font=\small},
]
% v0 = 31.3 m/s, g = 9.81: v0^2/g = 100 m
% y = x tanθ - x^2/(200 cos^2θ)
\addplot[popblue, very thick, domain=0:86.6, samples=80] {x*0.5774 - x*x/(200*0.75)};
\addlegendentry{$\theta = 30\degree$}
\addplot[poporange, very thick, domain=0:100, samples=80] {x - x*x/100};
\addlegendentry{$\theta = 45\degree$}
\addplot[popgreen, very thick, domain=0:86.6, samples=80] {x*1.7321 - x*x/50};
\addlegendentry{$\theta = 60\degree$}
\end{axis}
\end{tikzpicture}
\end{center}
\legende{Trajectories for the same launch speed $v_0$ (here $v_0^{\,2}/g = \SI{100}{m}$) at $\theta = 30\degree, 45\degree, 60\degree$. The $45\degree$ launch gives the greatest range; the $30\degree$ and $60\degree$ launches share the same range but very different heights and flight times.}
\end{popfigure}

\begin{attention}[The velocity at the top is NOT zero]
A very common error: "at the highest point the velocity is zero." \textbf{False.} Only the \emph{vertical} component vanishes: $v_y = 0$. The horizontal component $v_x = v_0\cos\theta$ is unchanged throughout the flight, so at the top the speed is $v_0\cos\theta$, directed horizontally. If the velocity were zero, the projectile would simply fall straight down from there!
\end{attention}

\begin{attention}[The range formula has a restricted domain of validity]
$R = \dfrac{v_0^{\,2}\sin 2\theta}{g}$ applies \textbf{only} when launch and landing are at the same height. If a ball is kicked from the top of a hill, or lands on a roof, you must go back to the parametric equations and solve $y(t) = h_{\text{landing}}$ for the time, then compute $x(T)$. Blindly applying the formula is a classic source of lost marks.
\end{attention}

\courssub{Horizontal Launch from a Height}

A stone is thrown \emph{horizontally} with speed $v_0$ from the top of a cliff of height $h$ (or a ball rolls off a table). This is the special case $\theta = 0\degree$, with launch point at height $h$.

Take the origin at the \textbf{base of the cliff}, $Oy$ upward. Initial conditions: $x(0) = 0$, $y(0) = h$, $v_x(0) = v_0$, $v_y(0) = 0$. Hence:
\[
x = v_0 t, \qquad y = h - \tfrac{1}{2}gt^2.
\]

\textbf{Time of fall:} set $y = 0$:
\[
t_f = \sqrt{\frac{2h}{g}} \quad\text{— exactly the free-fall time from height } h.
\]
The horizontal speed has \emph{no effect} on the time of fall — the independence principle again.

\textbf{Horizontal distance (range from the cliff base):}
\[
d = v_0\, t_f = v_0\sqrt{\frac{2h}{g}}.
\]

\textbf{Impact velocity:} at impact, $v_x = v_0$ and $v_y = -gt_f = -\sqrt{2gh}$, so the speed is $v = \sqrt{v_0^{\,2} + 2gh}$ — which energy conservation ($\tfrac12 m v^2 = \tfrac12 m v_0^2 + mgh$) confirms instantly.

\begin{popfigure}
\begin{center}
\begin{tikzpicture}[scale=0.9, >=stealth]
% cliff
\fill[popblueL] (0,0) rectangle (2.2,3.2);
\draw[popink, thick] (0,0) -- (2.2,0) (2.2,0) -- (2.2,3.2) (0,3.2) -- (2.2,3.2);
\draw[popink] (1.1,1.6) node {cliff};
\draw[<->, poppurple, thick] (0.3,0) -- (0.3,3.2) node[midway, left] {$h$};
% ground
\draw[popink, thick] (2.2,0) -- (9.2,0);
% trajectory: y = h - g x^2/(2 v0^2), h=3.2, fall at x=5.5
\draw[poporange, very thick, domain=0:5.5, samples=50] plot ({2.2+\x}, {3.2 - 3.2*\x*\x/30.25});
% initial velocity
\draw[->, popgreen, very thick] (2.2,3.2) -- (4.0,3.2) node[midway, above] {$\vec{v}_0$ (horizontal)};
\fill[popdark] (2.2,3.2) circle (2.5pt);
% impact
\fill[popblue] (7.7,0) circle (2.5pt) node[above right, popink] {impact};
\draw[<->, poppurple, thick] (2.2,-0.45) -- (7.7,-0.45) node[midway, below] {$d = v_0\sqrt{2h/g}$};
% vertical drop arrow
\draw[->, popmuted, dashed] (5.2,3.2) -- (5.2,1.4) node[midway, right] {free fall};
\end{tikzpicture}
\end{center}
\legende{Horizontal launch from a cliff of height $h$. The vertical motion is pure free fall from rest; the horizontal motion is uniform at speed $v_0$. The time of fall depends only on $h$.}
\end{popfigure}

\begin{exoresolu}[Ball rolling off a laboratory table]
A ball rolls horizontally off a table of height $h = \SI{1.25}{m}$ with speed $v_0 = \SI{2.0}{m.s^{-1}}$. Take $g = \SI{10}{m.s^{-2}}$. Find (a) the time of fall, (b) the horizontal distance from the foot of the table to the impact point, (c) the velocity at impact.
\tcblower
\textbf{(a) Time of fall.} Vertical motion: $y = h - \tfrac12 g t^2$. Setting $y = 0$:
\[
t_f = \sqrt{\frac{2h}{g}} = \sqrt{\frac{2\times 1.25}{10}} = \sqrt{0.25} = \SI{0.50}{s}.
\]
\textbf{(b) Horizontal distance.} Uniform horizontal motion:
\[
d = v_0 t_f = 2.0 \times 0.50 = \SI{1.0}{m}.
\]
\textbf{(c) Impact velocity.} Components at impact:
\[
v_x = v_0 = \SI{2.0}{m.s^{-1}}, \qquad v_y = -gt_f = -10\times 0.50 = \SI{-5.0}{m.s^{-1}}.
\]
\[
v = \sqrt{v_x^2 + v_y^2} = \sqrt{4.0 + 25} = \SI{5.4}{m.s^{-1}}, \qquad
\tan\alpha = \frac{|v_y|}{v_x} = 2.5 \;\Rightarrow\; \alpha \approx 68\degree \text{ below the horizontal.}
\]
Check with energy: $v = \sqrt{v_0^2 + 2gh} = \sqrt{4 + 25} = \SI{5.4}{m.s^{-1}}$. $\checkmark$
\end{exoresolu}

\courssub{Projectiles as Free Fall: the Monkey and the Hunter}

Here is a famous piece of reasoning that reveals the deep meaning of projectile motion. A hunter aims \emph{directly} at a monkey hanging from a branch. At the instant the dart leaves the blowpipe, the monkey lets go and falls. Does the dart hit the monkey?

\textbf{Yes — always} (provided the dart reaches the monkey's vertical line before the ground intervenes). The argument: imagine gravity were switched off. The dart would fly straight along the line of aim and hit the monkey at the aimed point $P$. Now switch gravity on. In time $t$, \emph{both} the dart and the monkey fall the \emph{same} extra distance $\tfrac12 g t^2$ below where they would have been. So at time $t$ the dart is exactly $\tfrac12 g t^2$ below $P$, and the monkey is exactly $\tfrac12 g t^2$ below its starting point — which is $P$. They meet.

\begin{aretenir}[The deep idea]
A projectile is in \cle{free fall} at every instant: relative to the straight-line path it would follow without gravity, it "falls" by $\tfrac12 g t^2$, whatever its launch speed or angle. Gravity accelerates \emph{everything} identically — dart, monkey, football, satellite.
\end{aretenir}

\courssub{From Projectiles to Satellites: Falling Around the Earth}

Newton himself proposed a beautiful thought experiment. Fire a projectile horizontally from the top of a very high mountain, above the atmosphere. At modest speed, it follows a parabola (locally) and lands at $A$. Fire it faster: it travels further before landing, at $B$, then $C$. But the Earth is round — its surface curves \emph{away} from the projectile. There exists a special speed for which the projectile's fall toward the Earth exactly matches the curving away of the Earth's surface: the projectile \textbf{never lands}. It is perpetually falling toward the Earth while never getting closer to the ground: it is in \cle{orbit}.

That special speed is precisely the orbital speed of Section 3: $v = \sqrt{\dfrac{GM_E}{R_E}} \approx \SI{7.9}{km.s^{-1}}$ for a low orbit. A satellite is nothing more than a projectile launched fast enough to \emph{fall around the Earth} forever. This closes the conceptual loop of the whole chapter: the same inverse-square law that holds the Moon also bends the path of a thrown stone.

\begin{savaistu}[Why astronauts float]
Astronauts in the International Space Station are \emph{not} beyond Earth's gravity — $g$ up there is still about $90\%$ of its surface value. They float because the station and everything in it are in continuous free fall together, exactly like the monkey and the dart. "Weightlessness" is really "fallingness."
\end{savaistu}

\courssub{Solving GCE-Style Projectile Problems}

\begin{methode}[A four-step strategy for any projectile problem]
\begin{enumerate}
\item \textbf{Choose axes and state the sign convention.} Usually $Ox$ horizontal in the direction of motion, $Oy$ vertically \emph{upward}, origin at the launch point (or at ground level for cliff problems). Then $a_x = 0$, $a_y = -g$.
\item \textbf{Resolve the initial velocity.} $v_{0x} = v_0\cos\theta$, $v_{0y} = v_0\sin\theta$. For a horizontal launch, $v_{0y} = 0$.
\item \textbf{Write the four equations.}
\[
v_x = v_{0x}, \quad v_y = v_{0y} - gt, \quad x = v_{0x}t, \quad y = v_{0y}t - \tfrac12 gt^2.
\]
\item \textbf{Translate the question into a condition and solve.} "Highest point" $\Rightarrow v_y = 0$. "Lands" $\Rightarrow y = $ landing height. "Hits a wall" $\Rightarrow x = $ wall distance. Then \textbf{check units and sign} of every answer.
\end{enumerate}
\end{methode}

\begin{exoresolu}[GCE-style: the football free kick]
A footballer kicks a ball from ground level with speed $v_0 = \SI{20}{m.s^{-1}}$ at $\theta = 30\degree$ above the horizontal. Take $g = \SI{10}{m.s^{-2}}$. Determine (a) the time of flight, (b) the maximum height, (c) the range, (d) the speed at the highest point.
\tcblower
Components: $v_{0x} = 20\cos 30\degree = 10\sqrt{3} \approx \SI{17.3}{m.s^{-1}}$, $v_{0y} = 20\sin 30\degree = \SI{10}{m.s^{-1}}$.

\textbf{(a)} $T = \dfrac{2v_0\sin\theta}{g} = \dfrac{2\times 10}{10} = \SI{2.0}{s}$.

\textbf{(b)} $H = \dfrac{v_0^{\,2}\sin^2\theta}{2g} = \dfrac{400\times 0.25}{20} = \SI{5.0}{m}$.

\textbf{(c)} $R = \dfrac{v_0^{\,2}\sin 2\theta}{g} = \dfrac{400\times\sin 60\degree}{10} = 40\times 0.866 = \SI{35}{m}$ (2 s.f.).

\textbf{(d)} At the top, $v_y = 0$ but $v_x$ persists: $v = v_{0x} = \SI{17.3}{m.s^{-1}}$ (horizontal). Note the answer is \emph{not} zero.
\end{exoresolu}

\begin{exercice}[Practice]
A stone is thrown horizontally with speed $\SI{15}{m.s^{-1}}$ from the top of a cliff $\SI{45}{m}$ high above the sea. Take $g = \SI{10}{m.s^{-2}}$ and neglect air resistance.
\begin{enumerate}
\item Calculate the time the stone takes to reach the sea.
\item Calculate the horizontal distance from the foot of the cliff to the point of impact.
\item Calculate the speed of the stone just before impact.
\item A second stone is simply \emph{dropped} from the top of the cliff at the same instant. Which stone hits the sea first? Justify your answer.
\end{enumerate}
\end{exercice}

\begin{corrige}[Exercise 1]
\textbf{1.} Vertical motion is free fall from rest: $t_f = \sqrt{2h/g} = \sqrt{2\times 45/10} = \sqrt{9} = \SI{3.0}{s}$.

\textbf{2.} $d = v_0 t_f = 15\times 3.0 = \SI{45}{m}$.

\textbf{3.} $v_x = \SI{15}{m.s^{-1}}$, $v_y = -gt_f = \SI{-30}{m.s^{-1}}$, so $v = \sqrt{15^2 + 30^2} = \sqrt{1125} \approx \SI{33.5}{m.s^{-1}}$.

\textbf{4.} Both hit the sea at the \textbf{same time}. The vertical motion of the thrown stone is identical to that of the dropped stone (both start with zero vertical velocity and fall $45$ m under gravity). The horizontal speed affects only \emph{where} it lands, not \emph{when}.
\end{corrige}

\begin{aretenir}[Chapter summary of Section 4]
\begin{itemize}
\item Near the Earth's surface, $\vec g$ is uniform: constant magnitude $\SI{9.81}{m.s^{-2}}$, constant (vertical) direction.
\item Projectile motion = uniform horizontal motion $+$ uniformly accelerated vertical motion; the two are \textbf{independent}.
\item Trajectory: $y = x\tan\theta - \dfrac{gx^2}{2v_0^{\,2}\cos^2\theta}$ — a parabola.
\item Same-level launch: $T = \dfrac{2v_0\sin\theta}{g}$, $H = \dfrac{v_0^{\,2}\sin^2\theta}{2g}$, $R = \dfrac{v_0^{\,2}\sin 2\theta}{g}$; maximum range at $\theta = 45\degree$; complementary angles give equal ranges.
\item Horizontal launch from height $h$: $t_f = \sqrt{2h/g}$ (independent of $v_0$), $d = v_0\sqrt{2h/g}$.
\item At the top of the path, only $v_y$ is zero; the speed is $v_0\cos\theta$.
\item A satellite is a projectile falling around the Earth — the conceptual bridge between this section and orbital motion.
\end{itemize}
\end{aretenir}

\newpage

\coursec{Integration activity}

\coursec{Integration Activity: Designing a Cameroonian Communications Satellite}

\courssub{Aims of the Activity}

This integration activity closes the chapter \emph{Gravitational Fields} by asking you to combine, in one authentic engineering-style project, every major result you have studied:

\begin{itemize}
\item \cle{Newton's law of gravitation} $F = \dfrac{GMm}{r^{2}}$ and the gravitational field strength $g(r) = \dfrac{GM}{r^{2}}$;
\item \cle{Kepler's third law} $T^{2} = \dfrac{4\pi^{2}}{GM}r^{3}$ applied to circular orbits;
\item \cle{Orbital speed} $v = \sqrt{\dfrac{GM}{r}}$ for a circular orbit;
\item \cle{Gravitational potential energy} $E_{p} = -\dfrac{GMm}{r}$ and energy balances between the Earth's surface and an orbit;
\item \cle{Escape velocity} $v_{\mathrm{esc}} = \sqrt{\dfrac{2GM}{R_{E}}}$ and its physical meaning;
\item scientific communication: presenting and defending a numerical design in front of an audience.
\end{itemize}

Beyond the calculations, the activity trains you to work like real engineers: choosing a model, stating its assumptions, computing with coherent units and sensible significant figures, comparing with accepted values, and discussing the limits of the model.

\courssub{Situation}

\textbf{Designing a Cameroonian Communications Satellite}\par

The (fictional) \emph{National Agency for Space Applications of Cameroon} wishes to provide permanent television broadcasting, telephone relay and internet coverage to the whole national territory, from the forests of the South Region to the savannas of the Far North. The agency's board decides that the most economical solution is a single \cle{geostationary} satellite: a satellite placed in a circular orbit in the plane of the equator, turning in the same sense as the Earth, with a period exactly equal to the Earth's rotation period. Seen from the ground in Yaoundé or Douala, such a satellite appears motionless in the sky, so that fixed parabolic antennas can point at it permanently.

Your class is divided into engineering teams. Each team receives the following mission order from the agency director:

\begin{quote}
\emph{``Determine the radius and altitude of the orbit our satellite must occupy, the speed it must be given there, the strength of the Earth's gravitational field at that altitude, the minimum energy per kilogram needed to raise it from the launch pad to its orbit, and the speed a rocket would need to escape the Earth completely. Present a one-page justified report.''}
\end{quote}

\textbf{Data provided to all teams}

\begin{center}
\begin{tabularx}{\linewidth}{|l|X|l|}
\hline
\rowcolor{popblueL} \textbf{Quantity} & \textbf{Symbol} & \textbf{Value} \\
\hline
Gravitational constant & $G$ & $\SI{6.67e-11}{N.m^{2}.kg^{-2}}$ \\
\hline
Mass of the Earth & $M_{E}$ & $\SI{5.97e24}{kg}$ \\
\hline
Mean radius of the Earth & $R_{E}$ & $\SI{6.37e6}{m}$ \\
\hline
Sidereal rotation period of the Earth & $T$ & $\SI{8.616e4}{s}$ (23 h 56 min) \\
\hline
Mass of the satellite & $m$ & $\SI{350}{kg}$ \\
\hline
\end{tabularx}
\end{center}

\textbf{Modelling assumptions}: the Earth is spherical and non-rotating for the field calculations; the orbit is circular and lies in the equatorial plane; only the Earth's gravity acts on the satellite (the Moon, the Sun and atmospheric drag are neglected); the launch pad is at sea level on the equator.

\begin{popfigure}
\begin{tikzpicture}[scale=0.55]
\def\RE{3.0} % Earth radius in cm
\def\r{5.2}  % Orbit radius in cm
% Earth
\fill[popblueL] (0,0) circle (\RE);
\draw[popblue, thick] (0,0) circle (\RE);
% Equatorial plane
\draw[popgreen, thick, dashed] (-\r-0.5,0) -- (\r+0.5,0);
\node[popgreen, above left, font=\small] at (-\r-0.5,0) {Equatorial plane};
% Orbit
\draw[poporange, thick, dashed] (0,0) circle (\r);
% Radius r
\draw[popdark, thick, <->] (0,0) -- (\r,0) node[midway, above, font=\small] {$r$};
% Radius RE
\draw[popdark, thick, <->] (0,0) -- (\RE,0) node[midway, below, font=\small] {$R_E$};
% Altitude h
\draw[popdark, thick, <->] (\RE, -0.4) -- (\r, -0.4) node[midway, below, font=\small] {$h$};
% Satellite
\fill[poporange] (\r,0) circle (0.12) node[right=2pt, font=\small] {Satellite ($m$)};
% Earth rotation arrow
\draw[->, poppink, thick] (1.2*\RE, 1.2*\RE) arc (45:135:1.7*\RE);
\node[poppink, font=\small] at (0, 2.2*\RE) {Earth rotation};
% Centre label
\node[popdark, font=\small] at (0,0) {Earth centre};
\end{tikzpicture}
\legende{Geometry of a geostationary orbit: circular, equatorial, radius $r = R_E + h$.}
\end{popfigure}

\courssub{Tasks}

\begin{exercice}[Mission Order]
\begin{enumerate}
\item \textbf{Orbit radius.} Starting from Newton's second law applied to the satellite in uniform circular motion (centripetal force supplied by gravity), derive Kepler's third law $T^{2} = \dfrac{4\pi^{2}}{GM_{E}}r^{3}$. Hence compute the orbit radius $r$ and the altitude $h = r - R_{E}$ of the geostationary orbit. Express $h$ in kilometres.

\item \textbf{Orbital speed and field strength.} Deduce the orbital speed $v$ of the satellite, in $\mathrm{m\,s^{-1}}$ then in $\mathrm{km\,h^{-1}}$. Compute the gravitational field strength $g(r)$ at the orbit and compare it with $g_{0} = \SI{9.81}{m.s^{-2}}$ at the surface. Is the satellite ``weightless''? Comment.

\item \textbf{Escape velocity.} Using an energy argument, derive the expression $v_{\mathrm{esc}} = \sqrt{\dfrac{2GM_{E}}{R_{E}}}$ and compute its value from the Earth's surface. Explain, in a few sentences, why the launcher must give the satellite a speed \emph{between} the orbital speed and the escape velocity, and why a speed equal to or greater than $v_{\mathrm{esc}}$ would be a failure for this mission.

\item \textbf{Energy cost.} Write the gravitational potential energy of the satellite on the launch pad ($r = R_{E}$) and on its orbit ($r = R_{E} + h$). Compute the minimum energy per kilogram, $\dfrac{\Delta E_{p}}{m}$, required to raise the satellite to its orbit (neglecting the kinetic energy of the orbital motion in a first approach). Then add the kinetic energy per kilogram $\dfrac{v^{2}}{2}$ needed on orbit and give the total energy per kilogram. Why is the real energy cost of a launch much higher?

\item \textbf{Comparison and report.} Compare your results with the accepted values (altitude $\approx \SI{35786}{km}$, orbital speed $\approx \SI{3.1}{km.s^{-1}}$). Identify possible sources of discrepancy. Write a one-page report (text, key formulas, one diagram) addressed to the agency director, justifying your design.
\end{enumerate}
\end{exercice}

\courssub{Model Solution}

\begin{exoresolu}[Complete Solution]

\textbf{1. Orbit radius from Kepler's third law.}\par
For a satellite of mass $m$ in a circular orbit of radius $r$, the gravitational force provides the necessary centripetal force:
\[
\frac{GM_{E}m}{r^{2}} = \frac{mv^{2}}{r}.
\]
The orbital speed is $v = \dfrac{2\pi r}{T}$, where $T$ is the orbital period. Substituting $v$:
\[
\frac{GM_{E}}{r^{2}} = \frac{4\pi^{2}r}{T^{2}}
\quad\Longrightarrow\quad
T^{2} = \frac{4\pi^{2}}{GM_{E}}\,r^{3},
\]
which is Kepler's third law for a circular orbit. For a geostationary satellite, $T$ must equal the Earth's sidereal rotation period. Solving for $r$:
\[
r = \left(\frac{GM_{E}T^{2}}{4\pi^{2}}\right)^{1/3}.
\]
Substitute the data ($G = \SI{6.67e-11}{N.m^{2}.kg^{-2}}$, $M_E = \SI{5.97e24}{kg}$, $T = \SI{8.616e4}{s}$):
\[
GM_{E} = \SI{3.982e14}{m^{3}.s^{-2}}, \qquad T^{2} = \SI{7.423e9}{s^{2}}.
\]
\[
r = \left(\frac{\SI{3.982e14}{m^{3}.s^{-2}} \times \SI{7.423e9}{s^{2}}}{4\pi^{2}}\right)^{1/3}
= \left(\frac{\SI{2.956e24}{m^{3}}}{\SI{39.48}{}}\right)^{1/3}
= \left(\SI{7.487e22}{m^{3}}\right)^{1/3}
\approx \SI{4.216e7}{m}.
\]
Hence the altitude:
\[
h = r - R_{E} = \SI{4.216e7}{m} - \SI{6.37e6}{m} = \SI{3.579e7}{m} \approx \SI{35790}{km}.
\]

\begin{aretenir}[Result 1]
Orbit radius $r \approx \SI{4.22e7}{m}$; Altitude $h \approx \SI{35790}{km}$ (3 significant figures).
\end{aretenir}

\textbf{2. Orbital speed and field strength.}\par
Orbital speed:
\[
v = \sqrt{\frac{GM_{E}}{r}} = \sqrt{\frac{\SI{3.982e14}{m^{3}.s^{-2}}}{\SI{4.216e7}{m}}}
= \sqrt{\SI{9.445e6}{m^{2}.s^{-2}}}
\approx \SI{3073}{m.s^{-1}} \approx \SI{1.11e4}{km.h^{-1}}.
\]
Gravitational field strength at the orbit:
\[
g(r) = \frac{GM_{E}}{r^{2}} = \frac{\SI{3.982e14}{m^{3}.s^{-2}}}{(\SI{4.216e7}{m})^{2}}
= \frac{\SI{3.982e14}{m^{3}.s^{-2}}}{\SI{1.777e15}{m^{2}}}
\approx \SI{0.224}{m.s^{-2}}.
\]
Comparison: $\dfrac{g(r)}{g_0} = \dfrac{0.224}{9.81} \approx 0.0228 \approx 2.3\%$.

\begin{attention}[Weightlessness?]
The satellite is \textbf{not} weightless. The gravitational field at that altitude is $2.3\%$ of the surface value; gravity is exactly what provides the centripetal acceleration keeping the satellite in orbit. The \emph{sensation} of weightlessness inside the satellite arises because the satellite and everything in it are in \cle{free fall} (they accelerate towards Earth at exactly $g(r)$).
\end{attention}

\textbf{3. Escape velocity.}\par
To escape to infinity with zero residual kinetic energy, the total mechanical energy at the surface must be at least zero:
\[
E_{\mathrm{total}} = \frac{1}{2}mv_{\mathrm{esc}}^{2} - \frac{GM_{E}m}{R_{E}} \ge 0.
\]
The minimum escape velocity corresponds to $E_{\mathrm{total}} = 0$:
\[
\frac{1}{2}mv_{\mathrm{esc}}^{2} = \frac{GM_{E}m}{R_{E}}
\quad\Longrightarrow\quad
v_{\mathrm{esc}} = \sqrt{\frac{2GM_{E}}{R_{E}}}.
\]
Numerically:
\[
v_{\mathrm{esc}} = \sqrt{\frac{2 \times \SI{3.982e14}{m^{3}.s^{-2}}}{\SI{6.37e6}{m}}}
= \sqrt{\SI{1.250e8}{m^{2}.s^{-2}}}
\approx \SI{11180}{m.s^{-1}} \approx \SI{11.2}{km.s^{-1}}.
\]

\begin{aretenir}[Result 3]
$v_{\mathrm{esc}} \approx \SI{11.2}{km.s^{-1}}$, while the orbital speed is $v \approx \SI{3.07}{km.s^{-1}}$.
\end{aretenir}

The launcher must impart a speed \emph{greater} than the low-orbit circular speed to reach the geostationary altitude, but \emph{strictly less} than $v_{\mathrm{esc}}$. If $v \ge v_{\mathrm{esc}}$, the satellite would have enough energy to leave the Earth's gravitational influence entirely and would never enter the desired closed orbit — the mission would fail. The useful injection window lies between the transfer-orbit perigee speed and the escape speed.

\textbf{4. Energy cost per kilogram.}\par
Gravitational potential energy per unit mass:
\[
\frac{E_{p}}{m} = -\frac{GM_{E}}{r}.
\]
At the surface ($r = R_E$):
\[
\frac{E_{p,\mathrm{surf}}}{m} = -\frac{\SI{3.982e14}{m^{3}.s^{-2}}}{\SI{6.37e6}{m}}
= -\SI{6.251e7}{J.kg^{-1}}.
\]
On the geostationary orbit ($r = \SI{4.216e7}{m}$):
\[
\frac{E_{p,\mathrm{orb}}}{m} = -\frac{\SI{3.982e14}{m^{3}.s^{-2}}}{\SI{4.216e7}{m}}
= -\SI{9.445e6}{J.kg^{-1}}.
\]
Minimum potential energy gain per kilogram (work against gravity):
\[
\frac{\Delta E_{p}}{m} = \frac{E_{p,\mathrm{orb}} - E_{p,\mathrm{surf}}}{m}
= GM_{E}\left(\frac{1}{R_{E}} - \frac{1}{r}\right)
= \SI{6.251e7}{J.kg^{-1}} - \SI{9.445e6}{J.kg^{-1}}
\approx \SI{5.307e7}{J.kg^{-1}}.
\]
Kinetic energy per kilogram required on the circular orbit:
\[
\frac{E_{k,\mathrm{orb}}}{m} = \frac{v^{2}}{2} = \frac{(\SI{3073}{m.s^{-1}})^{2}}{2}
\approx \SI{4.72e6}{J.kg^{-1}}.
\]
Total minimum mechanical energy per kilogram (ideal, no losses):
\[
\frac{E_{\mathrm{total}}}{m} = \frac{\Delta E_{p}}{m} + \frac{E_{k,\mathrm{orb}}}{m}
\approx \SI{5.307e7}{J.kg^{-1}} + \SI{4.72e6}{J.kg^{-1}}
= \SI{5.78e7}{J.kg^{-1}}.
\]
For the $350\ \mathrm{kg}$ satellite, this represents $\approx \SI{2.02e10}{J}$.

\begin{attention}[Real launch cost]
The real energy cost is \emph{far higher} because:
\begin{itemize}
\item The rocket must lift its own structure, engines, tanks and the vast majority of its mass as propellant (the \cle{tyranny of the rocket equation}).
\item Atmospheric drag during the first $\sim 100\ \mathrm{km}$ dissipates energy.
\item No engine is perfectly efficient (thermal losses, incomplete combustion).
\item Gravity losses: thrust spent fighting gravity while not yet at orbital speed.
\item Steering and plane-change manoeuvres.
\end{itemize}
Typical launchers require $\sim \SI{30}{MJ.kg^{-1}}$ to $\SI{50}{MJ.kg^{-1}}$ of payload to GEO, i.e., $50$--$100$ times the theoretical minimum.
\end{attention}

\textbf{5. Comparison, discussion and report outline.}\par
Our computed values:
\begin{itemize}
\item Altitude $h = \SI{35790}{km}$ vs. accepted $\SI{35786}{km}$ (difference $< 0.02\%$).
\item Orbital speed $v = \SI{3.07}{km.s^{-1}}$ vs. accepted $\SI{3.07}{km.s^{-1}}$ (excellent agreement).
\end{itemize}
Sources of the tiny residual discrepancies:
\begin{itemize}
\item Rounded values of $G$, $M_E$, $R_E$ (3 significant figures only).
\item Use of the \emph{sidereal} day ($86164\ \mathrm{s}$) rather than the solar day ($86400\ \mathrm{s}$) — correct for geostationary orbit.
\item Spherical Earth model (Earth's equatorial radius is $\SI{6378}{km}$, polar radius $\SI{6357}{km}$; $J_2$ oblateness perturbs the orbit slightly).
\item Neglect of lunar/solar perturbations and solar radiation pressure.
\end{itemize}

\begin{methode}[One-page report structure for the Agency Director]
\begin{enumerate}
\item \textbf{Title:} \emph{Geostationary Orbit Parameters for Cameroon Communications Satellite}.
\item \textbf{Summary:} One sentence giving $h$, $v$, $g(r)$, $\Delta E/m$, $v_{\mathrm{esc}}$.
\item \textbf{Assumptions:} Spherical Earth, equatorial circular orbit, only Earth gravity.
\item \textbf{Key formulas:} $T^2 = 4\pi^2 r^3/(GM_E)$, $v = \sqrt{GM_E/r}$, $g = GM_E/r^2$, $v_{\mathrm{esc}} = \sqrt{2GM_E/R_E}$, $\Delta E/m = GM_E(1/R_E - 1/r) + v^2/2$.
\item \textbf{Results table:} Symbol, Formula, Value, Unit.
\item \textbf{Diagram:} Earth, orbit, $R_E$, $h$, $r$, satellite, equatorial plane (as in the figure above).
\item \textbf{Conclusion:} Orbit validated; energy budget estimated; recommendation for launcher selection (e.g., Ariane 5/6, Falcon 9, Long March 3B capabilities).
\end{enumerate}
\end{methode}

\end{exoresolu}

\begin{attention}[Traps to Avoid in This Activity]
\begin{itemize}
\item Do not confuse \cle{radius} $r$ (from the Earth's centre) with \cle{altitude} $h$ (from the surface): every formula of this chapter uses $r$.
\item Convert the period to \emph{seconds} before substituting; mixing hours and metres is the most common source of absurd results.
\item In Kepler's third law, $T$ is the \emph{sidereal} day ($\SI{8.616e4}{s}$), not $24\ \mathrm{h}$: a geostationary satellite must turn with the Earth relative to the stars.
\item Potential energies are \emph{negative}; the energy to \emph{supply} is the positive difference $\Delta E_{p} = E_{p,\mathrm{orb}} - E_{p,\mathrm{surf}}$.
\item Escape velocity is derived from \emph{energy conservation}, not from force balance.
\end{itemize}
\end{attention}

\begin{savaistu}[Cameroon and Space Applications]
Cameroon does not yet have its own launch capability, but it is an active user of space services. The country hosts ground stations for satellite telemetry (e.g., in Kribi) and benefits from geostationary satellites like \emph{Intelsat}, \emph{Eutelsat}, \emph{RASCOM-QAF1} (the first pan-African satellite, launched 2007) and \emph{Belintersat-1} for broadcasting, VSAT internet, and telephony. The \emph{African Space Agency} (headquartered in Cairo, Egypt, since 2023) coordinates continental space policy. A future Cameroonian geostationary satellite would likely be procured from a manufacturer (Airbus, Thales Alenia Space, CAST) and launched on a commercial rocket (Arianespace, SpaceX, CNSA), exactly as modelled in this activity.
\end{savaistu}

\begin{aretenir}[Link to Projectile Motion (Lower Sixth Mechanics)]
The trajectory of a projectile near Earth's surface is a parabola (constant $g$ approximation). The trajectory of a satellite is a conic section (ellipse, parabola, hyperbola) governed by the inverse-square law. A geostationary orbit is a \emph{circular ellipse} ($e=0$). An escape trajectory is a \emph{parabola} ($e=1$, total energy $=0$). A hyperbolic trajectory ($e>1$, total energy $>0$) corresponds to interplanetary transfer. The unification of projectile and orbital motion under Newton's law of gravitation is one of the great syntheses of classical mechanics.
\end{aretenir}

\newpage

\coursec{Methods and worked examples}

\coursec{Gravitational Fields}

\courssub{Newton's Law of Universal Gravitation and Kepler's Laws}

\begin{definition}[Newton's Law of Universal Gravitation]
Every particle of matter in the universe attracts every other particle with a force that is directly proportional to the product of their masses and inversely proportional to the square of the distance between their centres. The force acts along the line joining the centres and is always attractive.
\end{definition}

\begin{propriete}[Mathematical Form of Newton's Law]
For two point masses $m_1$ and $m_2$ separated by a distance $r$ (centre-to-centre), the magnitude of the gravitational force is
\[ F = \frac{G\,m_1 m_2}{r^{2}}, \]
where $G = 6.67\times 10^{-11}\ \mathrm{N\,m^{2}\,kg^{-2}}$ is the universal gravitational constant. In vector form, the force on $m_2$ due to $m_1$ is
\[ \vec{F}_{12} = -\frac{G\,m_1 m_2}{r^{2}}\,\hat{r}_{12}, \]
where $\hat{r}_{12}$ is the unit vector from $m_1$ to $m_2$. The force on $m_1$ due to $m_2$ is $\vec{F}_{21} = -\vec{F}_{12}$ (Newton's third law).
\end{propriete}

\begin{aretenir}[Key Points for Newton's Law]
\begin{itemize}
\item $G$ is a universal constant; its value is the same everywhere in the universe.
\item The law applies exactly to point masses and to spherically symmetric bodies (where all mass can be considered concentrated at the centre).
\item For non-spherical or non-symmetric bodies, the law must be applied by integration over mass elements.
\item Gravitational force is a \cle{central force} (always directed along the line joining centres) and a \cle{conservative force}.
\item The inverse-square nature means: if distance doubles, force becomes one-quarter; if distance triples, force becomes one-ninth.
\end{itemize}
\end{aretenir}

\begin{methode}[Calculating a gravitational force between two bodies]
\begin{enumerate}
\item \textbf{Identify} the two masses $m_1$, $m_2$ and the \emph{centre-to-centre} distance $r$ (never the surface-to-surface distance!). Convert everything to SI units (kg, m).
\item \textbf{Write} Newton's law:
\[ F = \frac{G\,m_1 m_2}{r^{2}}, \qquad G = 6.67\times 10^{-11}\ \mathrm{N\,m^{2}\,kg^{-2}}. \]
\item \textbf{Compute} the magnitude. The force is always \cle{attractive} and acts along the line joining the centres.
\item For \textbf{extended spherical bodies} (planets, stars), treat all the mass as concentrated at the centre — this is exact for spherically symmetric bodies (Shell Theorem).
\item \textbf{Check} the order of magnitude: gravitational forces between laboratory-sized objects are tiny ($10^{-10}$ to $10^{-7}$ N); they become large only when at least one mass is astronomical.
\end{enumerate}
\textbf{Reflex:} if the distance is doubled, the force is divided by $4$ (inverse-square law). Use ratios to avoid heavy computation.
\end{methode}

\begin{exoresolu}[Force between the Earth and the Moon]
The Earth has mass $M_E = 5.98\times 10^{24}\ \mathrm{kg}$ and the Moon has mass $M_M = 7.35\times 10^{22}\ \mathrm{kg}$. Their centres are $r = 3.84\times 10^{8}\ \mathrm{m}$ apart.
\begin{enumerate}
\item Calculate the gravitational force exerted by the Earth on the Moon.
\item State, with justification, the force exerted by the Moon on the Earth.
\item Show that this force provides exactly the centripetal force needed if the Moon orbits the Earth with period $T = 27.3$ days.
\end{enumerate}
\tcblower
\textbf{1. Force exerted by the Earth on the Moon.}

Both bodies are nearly spherical, so we apply Newton's law with the centre-to-centre distance:
\[ F = \frac{G M_E M_M}{r^{2}} = \frac{(6.67\times 10^{-11})(5.98\times 10^{24})(7.35\times 10^{22})}{(3.84\times 10^{8})^{2}}. \]
Numerator: $(6.67\times 10^{-11})(5.98\times 10^{24}) = 3.989\times 10^{14}$; then $(3.989\times 10^{14})(7.35\times 10^{22}) = 2.932\times 10^{37}\ \mathrm{N\,m^{2}}$.
Denominator: $(3.84\times 10^{8})^{2} = 1.4746\times 10^{17}\ \mathrm{m^{2}}$.
\[ F = \frac{2.932\times 10^{37}}{1.4746\times 10^{17}} = 1.99\times 10^{20}\ \mathrm{N}. \]
The force is attractive, directed from the Moon towards the centre of the Earth.

\textbf{2. Force of the Moon on the Earth.}

By Newton's \cle{third law}, the Moon attracts the Earth with a force of the \emph{same magnitude} $1.99\times 10^{20}\ \mathrm{N}$, opposite in direction (towards the Moon). The forces form an action–reaction pair: same magnitude, same line of action, opposite senses, acting on different bodies.

\textbf{3. Centripetal force check.}

For a circular orbit of radius $r$ and period $T$, the required centripetal force on the Moon is
\[ F_c = M_M r \omega^{2} = M_M r \left(\frac{2\pi}{T}\right)^{2}. \]
Convert the period: $T = 27.3\times 24\times 3600 = 2.359\times 10^{6}\ \mathrm{s}$.
\[ F_c = (7.35\times 10^{22})(3.84\times 10^{8})\left(\frac{2\pi}{2.359\times 10^{6}}\right)^{2}. \]
Compute the angular speed: $\omega = \dfrac{2\pi}{2.359\times 10^{6}} = 2.663\times 10^{-6}\ \mathrm{rad\,s^{-1}}$, so $\omega^{2} = 7.09\times 10^{-12}\ \mathrm{s^{-2}}$.
\[ F_c = (7.35\times 10^{22})(3.84\times 10^{8})(7.09\times 10^{-12}) = 2.00\times 10^{20}\ \mathrm{N}. \]
This equals the gravitational force found in part 1 (to within rounding). Hence gravity alone accounts for the Moon's orbital motion — historically this was Newton's great test of his law.
\end{exoresolu}

\begin{propriete}[Kepler's Laws of Planetary Motion]
\begin{enumerate}
\item \textbf{First Law (Law of Ellipses):} The orbit of a planet around the Sun is an ellipse with the Sun at one focus.
\item \textbf{Second Law (Law of Equal Areas):} A line segment joining a planet and the Sun sweeps out equal areas during equal intervals of time. This is a consequence of conservation of angular momentum.
\item \textbf{Third Law (Harmonic Law):} The square of the orbital period $T$ of a planet is proportional to the cube of the semi-major axis $a$ of its orbit:
\[ T^{2} \propto a^{3} \quad \text{or} \quad \frac{T^{2}}{a^{3}} = \text{constant for all planets orbiting the same star}. \]
For circular orbits, $a = r$ (orbital radius).
\end{enumerate}
\end{propriete}

\begin{experience}[Verifying Kepler's Third Law with Solar System Data]
Using the mean orbital radii and periods of the planets (data from NASA), compute $T^{2}/r^{3}$ for each planet. You will find the ratio is constant (within observational uncertainty), confirming Kepler's third law. This law was later derived by Newton from his law of gravitation, showing that the constant is $4\pi^{2}/(GM_{\text{Sun}})$.
\end{experience}

\begin{exoresolu}[Applying Kepler's Third Law]
The Earth orbits the Sun at a mean distance $r_E = 1.50\times 10^{11}\ \mathrm{m}$ with period $T_E = 365.25$ days. Mars has an orbital period $T_M = 687$ days. Assuming circular orbits, find the mean orbital radius of Mars.
\tcblower
Using Kepler's third law as a ratio (same central mass, the Sun):
\[ \left(\frac{T_M}{T_E}\right)^{2} = \left(\frac{r_M}{r_E}\right)^{3} \;\Rightarrow\; r_M = r_E \left(\frac{T_M}{T_E}\right)^{2/3}. \]
\[ \frac{T_M}{T_E} = \frac{687}{365.25} = 1.881, \quad (1.881)^{2/3} = 1.524. \]
\[ r_M = 1.50\times 10^{11} \times 1.524 = 2.29\times 10^{11}\ \mathrm{m}. \]
This matches the accepted value for Mars' semi-major axis.
\end{exoresolu}

\courssub{Gravitational Field Strength and Variation of $g$}

\begin{definition}[Gravitational Field Strength]
The gravitational field strength $\vec{g}$ at a point in a gravitational field is the gravitational force per unit mass experienced by a small test mass placed at that point:
\[ \vec{g} = \frac{\vec{F}}{m}, \qquad \text{unit: } \mathrm{N\,kg^{-1}} \text{ (equivalent to } \mathrm{m\,s^{-2}}\text{)}. \]
It is a vector quantity, directed towards the centre of the mass creating the field.
\end{definition}

\begin{propriete}[Field Strength of a Spherical Mass]
For a point mass $M$ or a spherically symmetric body of mass $M$, the gravitational field strength at a distance $r$ from its centre is
\[ g = \frac{GM}{r^{2}} \quad (r \geq R), \]
where $R$ is the radius of the body. At the surface ($r = R$), $g_0 = GM/R^{2}$. At altitude $h$, $r = R + h$ and $g = GM/(R+h)^{2}$.
\end{propriete}

\begin{aretenir}[Field Strength vs. Weight]
\begin{itemize}
\item \cle{Mass} $m$ is an intrinsic property of a body (invariant, unit kg).
\item \cle{Weight} $W = mg$ is the gravitational force on the body at a specific location (depends on local $g$, unit N).
\item Field strength $g$ is a property of the \emph{field} at a point; weight is a property of the \emph{body} in that field.
\end{itemize}
\end{aretenir}

\begin{methode}[Using $g = GM/r^{2}$ and relating field strength to weight]
\begin{enumerate}
\item The \cle{gravitational field strength} at a point is the force per unit mass: $\vec{g} = \vec{F}/m$, unit $\mathrm{N\,kg^{-1}}$ (equivalent to $\mathrm{m\,s^{-2}}$).
\item At distance $r$ from the centre of a spherical planet of mass $M$:
\[ g = \frac{GM}{r^{2}}. \]
At the surface, $r = R$ (planet's radius); at altitude $h$, $r = R + h$.
\item \textbf{Weight} of a body of mass $m$ at that point: $W = mg$.
\item For \textbf{ratio problems} (comparing two planets or two altitudes), write
\[ \frac{g_1}{g_2} = \frac{M_1}{M_2}\left(\frac{r_2}{r_1}\right)^{2} \]
— constants cancel, no need for $G$.
\item \textbf{Reflex:} mass $m$ is invariant; weight $W$ depends on the local $g$.
\end{enumerate}
\end{methode}

\begin{exoresolu}[Weight of a mountaineer on Mount Cameroon]
Take $M_E = 5.98\times 10^{24}\ \mathrm{kg}$, mean Earth radius $R_E = 6.37\times 10^{6}\ \mathrm{m}$. A climber of mass $m = 70.0\ \mathrm{kg}$ stands at the summit of Mount Cameroon, altitude $h = 4.10\times 10^{3}\ \mathrm{m}$.
\begin{enumerate}
\item Calculate $g$ at the summit and at sea level.
\item Calculate the climber's weight at both places and comment.
\end{enumerate}
\tcblower
\textbf{1. Field strengths.}

At sea level, $r = R_E$:
\[ g_0 = \frac{GM_E}{R_E^{2}} = \frac{(6.67\times 10^{-11})(5.98\times 10^{24})}{(6.37\times 10^{6})^{2}} = \frac{3.989\times 10^{14}}{4.058\times 10^{13}} = 9.83\ \mathrm{N\,kg^{-1}}. \]
At the summit, $r = R_E + h = 6.37\times 10^{6} + 4.10\times 10^{3} = 6.3741\times 10^{6}\ \mathrm{m}$:
\[ g_h = \frac{3.989\times 10^{14}}{(6.3741\times 10^{6})^{2}} = \frac{3.989\times 10^{14}}{4.063\times 10^{13}} = 9.82\ \mathrm{N\,kg^{-1}}. \]

\textbf{2. Weights.}

At sea level: $W_0 = mg_0 = 70.0\times 9.83 = 688\ \mathrm{N}$.

At the summit: $W_h = mg_h = 70.0\times 9.82 = 687\ \mathrm{N}$.

\textbf{Comment.} The decrease is about $0.1\ \%$: altitude variations of $g$ on Earth are small because $h \ll R_E$. The climber's \emph{mass} is unchanged ($70.0\ \mathrm{kg}$ everywhere); only the weight varies. Note that the value $g = 9.81\ \mathrm{m\,s^{-2}}$ used in class is a standard reference value; the local value depends on altitude and latitude.
\end{exoresolu}

\begin{popfigure}
\begin{tikzpicture}[scale=1.2]
\begin{axis}[
    width=0.9\linewidth,
    height=7cm,
    xlabel={Distance from centre $r$ (m)},
    ylabel={Gravitational field strength $g$ (N/kg)},
    xmin=0, xmax=2.5e7,
    ymin=0, ymax=10.5,
    xtick={0,6.37e6,1.274e7,1.911e7,2.548e7},
    xticklabels={$0$, $R_E$, $2R_E$, $3R_E$, $4R_E$},
    ytick={0,2.46,4.92,7.37,9.83},
    grid=major,
    grid style={popline},
    axis lines=middle,
    enlargelimits=false,
    clip=false
]
% Inside Earth: linear increase
\addplot[domain=0:6.37e6, samples=100, popblue, thick] {9.83*x/6.37e6};
% Outside Earth: inverse square decrease
\addplot[domain=6.37e6:2.548e7, samples=200, poporange, thick] {3.989e14/x^2};
% Vertical line at surface
\draw[popmuted, dashed] (axis cs:6.37e6,0) -- (axis cs:6.37e6,9.83);
% Labels
\node[popblue, anchor=south west] at (axis cs:3e6,5) {$g \propto r$ (inside)};
\node[poporange, anchor=north west] at (axis cs:1.2e7,4) {$g \propto 1/r^2$ (outside)};
\end{axis}
\end{tikzpicture}
\legende{Variation of gravitational field strength $g$ with distance $r$ from Earth's centre. Inside the Earth ($r<R_E$), $g$ increases linearly with $r$ assuming uniform density. Outside ($r>R_E$), $g$ follows an inverse-square law.}
\end{popfigure}

\begin{methode}[Choosing the correct formula for $g(r)$ inside and outside a uniform sphere]
\begin{enumerate}
\item \textbf{Locate the point:} is $r \geq R$ (outside) or $r \leq R$ (inside)? Sketch the situation.
\item \textbf{Outside} ($r \geq R$): the whole mass $M$ acts as if at the centre (Shell Theorem):
\[ g(r) = \frac{GM}{r^{2}} \quad (\text{inverse-square decrease}). \]
\item \textbf{Inside} ($r \leq R$), assuming uniform density $\rho$: only the mass within radius $r$ contributes (shells outside exert zero net force). The enclosed mass is $M_{\text{in}} = \tfrac{4}{3}\pi r^{3}\rho = M(r/R)^{3}$, giving
\[ g(r) = \frac{GM_{\text{in}}}{r^{2}} = \frac{GM}{R^{3}}\,r = \frac{4}{3}\pi G\rho\, r \quad (\text{linear increase with } r). \]
\item \textbf{Continuity check:} both formulas give $g = GM/R^{2}$ at $r = R$. At the centre, $g = 0$.
\item \textbf{Graph:} $g$ rises linearly from $0$ at $r=0$ to a maximum at $r=R$, then falls as $1/r^{2}$.
\end{enumerate}
\end{methode}

\begin{exoresolu}[Gravity in a deep mine]
Assume the Earth is a uniform sphere of mass $M = 5.98\times 10^{24}\ \mathrm{kg}$ and radius $R = 6.37\times 10^{6}\ \mathrm{m}$, with surface field $g_0 = 9.83\ \mathrm{N\,kg^{-1}}$.
\begin{enumerate}
\item Show that at depth $d$ below the surface, $g = g_0\left(1 - \dfrac{d}{R}\right)$.
\item Calculate $g$ at the bottom of a mine shaft $d = 3.0\ \mathrm{km}$ deep.
\item Compare with the value at altitude $h = 3.0\ \mathrm{km}$ above the surface. Which changes $g$ more?
\end{enumerate}
\tcblower
\textbf{1. Derivation.}

At depth $d$, the distance from the centre is $r = R - d < R$, so we use the interior formula:
\[ g(r) = \frac{GM}{R^{3}}\,r = \frac{GM}{R^{3}}(R - d). \]
Since $g_0 = \dfrac{GM}{R^{2}}$, we can write $\dfrac{GM}{R^{3}} = \dfrac{g_0}{R}$, hence
\[ g = \frac{g_0}{R}(R - d) = g_0\left(1 - \frac{d}{R}\right). \qquad \blacksquare \]

\textbf{2. Value at $d = 3.0\ \mathrm{km}$.}

$\dfrac{d}{R} = \dfrac{3.0\times 10^{3}}{6.37\times 10^{6}} = 4.71\times 10^{-4}$.
\[ g = 9.83\,(1 - 4.71\times 10^{-4}) = 9.83 - 0.0046 = 9.825\ \mathrm{N\,kg^{-1}}. \]

\textbf{3. Comparison with altitude $h = 3.0\ \mathrm{km}$.}

Outside the Earth, $g = \dfrac{GM}{(R+h)^{2}} = g_0\left(\dfrac{R}{R+h}\right)^{2}$. Since $h \ll R$, use the binomial approximation $\left(1 + \dfrac{h}{R}\right)^{-2} \approx 1 - \dfrac{2h}{R}$:
\[ g \approx g_0\left(1 - \frac{2h}{R}\right) = 9.83\,(1 - 9.42\times 10^{-4}) = 9.821\ \mathrm{N\,kg^{-1}}. \]
\textbf{Conclusion:} ascending $3.0\ \mathrm{km}$ reduces $g$ by about \emph{twice} as much as descending $3.0\ \mathrm{km}$. Going up, you move away from \emph{all} the mass; going down, part of the mass is left above you and its pull cancels.
\end{exoresolu}

\courssub{Gravitational Potential, Potential Energy, Escape and Orbital Speeds}

\begin{definition}[Gravitational Potential]
The gravitational potential $V$ at a point in a gravitational field is the work done per unit mass by an external agent in bringing a test mass from infinity to that point, without acceleration. For a point mass $M$ or a spherically symmetric body:
\[ V = -\frac{GM}{r} \quad (\mathrm{J\,kg^{-1}}), \]
with the zero of potential defined at infinity ($r \to \infty$). The potential is always negative for an attractive field.
\end{definition}

\begin{definition}[Gravitational Potential Energy]
The gravitational potential energy $E_p$ of a system consisting of a mass $m$ in the field of a mass $M$ is the work done by an external agent to assemble the system from infinite separation:
\[ E_p = mV = -\frac{GMm}{r} \quad (\mathrm{J}). \]
It is a property of the \emph{system} (both masses), not of the test mass alone.
\end{definition}

\begin{propriete}[Properties of $V$ and $E_p$]
\begin{itemize}
\item Both are \emph{scalar} quantities (easier to work with than vector forces).
\item Both are \emph{negative} for attractive fields, with zero at infinity.
\item The potential $V$ characterises the \emph{field} itself (energy per unit mass); $E_p$ characterises the \emph{body–field system}.
\item Difference in potential between two points $A$ and $B$:
\[ \Delta V = V_B - V_A = GM\left(\frac{1}{r_A} - \frac{1}{r_B}\right). \]
\item Work done by an external agent moving mass $m$ slowly from $A$ to $B$:
\[ W_{A\to B} = \Delta E_p = m\Delta V = GMm\left(\frac{1}{r_A} - \frac{1}{r_B}\right). \]
\item For a body moving freely in the field, total mechanical energy is conserved:
\[ E = \frac{1}{2}mv^{2} - \frac{GMm}{r} = \text{constant}. \]
\end{itemize}
\end{propriete}

\begin{attention}[Common Traps with Potential and Potential Energy]
\begin{itemize}
\item \textbf{Sign convention:} $V$ and $E_p$ are \emph{negative}. Moving away from the planet ($r$ increases) makes them \emph{less negative} (i.e. they \emph{increase} towards zero).
\item \textbf{Do not confuse} $V$ (field property, $\mathrm{J\,kg^{-1}}$) with $E_p$ (system property, $\mathrm{J}$).
\item \textbf{Do not use} $E_p = mgh$ for large altitude changes — it assumes constant $g$ and is only valid for $h \ll R$.
\item \textbf{Work done by gravity} is $W_{\text{grav}} = -\Delta E_p$. Work done \emph{against} gravity is $W_{\text{ext}} = +\Delta E_p$.
\end{itemize}
\end{attention}

\begin{methode}[Working with potential $V$ and potential energy $E_p$]
\begin{enumerate}
\item \textbf{Definitions} (zero at infinity, both always \emph{negative} for attractive fields):
\[ V = -\frac{GM}{r} \quad (\mathrm{J\,kg^{-1}}), \qquad E_p = mV = -\frac{GMm}{r}. \]
\item \textbf{Work done by an external agent} moving mass $m$ slowly from $A$ to $B$:
\[ W_{A\to B} = \Delta E_p = E_{p,B} - E_{p,A} = GMm\left(\frac{1}{r_A} - \frac{1}{r_B}\right). \]
\item \textbf{Sign reflex:} moving \emph{away} from the planet ($r_B > r_A$) requires positive external work; $E_p$ increases (becomes less negative) towards $0$.
\item \textbf{Energy conservation} for a body moving freely in the field:
\[ \frac{1}{2}mv^{2} - \frac{GMm}{r} = \text{constant}. \]
\item \textbf{Trap:} do not mix $V$ (per unit mass, property of the field) with $E_p$ (property of the body–planet system). Do not use $E_p = mgh$ over large altitude changes — it assumes constant $g$.
\end{enumerate}
\end{methode}

\begin{exoresolu}[Lifting a satellite into orbit]
A satellite of mass $m = 500\ \mathrm{kg}$ is taken from the Earth's surface to a circular orbit of radius $r = 7.00\times 10^{6}\ \mathrm{m}$. Take $M = 5.98\times 10^{24}\ \mathrm{kg}$, $R = 6.37\times 10^{6}\ \mathrm{m}$.
\begin{enumerate}
\item Calculate the gravitational potential at the surface and at the orbit.
\item Calculate the change in the satellite's gravitational potential energy.
\item Hence find the minimum work the launch system must supply against gravity.
\end{enumerate}
\tcblower
\textbf{1. Potentials.}

At the surface:
\[ V_R = -\frac{GM}{R} = -\frac{(6.67\times 10^{-11})(5.98\times 10^{24})}{6.37\times 10^{6}} = -\frac{3.989\times 10^{14}}{6.37\times 10^{6}} = -6.26\times 10^{7}\ \mathrm{J\,kg^{-1}}. \]
At the orbit:
\[ V_r = -\frac{GM}{r} = -\frac{3.989\times 10^{14}}{7.00\times 10^{6}} = -5.70\times 10^{7}\ \mathrm{J\,kg^{-1}}. \]
Both are negative; the potential is \emph{higher} (less negative) in orbit.

\textbf{2. Change in potential energy.}
\[ \Delta E_p = m(V_r - V_R) = 500\times\left[(-5.70\times 10^{7}) - (-6.26\times 10^{7})\right] = 500\times(5.6\times 10^{6}). \]
\[ \boxed{\Delta E_p = +2.8\times 10^{9}\ \mathrm{J}} \]
The positive sign confirms that energy must be supplied to move the satellite away from the Earth.

\textbf{3. Minimum work against gravity.}

If the satellite is moved slowly (no change in kinetic energy), the external work equals $\Delta E_p$:
\[ W_{\min} = 2.8\times 10^{9}\ \mathrm{J}. \]
\textbf{Remark:} in reality the rocket must also supply the \emph{kinetic} energy of the orbit, $\tfrac{1}{2}mv^{2} = \dfrac{GMm}{2r} = \dfrac{(3.989\times 10^{14})(500)}{2\times 7.00\times 10^{6}} = 1.42\times 10^{10}\ \mathrm{J}$ — actually larger than the potential-energy gain! The total mechanical energy to be supplied is about $1.70\times 10^{10}\ \mathrm{J}$.
\end{exoresolu}

\begin{popfigure}
\begin{tikzpicture}[scale=1.1]
\begin{axis}[
    width=0.9\linewidth,
    height=7cm,
    xlabel={Distance from centre $r$ (m)},
    ylabel={Gravitational potential $V$ (MJ/kg)},
    xmin=0, xmax=2.5e7,
    ymin=-70, ymax=5,
    xtick={0,6.37e6,1.274e7,1.911e7,2.548e7},
    xticklabels={$0$, $R_E$, $2R_E$, $3R_E$, $4R_E$},
    ytick={-62.6,-31.3,0},
    grid=major,
    grid style={popline},
    axis lines=middle,
    enlargelimits=false,
    clip=false
]
% Inside Earth: parabolic (for uniform density)
\addplot[domain=0:6.37e6, samples=100, popblue, thick] {-62.6*(3/2 - 0.5*(x/6.37e6)^2)};
% Outside Earth: -1/r
\addplot[domain=6.37e6:2.548e7, samples=200, poporange, thick] {-3.989e14/x/1e6};
% Vertical line at surface
\draw[popmuted, dashed] (axis cs:6.37e6,-70) -- (axis cs:6.37e6,5);
% Labels
\node[popblue, anchor=north east] at (axis cs:3e6,-40) {$V = -\frac{GM}{2R}(3-r^2/R^2)$};
\node[poporange, anchor=south west] at (axis cs:1.2e7,-20) {$V = -GM/r$};
\node[popmuted, anchor=north] at (axis cs:6.37e6,2) {$R_E$};
\end{axis}
\end{tikzpicture}
\legende{Gravitational potential $V$ vs distance $r$ for a uniform sphere. Inside ($r<R_E$), $V$ is parabolic; outside, it follows $-1/r$. The potential is continuous and differentiable at $r=R_E$.}
\end{popfigure}

\begin{methode}[Deriving and using $v_{esc}$ and $v_{orb}$]
\begin{enumerate}
\item \textbf{Escape velocity} (from surface $R$): set total mechanical energy to zero (body just reaches infinity at rest):
\[ \frac{1}{2}mv_{esc}^{2} - \frac{GMm}{R} = 0 \;\Rightarrow\; v_{esc} = \sqrt{\frac{2GM}{R}} = \sqrt{2g_0 R}. \]
\item \textbf{Orbital speed} for a circular orbit of radius $r$: equate gravity to centripetal force:
\[ \frac{GMm}{r^{2}} = \frac{mv^{2}}{r} \;\Rightarrow\; v_{orb} = \sqrt{\frac{GM}{r}}. \]
\item \textbf{Memory hook:} $v_{esc} = \sqrt{2}\,v_{orb}$ at the same radius.
\item \textbf{Period:} $T = \dfrac{2\pi r}{v_{orb}} = 2\pi\sqrt{\dfrac{r^{3}}{GM}}$ — this is Kepler's third law, $T^{2} \propto r^{3}$.
\item \textbf{Reflex:} both speeds are \emph{independent of the satellite's mass} $m$ — it cancels.
\end{enumerate}
\end{methode}

\begin{exoresolu}[Escape from the Earth and from the Moon]
Take $M_E = 5.98\times 10^{24}\ \mathrm{kg}$, $R_E = 6.37\times 10^{6}\ \mathrm{m}$; for the Moon $M_M = 7.35\times 10^{22}\ \mathrm{kg}$, $R_M = 1.74\times 10^{6}\ \mathrm{m}$.
\begin{enumerate}
\item Derive the expression for escape velocity from the surface of a planet.
\item Calculate the escape velocities of the Earth and of the Moon.
\item A low Earth-orbit satellite circles at $r = 6.77\times 10^{6}\ \mathrm{m}$. Find its orbital speed and period.
\end{enumerate}
\tcblower
\textbf{1. Derivation.}

Conservation of mechanical energy between the surface (speed $v_{esc}$) and infinity (speed $0$, potential $0$):
\[ \underbrace{\frac{1}{2}mv_{esc}^{2} - \frac{GMm}{R}}_{\text{surface}} = \underbrace{0 + 0}_{\text{infinity}} \;\Rightarrow\; v_{esc} = \sqrt{\frac{2GM}{R}}. \qquad \blacksquare \]

\textbf{2. Numerical values.}

Earth:
\[ v_{esc,E} = \sqrt{\frac{2(6.67\times 10^{-11})(5.98\times 10^{24})}{6.37\times 10^{6}}} = \sqrt{\frac{7.977\times 10^{14}}{6.37\times 10^{6}}} = \sqrt{1.252\times 10^{8}} = 1.12\times 10^{4}\ \mathrm{m\,s^{-1}} \approx 11.2\ \mathrm{km\,s^{-1}}. \]
Moon:
\[ v_{esc,M} = \sqrt{\frac{2(6.67\times 10^{-11})(7.35\times 10^{22})}{1.74\times 10^{6}}} = \sqrt{\frac{9.806\times 10^{12}}{1.74\times 10^{6}}} = \sqrt{5.636\times 10^{6}} = 2.37\times 10^{3}\ \mathrm{m\,s^{-1}} \approx 2.4\ \mathrm{km\,s^{-1}}. \]
The Moon's low escape velocity explains why it cannot retain an atmosphere: fast gas molecules escape easily.

\textbf{3. Orbital speed and period.}
\[ v_{orb} = \sqrt{\frac{GM_E}{r}} = \sqrt{\frac{3.989\times 10^{14}}{6.77\times 10^{6}}} = \sqrt{5.892\times 10^{7}} = 7.68\times 10^{3}\ \mathrm{m\,s^{-1}} \approx 7.7\ \mathrm{km\,s^{-1}}. \]
\[ T = \frac{2\pi r}{v_{orb}} = \frac{2\pi(6.77\times 10^{6})}{7.68\times 10^{3}} = 5.54\times 10^{3}\ \mathrm{s} \approx 92\ \mathrm{min}. \]
A low-orbit satellite circles the Earth in about an hour and a half — consistent with observations of the International Space Station.
\end{exoresolu}

\courssub{Geostationary Satellites and Kepler's Third Law}

\begin{definition}[Geostationary Satellite]
A geostationary satellite orbits the Earth in the \emph{equatorial plane}, in the \emph{same direction} as the Earth's rotation, with an orbital period exactly equal to the Earth's sidereal rotation period ($T = 23\ \mathrm{h}\ 56\ \mathrm{min}\ 4\ \mathrm{s} \approx 8.616\times 10^{4}\ \mathrm{s}$). It appears stationary relative to a fixed point on the equator.
\end{definition}

\begin{propriete}[Geostationary Orbit Parameters]
For a satellite of mass $m$ in a circular geostationary orbit around Earth (mass $M_E$):
\begin{itemize}
\item Orbital radius: $r = \left(\dfrac{GM_E T^{2}}{4\pi^{2}}\right)^{1/3}$
\item Altitude: $h = r - R_E$
\item Orbital speed: $v = \dfrac{2\pi r}{T} = \sqrt{\dfrac{GM_E}{r}}$
\end{itemize}
The values are independent of the satellite's mass $m$.
\end{propriete}

\begin{methode}[Solving geostationary / Kepler problems]
\begin{enumerate}
\item A \cle{geostationary} satellite: circular orbit in the \emph{equatorial plane}, same sense as the Earth's rotation, period $T = 24\ \mathrm{h}$ (sidereal day, $\approx 8.616\times 10^{4}\ \mathrm{s}$).
\item \textbf{Find the radius} from Kepler's third law:
\[ \frac{T^{2}}{r^{3}} = \frac{4\pi^{2}}{GM} \;\Rightarrow\; r = \left(\frac{GMT^{2}}{4\pi^{2}}\right)^{1/3}. \]
\item \textbf{Altitude:} $h = r - R_E$. \textbf{Speed:} $v = 2\pi r / T$.
\item \textbf{Comparing two satellites} around the same planet: use the ratio
\[ \left(\frac{T_1}{T_2}\right)^{2} = \left(\frac{r_1}{r_2}\right)^{3} \quad\text{(no need for } G \text{ or } M). \]
\item \textbf{Trap:} $T^{2}/r^{3} = \text{constant}$ holds only for satellites orbiting the \emph{same} central mass.
\end{enumerate}
\end{methode}

\begin{exoresolu}[Placing a communications satellite over Cameroon]
A telecommunications company wants a geostationary satellite to serve Central Africa. Take $M_E = 5.98\times 10^{24}\ \mathrm{kg}$, $R_E = 6.37\times 10^{6}\ \mathrm{m}$, sidereal day $T = 8.616\times 10^{4}\ \mathrm{s}$.
\begin{enumerate}
\item Calculate the radius and the altitude of the geostationary orbit.
\item Calculate the orbital speed of the satellite.
\item A second satellite has an orbital period of $12\ \mathrm{h}$. Using Kepler's third law only, find the ratio of its orbital radius to that of the geostationary satellite.
\end{enumerate}
\tcblower
\textbf{1. Radius and altitude.}

From Kepler's third law:
\[ r^{3} = \frac{GM T^{2}}{4\pi^{2}} = \frac{(3.989\times 10^{14})(8.616\times 10^{4})^{2}}{4\pi^{2}}. \]
Compute: $T^{2} = 7.424\times 10^{9}\ \mathrm{s^{2}}$; numerator $= 3.989\times 10^{14}\times 7.424\times 10^{9} = 2.961\times 10^{24}$; divide by $4\pi^{2} = 39.48$:
\[ r^{3} = 7.50\times 10^{22}\ \mathrm{m^{3}} \;\Rightarrow\; r = (7.50\times 10^{22})^{1/3} = 4.22\times 10^{7}\ \mathrm{m}. \]
Altitude:
\[ h = r - R_E = 4.22\times 10^{7} - 6.37\times 10^{6} = 3.58\times 10^{7}\ \mathrm{m} \approx 35\,800\ \mathrm{km}. \]

\textbf{2. Orbital speed.}
\[ v = \frac{2\pi r}{T} = \frac{2\pi(4.22\times 10^{7})}{8.616\times 10^{4}} = 3.08\times 10^{3}\ \mathrm{m\,s^{-1}} \approx 3.1\ \mathrm{km\,s^{-1}}. \]
The satellite moves at about $3\ \mathrm{km\,s^{-1}}$ yet appears fixed in the sky, because the Earth rotates beneath it at the same angular rate — a dish antenna in Yaoundé, once pointed at it, never needs to move.

\textbf{3. Ratio for the 12-hour satellite.}

Both satellites orbit the Earth, so Kepler's third law applies as a ratio:
\[ \frac{r_2}{r_1} = \left(\frac{T_2}{T_1}\right)^{2/3} = \left(\frac{12}{24}\right)^{2/3} = (0.5)^{2/3} = 0.63. \]
So $r_2 \approx 0.63\times 4.22\times 10^{7} = 2.66\times 10^{7}\ \mathrm{m}$. (This is roughly the altitude class of GPS-type navigation satellites.)
\end{exoresolu}

\begin{popfigure}
\begin{tikzpicture}[scale=1.1]
% Earth
\fill[popblueL] (0,0) circle (1.5);
\draw[popblue, thick] (0,0) circle (1.5);
\node at (0,0) {Earth};
% Geostationary orbit
\draw[poporange, thick, dashed] (0,0) circle (5);
\node[poporange] at (5,0) [anchor=west] {Geostationary orbit ($r \approx 4.2\times 10^7$ m)};
% Low Earth orbit
\draw[popgreen, thick, dashed] (0,0) circle (2.2);
\node[popgreen] at (2.2,0) [anchor=west] {LEO ($r \approx 6.8\times 10^6$ m)};
% GPS orbit (12h)
\draw[poppurple, thick, dashed] (0,0) circle (3.15);
\node[poppurple] at (3.15,0) [anchor=west] {GPS / 12h orbit ($r \approx 2.7\times 10^7$ m)};
% Satellite symbols
\fill[poporange] (5,0) circle (0.08);
\fill[popgreen] (2.2,0) circle (0.06);
\fill[poppurple] (3.15,0) circle (0.06);
% Equatorial plane label
\draw[popmuted, <->] (-1.5,-1.5) -- (1.5,-1.5);
\node[popmuted] at (0,-1.8) {Equatorial plane};
\end{tikzpicture}
\legende{Comparison of satellite orbits around Earth (not to scale). Geostationary orbit is at $\approx 35\,800$ km altitude; GPS satellites orbit at $\approx 20\,200$ km (12 h period); Low Earth Orbit (LEO) satellites like the ISS orbit at $\approx 400$ km (90 min period).}
\end{popfigure}

\courssub{Motion in the Gravitational Field: Projectile Motion}

\begin{definition}[Projectile Motion in a Uniform Gravitational Field]
Near the Earth's surface, the gravitational field is approximately \cle{uniform} (constant magnitude $g \approx 9.81\ \mathrm{m\,s^{-2}}$, direction vertically downwards). A projectile is any body given an initial velocity and then moving under the influence of gravity alone (air resistance neglected). Its trajectory is a parabola.
\end{definition}

\begin{propriete}[Equations of Projectile Motion]
Choose axes: $Ox$ horizontal, $Oy$ vertical (upwards positive); origin at launch point. Initial velocity $\vec{v}_0$ at angle $\theta$ to the horizontal.
\begin{align*}
\text{Components:} &\quad v_{0x} = v_0\cos\theta, \quad v_{0y} = v_0\sin\theta. \\
\text{Acceleration:} &\quad a_x = 0, \quad a_y = -g. \\
\text{Velocity:} &\quad v_x(t) = v_0\cos\theta, \quad v_y(t) = v_0\sin\theta - gt. \\
\text{Position:} &\quad x(t) = v_0\cos\theta\; t, \quad y(t) = v_0\sin\theta\; t - \tfrac{1}{2}gt^{2}. \\
\text{Trajectory (eliminate $t$):} &\quad y = x\tan\theta - \frac{g\,x^{2}}{2v_0^{2}\cos^{2}\theta} \quad (\text{parabola}).
\end{align*}
For launch and landing at the same height ($y=0$):
\begin{align*}
\text{Time of flight:} &\quad T = \frac{2v_0\sin\theta}{g}. \\
\text{Range:} &\quad R = \frac{v_0^{2}\sin 2\theta}{g}. \\
\text{Maximum height:} &\quad H = \frac{v_0^{2}\sin^{2}\theta}{2g}.
\end{align*}
Maximum range on level ground occurs at $\theta = 45^{\circ}$.
\end{propriete}

\begin{aretenir}[Key Reflexes for Projectile Problems]
\begin{itemize}
\item Horizontal velocity is \emph{constant} ($a_x=0$); vertical motion is uniformly accelerated ($a_y=-g$).
\item Time $t$ is the \emph{same} for horizontal and vertical motions — it links the two.
\item At the top of the trajectory, $v_y = 0$ but $v_x \neq 0$; acceleration is still $g$ downwards.
\item The trajectory is symmetric about the maximum height \emph{only} if launch and landing heights are equal.
\item Energy conservation: $\frac{1}{2}mv^{2} + mgy = \frac{1}{2}mv_0^{2}$ (useful for speed at a given height).
\end{itemize}
\end{aretenir}

\begin{methode}[Solving projectile problems]
\begin{enumerate}
\item \textbf{Choose axes:} $Ox$ horizontal, $Oy$ vertical (upwards); origin at launch point. Near the Earth's surface $g$ is constant: the field is \cle{uniform}.
\item \textbf{Resolve the initial velocity:} $v_{0x} = v_0\cos\theta$, $v_{0y} = v_0\sin\theta$.
\item \textbf{Write the equations of motion} (only acceleration is $-g$ along $y$):
\[ x(t) = v_0\cos\theta\; t, \qquad y(t) = v_0\sin\theta\; t - \tfrac{1}{2}gt^{2}. \]
\item \textbf{Trajectory:} eliminate $t$:
\[ y = x\tan\theta - \frac{g\,x^{2}}{2v_0^{2}\cos^{2}\theta} \quad (\text{parabola}). \]
\item \textbf{Key results} (level ground):
\[ \text{range } R = \frac{v_0^{2}\sin 2\theta}{g}, \qquad \text{max height } H = \frac{v_0^{2}\sin^{2}\theta}{2g}, \qquad \text{time of flight } T = \frac{2v_0\sin\theta}{g}. \]
\item \textbf{Reflexes:} horizontal velocity is constant; at the summit $v_y = 0$ but $v_x \neq 0$; maximum range on level ground occurs at $\theta = 45^{\circ}$.
\end{enumerate}
\end{methode}

\begin{exoresolu}[A football free kick]
During a match in Bafoussam, a player kicks a ball from ground level with speed $v_0 = 20.0\ \mathrm{m\,s^{-1}}$ at $\theta = 30.0^{\circ}$ above the horizontal. Take $g = 9.81\ \mathrm{m\,s^{-2}}$ and neglect air resistance.
\begin{enumerate}
\item Write the time equations $x(t)$ and $y(t)$.
\item Calculate the time of flight and the range.
\item Calculate the maximum height reached.
\item Determine the velocity (magnitude and direction) of the ball at the top of its flight.
\end{enumerate}
\tcblower
\textbf{1. Time equations.}

$v_{0x} = 20.0\cos 30.0^{\circ} = 20.0\times 0.866 = 17.3\ \mathrm{m\,s^{-1}}$; \quad $v_{0y} = 20.0\sin 30.0^{\circ} = 10.0\ \mathrm{m\,s^{-1}}$.
\[ x(t) = 17.3\,t, \qquad y(t) = 10.0\,t - \tfrac{1}{2}(9.81)t^{2} = 10.0\,t - 4.91\,t^{2} \quad (\text{in metres, } t \text{ in s}). \]

\textbf{2. Time of flight and range.}

The ball lands when $y = 0$ (excluding $t = 0$):
\[ t(10.0 - 4.91\,t) = 0 \;\Rightarrow\; T = \frac{10.0}{4.91} = 2.04\ \mathrm{s}. \]
(Check with the formula: $T = \dfrac{2v_0\sin\theta}{g} = \dfrac{2\times 20.0\times 0.500}{9.81} = 2.04\ \mathrm{s}$. \checkmark)
\[ R = v_{0x}\,T = 17.3\times 2.04 = 35.3\ \mathrm{m}. \]
(Check: $R = \dfrac{v_0^{2}\sin 2\theta}{g} = \dfrac{400\times\sin 60.0^{\circ}}{9.81} = \dfrac{400\times 0.866}{9.81} = 35.3\ \mathrm{m}$. \checkmark)

\textbf{3. Maximum height.}

At the summit $v_y = 0$: $v_{0y} - gt = 0 \Rightarrow t = \dfrac{10.0}{9.81} = 1.02\ \mathrm{s}$ (half the flight time, as expected from symmetry).
\[ H = 10.0(1.02) - 4.91(1.02)^{2} = 10.2 - 5.11 = 5.10\ \mathrm{m}. \]
(Check: $H = \dfrac{v_0^{2}\sin^{2}\theta}{2g} = \dfrac{400\times 0.250}{19.62} = 5.10\ \mathrm{m}$. \checkmark)

\textbf{4. Velocity at the top.}

At the summit the vertical component vanishes; only the (constant) horizontal component remains:
\[ \vec{v}_{\text{top}} = v_{0x}\,\vec{\imath} = 17.3\ \mathrm{m\,s^{-1}}, \quad \text{directed horizontally, in the sense of motion.} \]
\begin{attention}[Common traps in projectile questions]
\begin{itemize}
\item The acceleration at the top is \emph{not} zero: it is $g = 9.81\ \mathrm{m\,s^{-2}}$ downwards throughout the flight. Only $v_y$ is zero at the summit.
\item Never mix components: time $t$ is the \emph{same} in $x(t)$ and $y(t)$ — it is the bridge between the two motions.
\item The range formula $R = v_0^{2}\sin 2\theta / g$ is valid \emph{only} when launch and landing are at the same height.
\item For launch from a height $h_0$ above ground, solve $y(t) = -h_0$ for $t$ to find the time of flight.
\end{itemize}
\end{attention}
\end{exoresolu}

\begin{popfigure}
\begin{tikzpicture}[scale=1.1]
\begin{axis}[
    width=0.9\linewidth,
    height=7cm,
    xlabel={Horizontal distance $x$ (m)},
    ylabel={Vertical height $y$ (m)},
    xmin=0, xmax=40,
    ymin=0, ymax=6,
    xtick={0,10,20,30,35.3},
    ytick={0,1,2,3,4,5,5.1},
    grid=major,
    grid style={popline},
    axis lines=middle,
    enlargelimits=false,
    clip=false
]
% Trajectory
\addplot[domain=0:35.3, samples=100, popblue, thick] {x*tan(30) - 9.81*x^2/(2*20^2*cos(30)^2)};
% Key points
\fill[poporange] (0,0) circle (0.08) node[below left] {Launch};
\fill[poporange] (17.65,5.1) circle (0.08) node[above] {Max height $H$};
\fill[poporange] (35.3,0) circle (0.08) node[below right] {Range $R$};
% Velocity vectors at launch
\draw[popgreen, ->, thick] (0,0) -- (1.73,1.0) node[above right] {$\vec{v}_0$};
\draw[popgreen, ->] (0,0) -- (1.73,0) node[below] {$v_{0x}$};
\draw[popgreen, ->] (0,0) -- (0,1.0) node[left] {$v_{0y}$};
% Velocity at top
\draw[poppurple, ->, thick] (17.65,5.1) -- (19.38,5.1) node[right] {$\vec{v}_{\text{top}} = v_{0x}$};
% Acceleration vector
\draw[popred, ->, thick] (17.65,5.1) -- (17.65,4.1) node[right] {$\vec{g}$};
\end{axis}
\end{tikzpicture}
\legende{Projectile trajectory for $v_0=20\ \mathrm{m/s}$, $\theta=30^\circ$. The velocity vector is tangent to the path; at the top it is purely horizontal. Acceleration $\vec{g}$ is constant and vertically downward everywhere.}
\end{popfigure}

\begin{savaistu}[Did you know?]
\begin{itemize}
\item The concept of a geostationary orbit was first proposed by Arthur C. Clarke in 1945 in a paper titled "Extra-Terrestrial Relays". Today, the geostationary orbit is often called the "Clarke Orbit" in his honour.
\item Cameroon (like all equatorial countries) benefits greatly from geostationary satellites for telecommunications, weather monitoring, and broadcasting. A dish pointed at a geostationary satellite in Yaoundé never needs to track it.
\item The escape velocity from the Solar System at Earth's orbit is about $42\ \mathrm{km/s}$, but Earth's orbital speed is $30\ \mathrm{km/s}$. Space probes use "gravity assists" (slingshot manoeuvres) around planets to gain the extra speed needed to leave the Solar System.
\item The value $g = 9.81\ \mathrm{m/s^2}$ is a standard average. At the poles, $g \approx 9.83\ \mathrm{m/s^2}$; at the equator, $g \approx 9.78\ \mathrm{m/s^2}$ due to Earth's rotation and equatorial bulge.
\end{itemize}
\end{savaistu}

\begin{exercice}[Practice Problems]
\begin{enumerate}
\item Two identical spheres of mass $m = 50\ \mathrm{kg}$ each are placed with their centres $1.0\ \mathrm{m}$ apart. Calculate the gravitational force between them. Compare this force with the weight of one sphere.
\item A planet has twice the mass of Earth and twice the radius of Earth. What is the gravitational field strength at its surface in terms of $g_E$?
\item A satellite of mass $200\ \mathrm{kg}$ is in a circular orbit of radius $8.0\times 10^{6}\ \mathrm{m}$ around Earth. Calculate its (a) orbital speed, (b) period, (c) kinetic energy, (d) potential energy, (e) total mechanical energy.
\item Derive an expression for the minimum energy required to move a satellite of mass $m$ from a circular orbit of radius $r_1$ to a higher circular orbit of radius $r_2$ ($r_2 > r_1$).
\item A projectile is launched from a cliff of height $h = 50\ \mathrm{m}$ with speed $v_0 = 30\ \mathrm{m/s}$ at $\theta = 40^\circ$ above horizontal. Find (a) the time of flight, (b) the horizontal range, (c) the maximum height above the cliff, (d) the speed just before impact.
\item A geostationary satellite orbits at $r = 4.22\times 10^{7}\ \mathrm{m}$. A spy satellite orbits at $r = 7.00\times 10^{6}\ \mathrm{m}$. How many orbits does the spy satellite complete in one day?
\end{enumerate}
\end{exercice}

\begin{corrige}[Exercise 1]
$F = \frac{G m^2}{r^2} = \frac{(6.67\times 10^{-11})(50)^2}{1^2} = 1.67\times 10^{-7}\ \mathrm{N}$. Weight $W = mg = 50\times 9.81 = 490\ \mathrm{N}$. The gravitational force is $\approx 3.4\times 10^{-10}$ times the weight — utterly negligible.
\end{corrige}

\begin{corrige}[Exercise 2]
$g_p = \frac{G(2M_E)}{(2R_E)^2} = \frac{2}{4}\frac{GM_E}{R_E^2} = \frac{1}{2}g_E$.
\end{corrige}

\begin{corrige}[Exercise 3]
(a) $v = \sqrt{GM/r} = \sqrt{3.989\times 10^{14}/8.0\times 10^6} = 7.06\times 10^3\ \mathrm{m/s}$.
(b) $T = 2\pi r/v = 7.12\times 10^3\ \mathrm{s} \approx 119\ \mathrm{min}$.
(c) $K = \frac{1}{2}mv^2 = 4.98\times 10^9\ \mathrm{J}$.
(d) $U = -GMm/r = -9.97\times 10^9\ \mathrm{J}$.
(e) $E = K+U = -4.98\times 10^9\ \mathrm{J} = -K = U/2$ (virial theorem for circular orbits).
\end{corrige}

\begin{corrige}[Exercise 4]
Total energy in circular orbit: $E = -\frac{GMm}{2r}$. Energy required: $\Delta E = E_2 - E_1 = \frac{GMm}{2}\left(\frac{1}{r_1} - \frac{1}{r_2}\right)$. This is positive since $r_2 > r_1$.
\end{corrige}

\begin{corrige}[Exercise 5]
(a) Solve $y(t) = -50 = 30\sin40^\circ\, t - 4.905 t^2 \Rightarrow 4.905 t^2 - 19.28 t - 50 = 0$. Positive root: $t = 5.96\ \mathrm{s}$.
(b) $R = v_{0x} t = 30\cos40^\circ \times 5.96 = 137\ \mathrm{m}$.
(c) $H_{\text{max}} = \frac{(30\sin40^\circ)^2}{2g} = 18.9\ \mathrm{m}$ above cliff.
(d) $v_y = v_{0y} - gt = 19.28 - 9.81\times 5.96 = -39.2\ \mathrm{m/s}$; $v_x = 22.98\ \mathrm{m/s}$; $v = \sqrt{v_x^2+v_y^2} = 45.5\ \mathrm{m/s}$.
\end{corrige}

\begin{corrige}[Exercise 6]
$T_{\text{geo}} = 24\ \mathrm{h}$. $T_{\text{spy}} = 2\pi\sqrt{r^3/GM} = 2\pi\sqrt{(7.00\times 10^6)^3/3.989\times 10^{14}} = 5.54\times 10^3\ \mathrm{s} = 1.54\ \mathrm{h}$. Orbits per day $= 24/1.54 = 15.6$ orbits.
\end{corrige}

\newpage

\coursec{Exercises}

\courssub{I check my knowledge}

\begin{exercice}[MCQ: Newton's law of gravitation]
Two point masses $m_1$ and $m_2$ are separated by a distance $r$. The gravitational force between them is $F$. If the distance is doubled and each mass is tripled, the new force is:
\begin{enumerate}
\item $\dfrac{3F}{4}$
\item $\dfrac{9F}{4}$
\item $\dfrac{9F}{2}$
\item $9F$
\end{enumerate}
\end{exercice}

\begin{exercice}[True or false?]
State whether each of the following statements is true or false, and justify your answer in one or two sentences.
\begin{enumerate}
\item The gravitational field strength at the centre of the Earth is zero.
\item The gravitational potential at a point near the Earth is a positive quantity.
\item A geostationary satellite can be placed in orbit directly above Douala.
\item The escape velocity from a planet depends on the mass of the object being launched.
\end{enumerate}
\end{exercice}

\begin{exercice}[Definitions]
Define each of the following terms precisely, giving the defining formula and SI unit where appropriate:
\begin{enumerate}
\item gravitational field strength at a point;
\item gravitational potential at a point;
\item escape velocity;
\item geostationary satellite.
\end{enumerate}
\end{exercice}

\begin{exercice}[Direct calculation: force between lead spheres]
Two small lead spheres of masses $m_1 = \SI{50}{kg}$ and $m_2 = \SI{80}{kg}$ have their centres $\SI{0.50}{m}$ apart. Take $G = \SI{6.67e-11}{N.m^2.kg^{-2}}$.
\begin{enumerate}
\item Calculate the gravitational force of attraction between the spheres.
\item Compare this force with the weight of a sheet of paper of mass $\SI{1.0}{g}$ (take $g = \SI{9.81}{m.s^{-2}}$).
\item Hence explain why gravitational forces between everyday objects are never felt in daily life.
\end{enumerate}
\end{exercice}

\courssub{I practise}

\begin{exercice}[Variation of $g$ with position]
Take the mass of the Earth $M_E = \SI{5.98e24}{kg}$, its radius $R_E = \SI{6.38e6}{m}$ and $G = \SI{6.67e-11}{N.m^2.kg^{-2}}$. Assume the Earth has uniform density.
\begin{enumerate}
\item Calculate the gravitational field strength $g_0$ at the Earth's surface.
\item Calculate the value of $g$ at an altitude $h = 2R_E$ above the surface.
\item Calculate the value of $g$ at a depth $d = R_E/2$ below the surface.
\item Sketch the graph of $g(r)$ for $r$ ranging from $0$ to $3R_E$, marking the values found above.
\end{enumerate}
\end{exercice}

\begin{exercice}[A strange planet]
A planet has twice the mass of the Earth and half its radius.
\begin{enumerate}
\item Show that the gravitational field strength at its surface is $g_P = 8\,g_0$, where $g_0$ is the field strength at the Earth's surface.
\item Show that the escape velocity from this planet is $v_P = 2\,v_E$, where $v_E$ is the escape velocity from the Earth.
\item Taking $g_0 = \SI{9.81}{m.s^{-2}}$ and $v_E = \SI{11.2}{km.s^{-1}}$, give the numerical values of $g_P$ and $v_P$.
\end{enumerate}
\end{exercice}

\begin{exercice}[A low-orbit satellite]
A satellite is in a circular orbit around the Earth with a period of $90$ minutes. Take $M_E = \SI{5.98e24}{kg}$ and $G = \SI{6.67e-11}{N.m^2.kg^{-2}}$.
\begin{enumerate}
\item Using Kepler's third law, calculate the radius of the orbit.
\item Calculate the orbital speed of the satellite.
\item State, with justification, whether this satellite can be geostationary.
\end{enumerate}
\end{exercice}

\begin{exercice}[Projectile from a cliff in Limbe]
A stone is thrown horizontally with a speed of $\SI{15}{m.s^{-1}}$ from the top of a cliff $\SI{45}{m}$ high overlooking the sea at Limbe. Air resistance is negligible; take $g = \SI{9.81}{m.s^{-2}}$.
\begin{enumerate}
\item Write down the equations of motion $x(t)$ and $y(t)$ of the stone, choosing axes clearly.
\item Calculate the time taken for the stone to reach the sea.
\item Calculate the horizontal distance from the foot of the cliff to the point of impact.
\item Determine the velocity (magnitude and direction below the horizontal) of the stone at impact.
\end{enumerate}
\end{exercice}

\courssub{I prepare for the GCE}

\begin{exercice}[Geostationary satellites (GCE style)]
\begin{enumerate}
\item State Kepler's three laws of planetary motion. \hfill (3 marks)
\item Define a \emph{geostationary} satellite and state two of its uses for a country such as Cameroon. \hfill (3 marks)
\item Starting from Newton's law of gravitation and the condition for uniform circular motion, show that for a satellite of mass $m$ in a circular orbit of radius $r$ around the Earth (mass $M_E$),
\[ T^2 = \frac{4\pi^2 r^3}{G M_E}. \] \hfill (4 marks)
\item Hence calculate the orbital radius, the altitude above the Earth's surface, and the orbital speed of a geostationary satellite. (Take $G = \SI{6.67e-11}{N.m^2.kg^{-2}}$, $M_E = \SI{5.98e24}{kg}$, $R_E = \SI{6.38e6}{m}$, $T = \SI{24}{h}$.) \hfill (6 marks)
\item Explain, with the aid of a diagram, why a network of only three geostationary satellites can provide communication coverage of (almost) the whole Earth, and state which regions are poorly covered. \hfill (4 marks)
\end{enumerate}
\end{exercice}

\begin{exercice}[Energy of an orbiting satellite (GCE style)]
A satellite of mass $m = \SI{100}{kg}$ is in a circular orbit of radius $r_1 = 2R_E$ around the Earth.
\begin{enumerate}
\item Define gravitational potential and gravitational potential energy. \hfill (3 marks)
\item Show that the total mechanical energy of a satellite in a circular orbit of radius $r$ is
\[ E = -\frac{G M_E m}{2r}. \] \hfill (4 marks)
\item Explain the physical significance of the negative sign of this energy. \hfill (2 marks)
\item Calculate the work that must be supplied to transfer the satellite from the orbit of radius $r_1 = 2R_E$ to a higher circular orbit of radius $r_2 = 3R_E$. (Take $G = \SI{6.67e-11}{N.m^2.kg^{-2}}$, $M_E = \SI{5.98e24}{kg}$, $R_E = \SI{6.38e6}{m}$.) \hfill (5 marks)
\item The satellite is then allowed to fall back to Earth. Explain qualitatively what happens to its kinetic energy, potential energy and total mechanical energy if air resistance is taken into account. \hfill (3 marks)
\end{enumerate}
\end{exercice}

\begin{exercice}[Projectile motion and data analysis (GCE style)]
\textbf{Part A.} A football is kicked from level ground with an initial speed of $\SI{20}{m.s^{-1}}$ at an angle of $40^{\circ}$ above the horizontal. Air resistance is negligible; take $g = \SI{9.81}{m.s^{-2}}$.
\begin{enumerate}
\item Resolve the initial velocity into horizontal and vertical components. \hfill (2 marks)
\item Calculate the time of flight of the ball. \hfill (3 marks)
\item Calculate the maximum height reached. \hfill (3 marks)
\item Calculate the horizontal range. \hfill (2 marks)
\item State the angle of projection that gives the maximum range for a given initial speed, and justify your answer using the range formula. \hfill (2 marks)
\end{enumerate}
\textbf{Part B.} The table below gives the mean orbital radius $r$ and the period $T$ of four moons of Jupiter.
\begin{center}
\begin{tabularx}{\linewidth}{|l|X|X|X|X|}
\hline
\rowcolor{popblueL} Moon & Io & Europa & Ganymede & Callisto \\
\hline $r$ ($\times 10^{8}$ m) & 4.22 & 6.71 & 10.7 & 18.8 \\
\hline $T$ (days) & 1.77 & 3.55 & 7.15 & 16.7 \\
\hline
\end{tabularx}
\end{center}
\begin{enumerate}
\item Copy and complete the table by computing $T^2$ (in $\mathrm{s^2}$) and $r^3$ (in $\mathrm{m^3}$). \hfill (3 marks)
\item Plot the graph of $T^2$ against $r^3$ and comment on its shape. \hfill (3 marks)
\item Use the gradient of the graph, together with Kepler's third law, to deduce the mass of Jupiter. (Take $G = \SI{6.67e-11}{N.m^2.kg^{-2}}$.) \hfill (4 marks)
\end{enumerate}
\end{exercice}

\newpage

\coursec{Solutions to the exercises}

\begin{corrige}[Exercise 1 — MCQ: Newton's law of gravitation]
The initial force is $F = \dfrac{G m_1 m_2}{r^2}$. After the changes, the masses become $3m_1$ and $3m_2$ and the separation becomes $2r$:
\[
F' = \frac{G(3m_1)(3m_2)}{(2r)^2} = \frac{9 G m_1 m_2}{4 r^2} = \frac{9}{4}F.
\]
\textbf{Correct answer: (b)} $\dfrac{9F}{4}$.

\textbf{Examiner's tip:} the masses contribute a factor $3\times 3 = 9$ (numerator), the distance a factor $2^2 = 4$ (denominator). A common error is to forget to square the distance factor, which would wrongly give $\frac{9F}{2}$.
\end{corrige}

\begin{corrige}[Exercise 2 — True or false?]
\begin{enumerate}
\item \textbf{True.} Inside a uniform spherical Earth, $g(r) = \dfrac{GM_E}{R_E^3}\,r$; at the centre $r = 0$, so $g = 0$. (Equivalently, by symmetry the pulls from all directions cancel.)
\item \textbf{False.} Gravitational potential is defined by $V = -\dfrac{GM}{r}$, with the zero chosen at infinity. Since the field is attractive, $V$ is \emph{negative} at every finite distance.
\item \textbf{False.} A geostationary satellite must orbit in the \emph{equatorial plane} (above latitude $0^{\circ}$). Douala lies at about $4^{\circ}$\,N, so no satellite can remain fixed above it.
\item \textbf{False.} From $\frac12 m v_e^2 = \dfrac{GMm}{R}$, the test mass $m$ cancels: $v_e = \sqrt{2GM/R}$ depends only on the planet's mass and radius, not on the mass of the launched object.
\end{enumerate}
\end{corrige}

\begin{corrige}[Exercise 3 — Definitions]
\begin{enumerate}
\item \textbf{Gravitational field strength} at a point is the gravitational force per unit mass acting on a small test mass placed at that point:
\[ g = \frac{F}{m} \qquad \text{SI unit: } \mathrm{N\,kg^{-1}}\ (\text{equivalent to } \mathrm{m\,s^{-2}}). \]
It is a vector, directed towards the attracting mass.
\item \textbf{Gravitational potential} at a point is the work done per unit mass by an external agent in bringing a small test mass from infinity to that point (zero of potential at infinity):
\[ V = -\frac{GM}{r} \qquad \text{SI unit: } \mathrm{J\,kg^{-1}}. \]
It is a scalar, always negative.
\item \textbf{Escape velocity} is the minimum speed with which an object must be launched from the surface of a planet so that it reaches infinity with zero residual speed:
\[ v_e = \sqrt{\frac{2GM}{R}} = \sqrt{2gR} \qquad \text{SI unit: } \mathrm{m\,s^{-1}}. \]
\item \textbf{Geostationary satellite:} a satellite in a circular orbit in the Earth's equatorial plane, with a period equal to the Earth's rotation period ($T = 24$ h) and revolving in the same sense as the Earth, so that it appears permanently fixed above one point of the equator (altitude $\approx \SI{3.6e4}{km}$).
\end{enumerate}
\end{corrige}

\begin{corrige}[Exercise 4 — Direct calculation: force between lead spheres]
\begin{enumerate}
\item Applying Newton's law:
\[
F = \frac{G m_1 m_2}{r^2} = \frac{(6.67\times10^{-11})(50)(80)}{(0.50)^2}
= \frac{2.668\times10^{-7}}{0.25}
\approx \SI{1.07e-6}{N}.
\]
\item Weight of the sheet of paper:
\[
W = mg = (1.0\times10^{-3})(9.81) = \SI{9.81e-3}{N}.
\]
Ratio: $\dfrac{W}{F} = \dfrac{9.81\times10^{-3}}{1.07\times10^{-6}} \approx 9.2\times10^{3}$. The force between the spheres is about \textbf{ten thousand times smaller} than the weight of a single sheet of paper.
\item Because $G$ is extremely small ($\sim 10^{-11}$ in SI units), gravitational forces between ordinary objects (kilogram masses at metre distances) are of order $10^{-7}$\,N or less — far below friction and other contact forces, and undetectable by our senses. Gravitation only becomes appreciable when at least one of the masses is astronomical (e.g.\ the Earth).
\end{enumerate}
\end{corrige}

\begin{corrige}[Exercise 5 — Variation of $g$ with position]
\begin{enumerate}
\item At the surface:
\[
g_0 = \frac{GM_E}{R_E^2} = \frac{(6.67\times10^{-11})(5.98\times10^{24})}{(6.38\times10^{6})^2}
= \frac{3.989\times10^{14}}{4.070\times10^{13}}
\approx \SI{9.80}{N.kg^{-1}}.
\]
\item At altitude $h = 2R_E$, the distance from the centre is $r = R_E + h = 3R_E$:
\[
g = \frac{GM_E}{(3R_E)^2} = \frac{g_0}{9} \approx \frac{9.80}{9} \approx \SI{1.09}{N.kg^{-1}}.
\]
\item Inside a uniform Earth, $g(r) = \dfrac{GM_E}{R_E^3}\,r = g_0\,\dfrac{r}{R_E}$. At depth $d = R_E/2$, $r = R_E - d = R_E/2$:
\[
g = g_0 \times \frac{1}{2} \approx \SI{4.90}{N.kg^{-1}}.
\]
\item Graph: $g$ rises linearly from $0$ at $r=0$ to $g_0 \approx 9.80$ at $r = R_E$, then falls as $1/r^2$, passing through $g_0/9 \approx 1.09$ at $r = 3R_E$.
\end{enumerate}
\begin{popfigure}
\begin{center}
\begin{tikzpicture}[scale=1.0]
\draw[->,thick] (0,0) -- (7.2,0) node[below left]{$r$};
\draw[->,thick] (0,0) -- (0,4.2) node[below left]{$g$};
\draw[thick,popblue] (0,0) -- (2,3);
\draw[thick,popblue,domain=2:6.8,samples=60] plot (\x,{3*4/(\x*\x)});
\draw[dashed,gray] (2,0) node[below]{$R_E$} -- (2,3) -- (0,3) node[left]{$g_0$};
\draw[dashed,gray] (6,0) node[below]{$3R_E$} -- (6,0.333) -- (0,0.333) node[left]{$g_0/9$};
\draw[dashed,gray] (1,0) node[below]{$R_E/2$} -- (1,1.5);
\fill[poporange] (1,1.5) circle (1.5pt);
\fill[poporange] (2,3) circle (1.5pt);
\fill[poporange] (6,0.333) circle (1.5pt);
\node[popblue] at (4.6,2.2) {$g\propto \dfrac{1}{r^2}$};
\node[popblue] at (0.85,2.4) {$g\propto r$};
\end{tikzpicture}
\end{center}
\legende{Variation of $g$ with distance $r$ from the centre of a uniform Earth.}
\end{popfigure}
\end{corrige}

\begin{corrige}[Exercise 6 — A strange planet]
\begin{enumerate}
\item Surface field strength:
\[
g_P = \frac{G M_P}{R_P^2} = \frac{G(2M_E)}{(R_E/2)^2} = \frac{2GM_E}{R_E^2/4} = 8\,\frac{GM_E}{R_E^2} = 8g_0. \qquad \blacksquare
\]
\item Escape velocity:
\[
v_P = \sqrt{\frac{2GM_P}{R_P}} = \sqrt{\frac{2G(2M_E)}{R_E/2}} = \sqrt{4\,\frac{2GM_E}{R_E}} = 2\sqrt{\frac{2GM_E}{R_E}} = 2v_E. \qquad \blacksquare
\]
\item Numerical values:
\[
g_P = 8 \times 9.81 = \SI{78.5}{m.s^{-2}}, \qquad
v_P = 2 \times 11.2 = \SI{22.4}{km.s^{-1}}.
\]
\end{enumerate}
\textbf{Remark:} halving the radius \emph{increases} both $g$ and $v_e$ because they grow as the radius shrinks ($g \propto 1/R^2$, $v_e \propto 1/\sqrt{R}$).
\end{corrige}

\begin{corrige}[Exercise 7 — A low-orbit satellite]
\begin{enumerate}
\item Kepler's third law: $T^2 = \dfrac{4\pi^2 r^3}{GM_E}$, so $r = \left(\dfrac{GM_E T^2}{4\pi^2}\right)^{1/3}$. With $T = 90\ \mathrm{min} = \SI{5400}{s}$:
\[
r^3 = \frac{(6.67\times10^{-11})(5.98\times10^{24})(5400)^2}{4\pi^2}
= \frac{(3.989\times10^{14})(2.916\times10^{7})}{39.48}
\approx 2.946\times10^{20}\ \mathrm{m^3},
\]
\[
r \approx (2.946\times10^{20})^{1/3} \approx \SI{6.66e6}{m} \quad (\approx 6\,660\ \mathrm{km}).
\]
\item Orbital speed:
\[
v = \frac{2\pi r}{T} = \frac{2\pi(6.66\times10^{6})}{5400} \approx \SI{7.75e3}{m.s^{-1}} \approx \SI{7.7}{km.s^{-1}}.
\]
(Check: $v = \sqrt{GM_E/r} = \sqrt{3.989\times10^{14}/6.66\times10^{6}} \approx \SI{7.74e3}{m.s^{-1}}$.)
\item \textbf{No.} A geostationary satellite must have $T = 24$ h and orbit in the equatorial plane at $r \approx \SI{4.22e7}{m}$. This satellite has $T = 90$ min $\neq 24$ h (and its altitude is only $\approx 280$ km), so it cannot be geostationary.
\end{enumerate}
\end{corrige}

\begin{corrige}[Exercise 8 — Projectile from a cliff in Limbe]
Choose the origin $O$ at the top of the cliff, $x$-axis horizontal in the direction of the throw, $y$-axis vertically \emph{downwards} (so that $g$ is positive). Initial conditions: $v_{0x} = \SI{15}{m.s^{-1}}$, $v_{0y} = 0$.
\begin{enumerate}
\item Equations of motion:
\[
x(t) = 15\,t, \qquad y(t) = \tfrac{1}{2}g\,t^2 = 4.905\,t^2 \quad (\text{SI units}).
\]
\item The stone reaches the sea when $y = 45$ m:
\[
t = \sqrt{\frac{2h}{g}} = \sqrt{\frac{2\times 45}{9.81}} = \sqrt{9.174} \approx \SI{3.03}{s}.
\]
\item Horizontal distance:
\[
x = v_{0x}\,t = 15 \times 3.03 \approx \SI{45.4}{m}.
\]
\item Velocity components at impact:
\[
v_x = 15\ \mathrm{m\,s^{-1}}, \qquad v_y = gt = 9.81 \times 3.03 \approx 29.7\ \mathrm{m\,s^{-1}}.
\]
Magnitude:
\[
v = \sqrt{v_x^2 + v_y^2} = \sqrt{15^2 + 29.7^2} = \sqrt{225 + 882} \approx \SI{33.3}{m.s^{-1}}.
\]
Direction below the horizontal:
\[
\tan\theta = \frac{v_y}{v_x} = \frac{29.7}{15} = 1.98 \quad\Longrightarrow\quad \theta \approx 63.2^{\circ}.
\]
\end{enumerate}
\end{corrige}

\begin{corrige}[Exercise 9 — Geostationary satellites (GCE style)]
\begin{enumerate}
\item \textbf{Kepler's laws} (3 marks — 1 per law):
\begin{itemize}
\item \textbf{1st law (orbits):} each planet moves on an ellipse with the Sun at one focus.
\item \textbf{2nd law (areas):} the line joining the planet to the Sun sweeps out equal areas in equal times.
\item \textbf{3rd law (periods):} the ratio $\dfrac{T^2}{r^3}$ is the same constant for all planets orbiting the Sun.
\end{itemize}
\item \textbf{Definition} (1 mark): a geostationary satellite is a satellite in a circular equatorial orbit whose period equals the Earth's rotation period (24 h) and which revolves in the same sense, so it appears fixed above one point of the equator. \textbf{Uses} (any 2, 1 mark each): television and radio broadcasting; telephone and internet relay; weather observation; GPS-type positioning support; monitoring of crops and forests.
\item \textbf{Proof} (4 marks). The gravitational attraction supplies the centripetal force:
\[
\frac{GM_E m}{r^2} = m r \omega^2 \quad \text{(1 mark)}, \qquad \omega = \frac{2\pi}{T} \quad \text{(1 mark)}.
\]
Substituting:
\[
\frac{GM_E}{r^2} = r\,\frac{4\pi^2}{T^2}
\quad\Longrightarrow\quad
T^2 = \frac{4\pi^2 r^3}{GM_E} \quad \text{(2 marks)}. \qquad \blacksquare
\]
\item With $T = 24\ \mathrm{h} = \SI{86400}{s}$:
\[
r^3 = \frac{GM_E T^2}{4\pi^2} = \frac{(6.67\times10^{-11})(5.98\times10^{24})(86400)^2}{4\pi^2}
= \frac{(3.989\times10^{14})(7.465\times10^{9})}{39.48}
\approx 7.54\times10^{22}\ \mathrm{m^3},
\]
\[
r \approx \SI{4.22e7}{m} \approx 4.22\times10^{4}\ \mathrm{km} \quad \text{(2 marks)}.
\]
Altitude: $h = r - R_E = 4.22\times10^{7} - 6.38\times10^{6} \approx \SI{3.59e7}{m} \approx 3.6\times10^{4}$ km (2 marks).

Speed: $v = \dfrac{2\pi r}{T} = \dfrac{2\pi(4.22\times10^{7})}{86400} \approx \SI{3.07e3}{m.s^{-1}} \approx \SI{3.1}{km.s^{-1}}$ (2 marks).
\item Each geostationary satellite "sees" a disc covering roughly one third of the globe (about $120^{\circ}$ of longitude). Three satellites spaced $120^{\circ}$ apart in longitude therefore cover the whole equatorial and temperate belt (diagram: circle for the Earth, three satellites at $120^{\circ}$, tangent rays from each satellite to the Earth's surface). The \textbf{polar regions} (latitudes above about $80^{\circ}$) are poorly covered because the satellites lie near the horizon there. (4 marks: diagram 1, idea of $120^{\circ}$ spacing 1, coverage of the globe 1, polar regions 1.)
\end{enumerate}
\end{corrige}

\begin{corrige}[Exercise 10 — Energy of an orbiting satellite (GCE style)]
\begin{enumerate}
\item \textbf{Gravitational potential} at a point: work done per unit mass in bringing a small test mass from infinity to that point; $V = -\dfrac{GM}{r}$ (J\,kg$^{-1}$). \textbf{Gravitational potential energy} of a mass $m$ at that point: $E_p = mV = -\dfrac{GMm}{r}$ (J) — the work needed to bring $m$ from infinity to the point. (3 marks)
\item \textbf{Proof} (4 marks). For a circular orbit, gravity provides the centripetal force:
\[
\frac{GM_E m}{r^2} = \frac{mv^2}{r}
\quad\Longrightarrow\quad
E_k = \tfrac12 mv^2 = \frac{GM_E m}{2r} \quad \text{(2 marks)}.
\]
Total energy:
\[
E = E_k + E_p = \frac{GM_E m}{2r} - \frac{GM_E m}{r} = -\frac{GM_E m}{2r} \quad \text{(2 marks)}. \qquad \blacksquare
\]
\item The negative sign means the satellite is in a \textbf{bound state}: its total energy is below the zero level (at infinity), so energy must be \emph{supplied} to it ($\Delta E = +\frac{GM_Em}{2r}$) to free it from the Earth's attraction. (2 marks)
\item Work supplied = increase of total energy:
\[
W = E_2 - E_1 = -\frac{GM_Em}{2r_2} + \frac{GM_Em}{2r_1}
= \frac{GM_Em}{2}\left(\frac{1}{r_1} - \frac{1}{r_2}\right)
= \frac{GM_Em}{2R_E}\left(\frac{1}{2} - \frac{1}{3}\right)
= \frac{GM_Em}{12R_E}.
\]
Numerically:
\[
W = \frac{(6.67\times10^{-11})(5.98\times10^{24})(100)}{12(6.38\times10^{6})}
= \frac{3.989\times10^{16}}{7.656\times10^{7}}
\approx \SI{5.2e8}{J}.
\]
(5 marks: expression of $\Delta E$ 2, simplification 1, numerical value with unit 2.)
\item As the satellite descends, its potential energy decreases (becomes more negative) and its kinetic energy would increase correspondingly; but air resistance (a non-conservative force) converts mechanical energy into heat, so the \textbf{total mechanical energy decreases} continuously. The satellite spirals inwards, heats up, and may burn up in the atmosphere. (3 marks: one per correct statement about $E_k$, $E_p$, $E$.)
\end{enumerate}
\end{corrige}

\begin{corrige}[Exercise 11 — Projectile motion and data analysis (GCE style)]
\textbf{Part A.}
\begin{enumerate}
\item Components (2 marks):
\[
v_{0x} = 20\cos 40^{\circ} = \SI{15.3}{m.s^{-1}}, \qquad
v_{0y} = 20\sin 40^{\circ} = \SI{12.9}{m.s^{-1}}.
\]
\item Time of flight — the ball lands when $y = v_{0y}t - \frac12 gt^2 = 0$, giving (3 marks):
\[
T = \frac{2v_{0y}}{g} = \frac{2 \times 12.86}{9.81} \approx \SI{2.62}{s}.
\]
\item Maximum height (3 marks):
\[
H = \frac{v_{0y}^2}{2g} = \frac{(12.86)^2}{2\times 9.81} = \frac{165.4}{19.62} \approx \SI{8.43}{m}.
\]
\item Range (2 marks):
\[
R = v_{0x}\,T = 15.32 \times 2.62 \approx \SI{40.2}{m}.
\]
\item The range formula $R = \dfrac{v_0^2 \sin 2\theta}{g}$ is maximal when $\sin 2\theta = 1$, i.e.\ $2\theta = 90^{\circ}$, so $\theta = 45^{\circ}$. (2 marks)
\end{enumerate}
\textbf{Part B.}
\begin{enumerate}
\item Converting $T$ to seconds ($1$ day $= 86400$ s) and computing (values to 3 s.f.):
\begin{center}
\begin{tabularx}{\linewidth}{|l|X|X|X|X|}
\hline
\rowcolor{popblueL} Moon & Io & Europa & Ganymede & Callisto \\
\hline $r^3$ ($\times 10^{25}$ m$^3$) & 7.51 & 30.2 & 122 & 664 \\
\hline $T^2$ ($\times 10^{9}$ s$^2$) & 2.34 & 9.41 & 38.2 & 208 \\
\hline
\end{tabularx}
\end{center}
(3 marks: correct conversion 1, correct powers 2.)
\item The graph of $T^2$ against $r^3$ is a \textbf{straight line through the origin} (3 marks: axes labelled with units 1, correct points 1, straight line and comment 1). This confirms Kepler's third law: $T^2 \propto r^3$.
\item Kepler's third law gives $T^2 = \dfrac{4\pi^2}{GM_J}\,r^3$, so the gradient of the graph is $k = \dfrac{4\pi^2}{GM_J}$. From the table, using e.g.\ Callisto:
\[
k = \frac{T^2}{r^3} = \frac{2.08\times10^{11}}{6.64\times10^{26}} \approx 3.13\times10^{-16}\ \mathrm{s^2\,m^{-3}}.
\]
Hence
\[
M_J = \frac{4\pi^2}{Gk} = \frac{39.48}{(6.67\times10^{-11})(3.13\times10^{-16})}
\approx \SI{1.9e27}{kg}.
\]
(4 marks: linking gradient to $4\pi^2/GM_J$ 1, reading the gradient 1, correct algebra 1, value $\approx 1.9\times10^{27}$ kg 1.) This is in excellent agreement with the accepted mass of Jupiter.
\end{enumerate}
\end{corrige}

\newpage

\coursec{Summary sheet}

\begin{popfigure}
\begin{tikzpicture}[mindmap, grow cyclic,
every node/.style={concept, font=\small},
root concept/.append style={concept color=popblue, font=\bfseries\small},
level 1/.append style={level distance=3.4cm, sibling angle=60},
level 2/.append style={level distance=2.2cm, sibling angle=45, font=\scriptsize},
concept connection/.append style={color=popmuted}]
\node[root concept]{Gravitational Fields}
  child[concept color=poporange]{ node {Newton's Law \& Kepler}
    child { node {$F=\dfrac{Gm_1m_2}{r^2}$} }
    child { node {$T^2\propto r^3$} }
  }
  child[concept color=popgreen]{ node {Field Strength $g$}
    child { node {$g=\dfrac{GM}{r^2}$} }
    child { node {Inside: $g\propto r$} }
  }
  child[concept color=poppurple]{ node {Potential \& Energy}
    child { node {$V=-\dfrac{GM}{r}$} }
    child { node {$E_p=-\dfrac{GMm}{r}$} }
  }
  child[concept color=poppink]{ node {Escape \& Orbits}
    child { node {$v_{esc}=\sqrt{\dfrac{2GM}{R}}$} }
    child { node {$v_{orb}=\sqrt{\dfrac{GM}{r}}$} }
  }
  child[concept color=popgold]{ node {Satellites}
    child { node {Geostationary: $T=24$ h} }
    child { node {$E=-\dfrac{GMm}{2r}$} }
  }
  child[concept color=popblueL]{ node {Projectile Motion}
    child { node {Parabola, $g$ constant} }
    child { node {Range $R=\dfrac{u^2\sin2\theta}{g}$} }
  };
\end{tikzpicture}
\legende{Mind map of the chapter: Gravitational Fields.}
\end{popfigure}

\begin{bilan}
\textbf{Key definitions}
\begin{itemize}
\item \cle{Gravitational field}: region of space around a mass in which another mass experiences an attractive force.
\item \cle{Gravitational field strength} $g$: force per unit mass placed at the point, $\vec{g}=\vec{F}/m$; unit $\mathrm{N\,kg^{-1}}$ (equivalent to $\mathrm{m\,s^{-2}}$). It is a \emph{vector}, directed towards the attracting mass.
\item \cle{Gravitational potential} $V$ at a point: work done per unit mass by an external agent in bringing a small test mass from infinity to that point; unit $\mathrm{J\,kg^{-1}}$. It is a \emph{scalar}, always negative (zero at infinity).
\item \cle{Gravitational potential energy} $E_p$: work done to bring a mass $m$ from infinity to the point; $E_p=mV$.
\item \cle{Escape velocity}: minimum launch speed for a body to leave a planet's attraction and reach infinity with zero residual speed.
\item \cle{Geostationary satellite}: satellite in a circular equatorial orbit whose period equals the Earth's rotation period ($T=24$ h), so it appears fixed above one point of the equator.
\end{itemize}

\textbf{Laws and formulas to know}
\begin{itemize}
\item \cle{Newton's law of gravitation} (inverse-square law):
\[ F=\frac{Gm_1m_2}{r^2},\qquad G=\SI{6.67e-11}{N.m^2.kg^{-2}}. \]
Attractive, acts along the line joining the centres; valid for point masses and spherical bodies ($r$ measured centre to centre).
\item \cle{Kepler's laws}: (1) planets move on ellipses with the Sun at one focus; (2) the radius vector sweeps equal areas in equal times; (3) $T^2\propto r^3$, i.e.\ $\dfrac{T^2}{r^3}=\dfrac{4\pi^2}{GM}$ for circular orbits.
\item Field strength of a spherical mass $M$: $g=\dfrac{GM}{r^2}$ outside ($r\ge R$); at the surface $g_0=\dfrac{GM}{R^2}$.
\item \cle{Variation of $g$}: at altitude $h$, $g_h=\dfrac{GM}{(R+h)^2}=g_0\dfrac{R^2}{(R+h)^2}$; at depth $d$ (uniform Earth), $g_d=g_0\left(1-\dfrac{d}{R}\right)$, since only the inner sphere of radius $(R-d)$ attracts. Inside a uniform sphere, $g\propto r$; at the centre $g=0$.
\item Potential and energy:
\[ V=-\frac{GM}{r},\qquad E_p=-\frac{GMm}{r},\qquad g=-\frac{dV}{dr}. \]
\item Escape and orbital speeds:
\[ v_{esc}=\sqrt{\frac{2GM}{R}}=\sqrt{2g_0R},\qquad v_{orb}=\sqrt{\frac{GM}{r}}=\sqrt{\frac{g_0R^2}{r}},\qquad T=2\pi\sqrt{\frac{r^3}{GM}}. \]
\item Energy of an orbiting satellite: $E_k=\dfrac{GMm}{2r}$, $E_p=-\dfrac{GMm}{r}$, total $E=-\dfrac{GMm}{2r}$ (bound, negative).
\item Projectile motion (uniform $g$): horizontal $x=ut\cos\theta$; vertical $y=ut\sin\theta-\tfrac{1}{2}gt^2$; trajectory $y=x\tan\theta-\dfrac{gx^2}{2u^2\cos^2\theta}$ (parabola); time of flight $t_f=\dfrac{2u\sin\theta}{g}$; maximum height $H=\dfrac{u^2\sin^2\theta}{2g}$; range $R=\dfrac{u^2\sin2\theta}{g}$, maximum at $\theta=45^\circ$.
\end{itemize}

\textbf{Methods}
\begin{itemize}
\item Satellite problems: equate gravitational force to centripetal force, $\dfrac{GMm}{r^2}=\dfrac{mv^2}{r}=mr\omega^2$, then extract $v$, $\omega$ or $T$.
\item Energy problems: apply conservation of mechanical energy $E_k+E_p=\text{constant}$ between two points (e.g.\ surface and infinity for escape velocity).
\item Projectiles: resolve the initial velocity into components; treat horizontal (uniform) and vertical (uniformly accelerated) motions independently; link them through the common time $t$.
\item Always measure $r$ from the \emph{centre} of the Earth: $r=R+h$.
\end{itemize}

\textbf{Common mistakes}
\begin{itemize}
\item Using the altitude $h$ instead of $r=R+h$ in $g=GM/r^2$ or $v_{orb}$.
\item Forgetting the negative sign of $V$ and $E_p$; potential \emph{increases} (becomes less negative) with distance.
\item Confusing $v_{esc}=\sqrt{2gR}$ with $v_{orb}=\sqrt{gR}$ at the surface: they differ by a factor $\sqrt{2}$.
\item Writing $T^2\propto r^2$ instead of $T^2\propto r^3$ (Kepler's third law).
\item Mixing units: convert km to m, hours to seconds, before substituting.
\item Believing a satellite is ``weightless'' because there is no gravity: gravity supplies the centripetal force; the astronaut is in free fall.
\end{itemize}

\textbf{GCE tips}
\begin{itemize}
\item Quote Newton's law in words \emph{and} symbols; state that $G$ is the universal gravitational constant (do not confuse $G$ with $g$).
\item Show the derivation $mg=\dfrac{GMm}{R^2}\Rightarrow g=\dfrac{GM}{R^2}$: it links field strength to the inverse-square law and earns method marks.
\item For geostationary satellites, state the three conditions: circular orbit, equatorial plane, period $T=24$ h in the direction of the Earth's rotation.
\item Carry units through every numerical answer and give sensible significant figures (usually 3 s.f.).
\item Sketch the variation of $g$ with $r$: linear rise from the centre to $R$, then inverse-square decrease for $r>R$.
\end{itemize}
\end{bilan}

\findecours

\end{document}
